\documentclass[10pt]{article}
\usepackage[utf8]{inputenc}
\usepackage[T1]{fontenc}
\usepackage{graphicx}
\usepackage[export]{adjustbox}
\graphicspath{ {./images/} }
\usepackage{hyperref}
\hypersetup{colorlinks=true, linkcolor=blue, filecolor=magenta, urlcolor=cyan,}
\urlstyle{same}
\usepackage{amsmath}
\usepackage{amsfonts}
\usepackage{amssymb}
\usepackage[version=4]{mhchem}
\usepackage{stmaryrd}
\usepackage{bbold}
\usepackage{mathrsfs}
\usepackage{caption}
\usepackage{multirow}

%New command to display footnote whose markers will always be hidden
\let\svthefootnote\thefootnote
\newcommand\blfootnotetext[1]{%
  \let\thefootnote\relax\footnote{#1}%
  \addtocounter{footnote}{-1}%
  \let\thefootnote\svthefootnote%
}
%Overriding the \footnotetext command to hide the marker if its value is `0`
\let\svfootnotetext\footnotetext
\renewcommand\footnotetext[2][?]{%
  \if\relax#1\relax%
    \ifnum\value{footnote}=0\blfootnotetext{#2}\else\svfootnotetext{#2}\fi%
  \else%
    \if?#1\ifnum\value{footnote}=0\blfootnotetext{#2}\else\svfootnotetext{#2}\fi%
    \else\svfootnotetext[#1]{#2}\fi%
  \fi
}
\DeclareUnicodeCharacter{22C5}{\ifmmode\cdot\else{$\cdot$}\fi}
\DeclareUnicodeCharacter{0394}{\ifmmode\Delta\else{$\Delta$}\fi}
\DeclareUnicodeCharacter{22EE}{\ifmmode\vdots\else{$\vdots$}\fi}
\title{Chapter 15: Binary Response Models}

\author{}
\date{}

\begin{document}
\maketitle

\subsection*{15.1 Introduction}
In binary response models, the variable to be explained, $y$, is a random variable taking on the values zero and one, which indicate whether or not a certain event has occurred. For example, $y=1$ if a person is employed, $y=0$ otherwise; $y=1$ if a family contributes to charity during a particular year, $y=0$ otherwise; $y=1$ if a firm has a particular type of pension plan, $y=0$ otherwise. Regardless of the definition of $y$, it is traditional to refer to $y=1$ as a success and $y=0$ as a failure.

As in the case of linear models, we often call $y$ the explained variable, the response variable, the dependent variable, or the endogenous variable; $\mathbf{x} \equiv\left(x_{1}, x_{2}, \ldots, x_{K}\right)$ is the vector of explanatory variables, regressors, independent variables, exogenous variables, or covariates.

In binary response models, interest lies primarily in the response probability,


\begin{equation*}
p(\mathbf{x}) \equiv \mathrm{P}(y=1 \mid \mathbf{x})=\mathrm{P}\left(y=1 \mid x_{1}, x_{2}, \ldots, x_{K}\right), \tag{15.1}
\end{equation*}


for various values of $\mathbf{x}$. For example, when $y$ is an employment indicator, $\mathbf{x}$ might contain various individual characteristics such as education, age, marital status, and other factors that affect employment status, such as a binary indicator variable for participation in a recent job training program, or measures of past criminal behavior.

For a continuous variable, $x_{j}$, the partial effect of $x_{j}$ on the response probability is


\begin{equation*}
\frac{\partial \mathrm{P}(y=1 \mid \mathbf{x})}{\partial x_{j}}=\frac{\partial p(\mathbf{x})}{\partial x_{j}} . \tag{15.2}
\end{equation*}


When multiplied by $\Delta x_{j}$, equation (15.2) gives the approximate change in $\mathrm{P}(y=1 \mid \mathbf{x})$ when $x_{j}$ increases by $\Delta x_{j}$, holding all other variables fixed (for "small" $\Delta x_{j}$ ). Of course if, say, $x_{1} \equiv z$ and $x_{2} \equiv z^{2}$ for some variable $z$ (for example, $z$ could be work experience), we would be interested in $\partial p(\mathbf{x}) / \partial z$.

If $x_{K}$ is a binary variable, interest lies in


\begin{equation*}
p\left(x_{1}, x_{2}, \ldots, x_{K-1}, 1\right)-p\left(x_{1}, x_{2}, \ldots, x_{K-1}, 0\right), \tag{15.3}
\end{equation*}


which is the difference in response probabilities when $x_{K}=1$ and $x_{K}=0$. For most of the models we consider, whether a variable $x_{j}$ is continuous or discrete, the partial effect of $x_{j}$ on $p(\mathbf{x})$ depends on all of $\mathbf{x}$.

In studying binary response models, we need to recall some basic facts about Bernoulli (zero-one) random variables. The only difference between the setup here and that in basic statistics is the conditioning on $\mathbf{x}$. If $\mathrm{P}(y=1 \mid \mathbf{x})=p(\mathbf{x})$ then $\mathrm{P}(y=0 \mid \mathbf{x})=1-p(\mathbf{x}), \mathrm{E}(y \mid \mathbf{x})=p(\mathbf{x})$, and $\operatorname{Var}(y \mid \mathbf{x})=p(\mathbf{x})[1-p(\mathbf{x})]$.

\subsection*{15.2 Linear Probability Model for Binary Response}
The linear probability model (LPM) for binary response $y$ is specified as\\
$\mathrm{P}(y=1 \mid \mathbf{x})=\beta_{0}+\beta_{1} x_{1}+\beta_{2} x_{2}+\cdots+\beta_{K} x_{K}$.

As usual, the $x_{j}$ can be functions of underlying explanatory variables, which would simply change the interpretations of the $\beta_{j}$. Assuming that $x_{1}$ is not functionally related to the other explanatory variables, $\beta_{1}=\partial \mathrm{P}(y=1 \mid \mathbf{x}) / \partial x_{1}$. Therefore, $\beta_{1}$ is the change in the probability of success given a one-unit increase in $x_{1}$. If $x_{1}$ is a binary explanatory variable, $\beta_{1}$ is just the difference in the probability of success when $x_{1}=1$ and $x_{1}=0$, holding the other $x_{j}$ fixed.

Using functions such as quadratics, logarithms, and so on among the independent variables causes no new difficulties. The important point is that the $\beta_{j}$ now measures the effects of the explanatory variables $x_{j}$ on a particular probability.

In deciding on an appropriate estimation technique, it is useful to derive the conditional mean and variance of $y$. Since $y$ is a Bernoulli random variable, these are simply


\begin{equation*}
\mathrm{E}(y \mid \mathbf{x})=\beta_{0}+\beta_{1} x_{1}+\beta_{2} x_{2}+\cdots+\beta_{K} x_{K} \tag{15.5}
\end{equation*}


$\operatorname{Var}(y \mid \mathbf{x})=\mathbf{x} \boldsymbol{\beta}(1-\mathbf{x} \boldsymbol{\beta})$,\\
where $\mathbf{x} \boldsymbol{\beta}$ is shorthand for the right-hand side of equation (15.5).\\
Equation (15.5) implies that, given a random sample, the ordinary least squares (OLS) regression of $y$ on $1, x_{1}, x_{2}, \ldots, x_{K}$ produces consistent and even unbiased estimators of the $\beta_{j}$. Equation (15.6) means that heteroskedasticity is present unless all of the slope coefficients $\beta_{1}, \ldots, \beta_{K}$ are zero. A nice way to deal with this issue is to use standard heteroskedasticity-robust standard errors and $t$ statistics. Further, robust tests of multiple restrictions should also be used. There is one case where the usual $F$ statistic can be used, and that is to test for joint significance of all variables (leaving the constant unrestricted). This test is asymptotically valid because $\operatorname{Var}(y \mid \mathbf{x})$ is constant under this particular null hypothesis.

If we operate under the assumption that $\mathrm{P}(y=1 \mid \mathbf{x})$ is given by equation (15.4), then we can obtain an asymptotically more efficient estimator by applying weighted least squares (WLS). Let $\hat{\boldsymbol{\beta}}$ be the OLS estimator, and let $\hat{y}_{i}$ denote the OLS fitted values. Then, provided $0<\hat{y}_{i}<1$ for all observations $i$, define the estimated standard deviation as $\hat{\sigma}_{i} \equiv\left[\hat{y}_{i}\left(1-\hat{y}_{i}\right)\right]^{1 / 2}$. Then the WLS estimator, $\boldsymbol{\beta}^{*}$, is obtained from the OLS regression\\
$y_{i} / \hat{\sigma}_{i}$ on $1 / \hat{\sigma}_{i}, x_{i 1} / \hat{\sigma}_{i}, \ldots, x_{i K} / \hat{\sigma}_{i}, \quad i=1,2, \ldots, N$.

The usual standard errors from this regression are valid, as follows from the treatment of weighted least squares in Chapter 12. In addition, all other testing can be done using $F$ statistics or $L M$ statistics using weighted regressions.

If some of the OLS fitted values are not between zero and one, WLS analysis is not possible without ad hoc adjustments to bring deviant fitted values into the unit interval. Further, since the OLS fitted value $\hat{y}_{i}$ is an estimate of the conditional probability $\mathrm{P}\left(y_{i}=1 \mid \mathbf{x}_{i}\right)$, it is somewhat awkward if the predicted probability is negative or above unity.

Aside from the issue of fitted values being outside the unit interval, the LPM implies that a ceteris paribus unit increase in $x_{j}$ always changes $\mathrm{P}(y=1 \mid \mathbf{x})$ by the same amount, regardless of the initial value of $x_{j}$. This feature of the LPM cannot literally be true because continually increasing one of the $x_{j}$ would eventually drive $\mathrm{P}(y=1 \mid \mathbf{x})$ to be less than zero or greater than one.

A sensible posture is to simply view the LPM as the linear projection of $y$ on the explanatory variables. Recall from Chapter 2 that the linear projection provides the best least squares fit among linear functions of $\mathbf{x}$ (although $\mathbf{x}$ itself might include nonlinear functions of underlying explanatory variables). If we operate within this scenario, OLS estimation of the LPM is attractive because it consistently estimates the parameters in the linear projection. The WLS estimator can be viewed as a linear projection on weighted variables, which may not be of as much interest, particularly because the weights are not obtained from $\operatorname{Var}(y \mid \mathbf{x})$ if (15.4) fails.

If the main purpose of estimating a binary response model is to approximate the partial effects of the explanatory variables, averaged across the distribution of $\mathbf{x}$, then the LPM often does a very good job. (Some evidence on how well it does can be obtained by comparing the OLS coefficients with the average partial effects from the nonlinear models we turn to in Section 15.3.) The fact that some predicted probabilities are outside the unit interval need not be a serious concern. But there is no guarantee that the LPM provides good estimates of the partial effects for a wide range of covariate values, especially for extreme values of $\mathbf{x}$.

Example 15.1 (Married Women's Labor Force Participation): We use the data from MROZ.RAW to estimate a linear probability model for labor force participation (inlf) of married women. Of the 753 women in the sample, 428 report working nonzero hours during the year. The variables we use to explain labor force participation are age, education, experience, nonwife income in thousands (nwifeinc), number of children less than six years of age (kidslt6), and number of children between 6 and 18\\
inclusive (kidsge6); 606 women report having no young children, while 118 report having exactly one young child. The usual OLS standard errors are in parentheses, while the heteroskedasticity-robust standard errors are in brackets:

$$
\begin{aligned}
& \widehat{\text { inlf }}=.586-.0034 \text { nwifeinc }+.038 \text { educ }+.039 \text { exper }-.00060 \text { exper }^{2} \\
& \text { (.154) (.0014) (.007) (.006) (.00018) } \\
& \text { [.151] [.0015] [.007] [.006] [.00019] } \\
& -.016 \text { age - . } 262 \text { kidslt6 + . } 013 \text { kidsge6 } \\
& \text { (.002) (.034) (.013) } \\
& \text { [.002] [.032] [.013] } \\
& N=753, \quad R^{2}=.264
\end{aligned}
$$

With the exception of kidsge6, all coefficients have sensible signs and are statistically significant; kidsge6 is neither statistically significant nor practically important. The coefficient on nwifeinc means that if nonwife income increases by $10(\$ 10,000)$, the probability of being in the labor force is predicted to fall by .034 . This is a small effect given that an increase in income by $\$ 10,000$ in 1975 dollars is very large in this sample. (The average of nwifeinc is about $\$ 20,129$ with standard deviation $\$ 11,635$.) Having one more small child is estimated to reduce the probability of inlf $=1$ by about .262, which is a fairly large effect.

Of the 753 fitted probabilities, 33 are outside the unit interval. Rather than using some adjustment to those 33 fitted values and applying WLS, we just use OLS and report heteroskedasticity-robust standard errors. Interestingly, these differ in practically unimportant ways from the usual OLS standard errors.

The case for the LPM is even stronger if most of the $x_{j}$ are discrete and take on only a few values. In the previous example, to allow a diminishing effect of young children on the probability of labor force participation, we can break kidslto into three binary indicators: no young children, one young child, and two or more young children. The last two indicators can be used in place of kidslto to allow the first young child to have a larger effect than subsequent young children. (Interestingly, when this method is used, the marginal effects of the first and second young children are virtually the same. The estimated effect of the first child is about -.263 , and the additional reduction in the probability of labor force participation for the next child is about -.274.)

In the extreme case where the model is saturated-that is, $\mathbf{x}$ contains dummy variables for mutually exclusive and exhaustive categories-the LPM is completely general. The fitted probabilities are simply the average $y_{i}$ within each cell defined by the\\
different values of $\mathbf{x}$; we need not worry about fitted probabilities less than zero or greater than one. See Problem 15.1.

\subsection*{15.3 Index Models for Binary Response: Probit and Logit}
We now study binary response models of the form\\
$\mathrm{P}(y=1 \mid \mathbf{x})=G(\mathbf{x} \boldsymbol{\beta}) \equiv p(\mathbf{x})$,\\
where $\mathbf{x}$ is $1 \times K, \boldsymbol{\beta}$ is $K \times 1$, and we take the first element of $\mathbf{x}$ to be unity. Examples where $\mathbf{x}$ does not contain unity are rare in practice. For the linear probability model, $G(z)=z$ is the identity function, which means that the response probabilities cannot be between 0 and 1 for all $\mathbf{x}$ and $\boldsymbol{\beta}$. In this section we assume that $G(\cdot)$ takes on values in the open unit interval: $0<G(z)<1$ for all $z \in \mathbb{R}$.

The model in equation (15.8) is generally called an index model because it restricts the way in which the response probability depends on $\mathbf{x}: p(\mathbf{x})$ is a function of $\mathbf{x}$ only through the index $\mathbf{x} \boldsymbol{\beta}=\beta_{1}+\beta_{2} x_{2}+\cdots+\beta_{K} x_{K}$. The function $G$ maps the index into the response probability.

In most applications, $G$ is a cumulative distribution function (cdf) whose specific form can sometimes be derived from an underlying economic model. For example, in Problem 15.2 you are asked to derive an index model from a utility-based model of charitable giving. The binary indicator $y$ equals unity if a family contributes to charity and zero otherwise. The vector $\mathbf{x}$ contains family characteristics, income, and the price of a charitable contribution (as determined by marginal tax rates). Under a normality assumption on a particular unobservable taste variable, $G$ is the standard normal cdf.

Index models where $G$ is a cdf can be derived more generally from an underlying latent variable model, as in Example 13.1:\\
$y^{*}=\mathbf{x} \boldsymbol{\beta}+e, \quad y=1\left[y^{*}>0\right]$,\\
where $e$ is a continuously distributed variable independent of $\mathbf{x}$ and the distribution of $e$ is symmetric about zero; recall from Chapter 13 that $1[\cdot]$ is the indicator function. If $G$ is the cdf of $e$, then, because the pdf of $e$ is symmetric about zero, $1-G(-z)= G(z)$ for all real numbers $z$. Therefore,\\
$\mathrm{P}(y=1 \mid \mathbf{x})=\mathrm{P}\left(y^{*}>0 \mid \mathbf{x}\right)=\mathrm{P}(e>-\mathbf{x} \boldsymbol{\beta} \mid \mathbf{x})=1-G(-\mathbf{x} \boldsymbol{\beta})=G(\mathbf{x} \boldsymbol{\beta})$,\\
which is exactly equation (15.8).\\
There is no particular reason for requiring $e$ to be symmetrically distributed in the latent variable model, but this happens to be the case for the binary response models applied most often.

In most applications of binary response models, the primary goal is to explain the effects of the $x_{j}$ on the response probability $\mathrm{P}(y=1 \mid \mathbf{x})$. The latent variable formulation tends to give the impression that we are primarily interested in the effects of each $x_{j}$ on $y^{*}$. As we will see, the direction of the effects of $x_{j}$ on $\mathrm{E}\left(y^{*} \mid \mathbf{x}\right)=\mathbf{x} \boldsymbol{\beta}$ and on $\mathrm{E}(y \mid \mathbf{x})=\mathrm{P}(y=1 \mid \mathbf{x})=G(\mathbf{x} \boldsymbol{\beta})$ are the same. But the latent variable $y^{*}$ rarely has a well-defined unit of measurement (for example, $y^{*}$ might be measured in utility units). Therefore, the magnitude of $\beta_{j}$ is not especially meaningful except in special cases.

The probit model is the special case of equation (15.8) with


\begin{equation*}
G(z)=\Phi(z) \equiv \int_{-\infty}^{z} \phi(v) \mathrm{d} v \tag{15.10}
\end{equation*}


where $\phi(z)$ is the standard normal density


\begin{equation*}
\phi(z)=(2 \pi)^{-1 / 2} \exp \left(-z^{2} / 2\right) . \tag{15.11}
\end{equation*}


The probit model can be derived from the latent variable formulation when $e$ has a standard normal distribution.

The logit model is a special case of equation (15.8) with


\begin{equation*}
G(z)=\Lambda(z) \equiv \exp (z) /[1+\exp (z)] \tag{15.12}
\end{equation*}


This model arises from the model (15.9) when $e$ has a standard logistic distribution.\\
The general specification (15.8) allows us to cover probit, logit, and a number of other binary choice models in one framework. In fact, in what follows we do not even need $G$ to be a cdf, but we do assume that $G(z)$ is strictly between zero and unity for all real numbers $z$.

In order to successfully apply probit and logit models, it is important to know how to interpret the $\beta_{j}$ on both continuous and discrete explanatory variables. First, if $x_{j}$ is continuous,


\begin{equation*}
\frac{\partial p(\mathbf{x})}{\partial x_{j}}=g(\mathbf{x} \boldsymbol{\beta}) \beta_{j}, \quad \text { where } g(z) \equiv \frac{\mathrm{d} G}{\mathrm{~d} z}(z) \tag{15.13}
\end{equation*}


Therefore, the partial effect of $x_{j}$ on $p(\mathbf{x})$ depends on $\mathbf{x}$ through $g(\mathbf{x} \boldsymbol{\beta})$. If $G(\cdot)$ is a strictly increasing cdf, as in the probit and logit cases, $g(z)>0$ for all $z$. Therefore, the sign of the effect is given by the sign of $\beta_{j}$. Also, the relative effects do not depend on $\mathbf{x}$ : for continuous variables $x_{j}$ and $x_{h}$, the ratio of the partial effects is constant and given by the ratio of the corresponding coefficients: $\frac{\partial p(\mathbf{x}) / \partial x_{j}}{\partial p(\mathbf{x}) / \partial x_{h}}=\beta_{j} / \beta_{h}$. In the typical\\
case that $g$ is a symmetric density about zero, with unique mode at zero, the largest effect is when $\mathbf{x} \boldsymbol{\beta}=0$. For example, in the probit case with $g(z)=\phi(z), g(0)=\phi(0) =1 / \sqrt{2 \pi} \approx .399$. In the logit case, $g(z)=\exp (z) /[1+\exp (z)]^{2}$, and so $g(0)=.25$.

If $x_{K}$ is a binary explanatory variable, then the partial effect from changing $x_{K}$ from zero to one, holding all other variables fixed, is simply\\
$G\left(\beta_{1}+\beta_{2} x_{2}+\cdots+\beta_{K-1} x_{K-1}+\beta_{K}\right)-G\left(\beta_{1}+\beta_{2} x_{2}+\cdots+\beta_{K-1} x_{K-1}\right)$.

Again, this expression depends on all other values of the other $x_{j}$. For example, if $y$ is an employment indicator and $x_{j}$ is a dummy variable indicating participation in a job training program, then expression (15.14) is the change in the probability of employment due to the job training program; this depends on other characteristics that affect employability, such as education and experience. Knowing the sign of $\beta_{K}$ is enough to determine whether the program had a positive or negative effect. But to find the magnitude of the effect, we have to estimate expression (15.14).

We can also use the difference in expression (15.14) for other kinds of discrete variables (such as number of children). If $x_{K}$ denotes this variable, then the effect on the probability of $x_{K}$ going from $c_{K}$ to $c_{K}+1$ is simply


\begin{align*}
G\left[\beta_{1}\right. & \left.+\beta_{2} x_{2}+\cdots+\beta_{K-1} x_{K-1}+\beta_{K}\left(c_{K}+1\right)\right] \\
& -G\left(\beta_{1}+\beta_{2} x_{2}+\cdots+\beta_{K-1} x_{K-1}+\beta_{K} c_{K}\right) \tag{15.15}
\end{align*}


It is straightforward to include standard functional forms among the explanatory variables. For example, in the model\\
$\mathrm{P}(y=1 \mid \mathbf{z})=G\left[\beta_{0}+\beta_{1} z_{1}+\beta_{2} z_{1}^{2}+\beta_{3} \log \left(z_{2}\right)+\beta_{4} z_{3}\right]$,\\
the partial effect of $z_{1}$ on $\mathrm{P}(y=1 \mid \mathbf{z})$ is $\partial \mathrm{P}(y=1 \mid \mathbf{z}) / \partial z_{1}=g(\mathbf{x} \boldsymbol{\beta})\left(\beta_{1}+2 \beta_{2} z_{1}\right)$, where $\mathbf{x} \boldsymbol{\beta}=\beta_{0}+\beta_{1} z_{1}+\beta_{2} z_{1}^{2}+\beta_{3} \log \left(z_{2}\right)+\beta_{4} z_{3}$. It follows that if the quadratic in $z_{1}$ has a hump shape or a $\mathbf{U}$ shape, the turning point in the response probability is $\left|\beta_{1} /\left(2 \beta_{2}\right)\right|$ (because $g(\mathbf{x} \boldsymbol{\beta})>0$ ). Also, $\partial \mathrm{P}(y=1 \mid \mathbf{z}) / \partial \log \left(z_{2}\right)=g(\mathbf{x} \boldsymbol{\beta}) \beta_{3}$, and so $g(\mathbf{x} \boldsymbol{\beta})\left(\beta_{3} / 100\right)$ is the approximate change in $\mathrm{P}(y=1 \mid \mathbf{z})$ given a 1 percent increase in $z_{2}$. Models with interactions among explanatory variables, including interactions between discrete and continuous variables, are handled similarly. When measuring the effects of discrete variables, we should use expression (15.15).

\subsection*{15.4 Maximum Likelihood Estimation of Binary Response Index Models}
Assume we have $N$ independent, identically distributed observations following the model (15.8). Since we essentially covered the case of probit in Chapter 13, the\\
discussion here will be brief. To estimate the model by (conditional) maximum likelihood (MLE), we need the log-likelihood function for each $i$. The density of $y_{i}$ given $\mathbf{x}_{i}$ can be written as\\
$f\left(y \mid \mathbf{x}_{i} ; \boldsymbol{\beta}\right)=\left[G\left(\mathbf{x}_{i} \boldsymbol{\beta}\right)\right]^{y}\left[1-G\left(\mathbf{x}_{i} \boldsymbol{\beta}\right)\right]^{1-y}, \quad y=0,1$.

The log likelihood for observation $i$ is a function of the $K \times 1$ vector of parameters and the data $\left(\mathbf{x}_{i}, y_{i}\right)$ :\\
$\ell_{i}(\boldsymbol{\beta})=y_{i} \log \left[G\left(\mathbf{x}_{i} \boldsymbol{\beta}\right)\right]+\left(1-y_{i}\right) \log \left[1-G\left(\mathbf{x}_{i} \boldsymbol{\beta}\right)\right]$.\\
(Recall from Chapter 13 that, technically speaking, we should distinguish the "true" value of beta, $\boldsymbol{\beta}_{\mathrm{o}}$, from a generic value. For conciseness we do not do so here.) Restricting $G(\cdot)$ to be strictly between zero and one ensures that $\ell_{i}(\boldsymbol{\beta})$ is well defined for all values of $\boldsymbol{\beta}$.

As usual, the $\log$ likelihood for a sample size of $N$ is $\mathscr{L}(\boldsymbol{\beta})=\sum_{i=1}^{N} \ell_{i}(\boldsymbol{\beta})$, and the MLE of $\boldsymbol{\beta}$, denoted $\hat{\boldsymbol{\beta}}$, maximizes this log likelihood. If $G(\cdot)$ is the standard normal cdf, then $\hat{\boldsymbol{\beta}}$ is the probit estimator; if $G(\cdot)$ is the logistic cdf, then $\hat{\boldsymbol{\beta}}$ is the logit estimator. From the general maximum likelihood results we know that $\hat{\boldsymbol{\beta}}$ is consistent and asymptotically normal. We can also easily estimate the asymptotic variance $\hat{\boldsymbol{\beta}}$.

We assume that $G(\cdot)$ is twice continuously differentiable, an assumption that is usually satisfied in applications (and, in particular, for probit and logit). As before, the function $g(z)$ is the derivative of $G(z)$. For the probit model, $g(z)=\phi(z)$, and for the logit model, $g(z)=\exp (z) /[1+\exp (z)]^{2}$.

Using the same calculations for the probit example as in Chapter 13, the score of the conditional log likelihood for observation $i$ can be shown to be


\begin{equation*}
\mathbf{s}_{i}(\boldsymbol{\beta}) \equiv \frac{g\left(\mathbf{x}_{i} \boldsymbol{\beta}\right) \mathbf{x}_{i}^{\prime}\left[y_{i}-G\left(\mathbf{x}_{i} \boldsymbol{\beta}\right)\right]}{G\left(\mathbf{x}_{i} \boldsymbol{\beta}\right)\left[1-G\left(\mathbf{x}_{i} \boldsymbol{\beta}\right)\right]} . \tag{15.18}
\end{equation*}


Similarly, the expected value of the Hessian conditional on $\mathbf{x}_{i}$ is


\begin{equation*}
-\mathrm{E}\left[\mathbf{H}_{i}(\boldsymbol{\beta}) \mid \mathbf{x}_{i}\right]=\frac{\left[g\left(\mathbf{x}_{i} \boldsymbol{\beta}\right)\right]^{2} \mathbf{x}_{i}^{\prime} \mathbf{x}_{i}}{\left\{G\left(\mathbf{x}_{i} \boldsymbol{\beta}\right)\left[1-G\left(\mathbf{x}_{i} \boldsymbol{\beta}\right)\right]\right\}} \equiv \mathbf{A}\left(\mathbf{x}_{i}, \boldsymbol{\beta}\right), \tag{15.19}
\end{equation*}


which is a $K \times K$ positive semidefinite (p.s.d.) matrix for each $i$. From the general conditional MLE results in Chapter 13, $\operatorname{Avar}(\hat{\boldsymbol{\beta}})$ is estimated as\\
$\widehat{\operatorname{Avar}(\hat{\boldsymbol{\beta}})} \equiv\left\{\sum_{i=1}^{N} \frac{\left[g\left(\mathbf{x}_{i} \hat{\boldsymbol{\beta}}\right)\right]^{2} \mathbf{x}_{i}^{\prime} \mathbf{x}_{i}}{G\left(\mathbf{x}_{i} \hat{\boldsymbol{\beta}}\right)\left[1-G\left(\mathbf{x}_{i} \hat{\boldsymbol{\beta}}\right)\right]}\right\}^{-1} \equiv \hat{\mathbf{V}}$.

In most cases the inverse exists, and when it does, $\hat{\mathbf{V}}$ is positive definite. If the matrix in equation (15.20) is not invertible, then perfect collinearity probably exists among the regressors.

As usual, we treat $\hat{\boldsymbol{\beta}}$ as being normally distributed with mean zero and variance matrix in equation (15.20). The (asymptotic) standard error of $\hat{\beta}_{j}$ is the square root of the $j$ th diagonal element of $\hat{\mathbf{V}}$. These can be used to construct $t$ statistics, which have a limiting standard normal distribution, and to construct approximate confidence intervals for each population parameter. These are reported with the estimates for packages that perform logit and probit. We discuss multiple hypothesis testing in the next section.

Some packages also compute Huber-White standard errors as an option for probit and logit analysis, using the sandwich form (13.71). While the robust variance matrix is consistent, using it in equation (15.20) means we must think that the binary response model is incorrectly specified. Unlike with nonlinear regression, in a binary response model it is not possible to correctly specify $\mathrm{E}(y \mid \mathbf{x})$ but to misspecify $\operatorname{Var}(y \mid \mathbf{x})$. Once we have specified $\mathrm{P}(y=1 \mid \mathbf{x})$, we have specified all conditional moments of $y$ given $\mathbf{x}$. Nevertheless, as we discussed in Section 13.11, it may be prudent to act as if all models are merely approximations to the truth, in which case inference should be based on the sandwich estimator in equation (13.71). (The sandwich estimator that uses the expected Hessian, as in equation (12.49), makes no sense in the binary response context because the expected Hessian cannot be computed without assumption (15.8).)

In Section 15.8 we will see that, when using binary response models with panel data, it is sometimes important to compute variance matrix estimators that are robust to serial dependence. But this need arises as a result of dependence across time or subgroup, and not because the response probability is misspecified.

\subsection*{15.5 Testing in Binary Response Index Models}
Any of the three tests from general MLE analysis-the Wald, likelihood ratio (LR), or Lagrange multiplier (LM) test-can be used to test hypotheses in binary response contexts. Since the tests are all asymptotically equivalent under local alternatives, the choice of statistic usually depends on computational simplicity (since finite sample comparisons must be limited in scope). In the following subsections we discuss some testing situations that often arise in binary choice analysis, and we recommend particular tests for their computational advantages.

\subsection*{15.5.1 Testing Multiple Exclusion Restrictions}
Consider the model\\
$\mathrm{P}(y=1 \mid \mathbf{x}, \mathbf{z})=G(\mathbf{x} \boldsymbol{\beta}+\mathbf{z} \boldsymbol{\gamma})$,\\
where $\mathbf{x}$ is $1 \times K$ and $\mathbf{z}$ is $1 \times Q$. We wish to test the null hypothesis $\mathrm{H}_{0}: \boldsymbol{\gamma}=\mathbf{0}$, so we are testing $Q$ exclusion restrictions. The elements of $\mathbf{z}$ can be functions of $\mathbf{x}$, such as quadratics and interactions-in which case the test is a pure functional form test. Or, the $\mathbf{z}$ can be additional explanatory variables. For example, $\mathbf{z}$ could contain dummy variables for occupation or region. In any case, the form of the test is the same.

Some packages, such as Stata, compute the Wald statistic for exclusion restrictions using a simple command following estimation of the general model. This capability makes it very easy to test multiple exclusion restrictions, provided the dimension of $(\mathbf{x}, \mathbf{z})$ is not so large as to make probit estimation difficult.

The likelihood ratio statistic is also easy to use. Let $\mathscr{L}_{u r}$ denote the value of the loglikelihood function from probit of $y$ on $\mathbf{x}$ and $\mathbf{z}$ (the unrestricted model), and let $\mathscr{L}_{r}$ denote the value of the likelihood function from probit of $y$ on $\mathbf{x}$ (the restricted model). Then the LR test of $\mathrm{H}_{0}: \boldsymbol{\gamma}=\mathbf{0}$ is simply $2\left(\mathscr{L}_{u r}-\mathscr{L}_{r}\right)$, which has an asymptotic $\chi_{Q}^{2}$ distribution under $\mathrm{H}_{0}$. This is analogous to the usual $F$ statistic in OLS analysis of a linear model.

The score or LM test is attractive if the unrestricted model is difficult to estimate. In this section, let $\hat{\boldsymbol{\beta}}$ denote the restricted estimator of $\boldsymbol{\beta}$, that is, the probit or logit estimator with $\mathbf{z}$ excluded from the model. The LM statistic using the estimated expected Hessian, $\hat{\mathbf{A}}_{i}$ (see equation (15.20) and Section 12.6.2), can be shown to be numerically identical to the following: (1) Define $\hat{u}_{i} \equiv y_{i}-G\left(\mathbf{x}_{i} \hat{\boldsymbol{\beta}}\right), \hat{G}_{i} \equiv G\left(\mathbf{x}_{i} \hat{\boldsymbol{\beta}}\right)$, and $\hat{g}_{i} \equiv g\left(\mathbf{x}_{i} \hat{\boldsymbol{\beta}}\right)$. These are all obtainable after estimating the model without $\mathbf{z}$. (2) Use all $N$ observations to run the auxiliary OLS regression


\begin{equation*}
\frac{\hat{u}_{i}}{\sqrt{\hat{G}_{i}\left(1-\hat{G}_{i}\right)}} \text { on } \frac{\hat{g}_{i}}{\sqrt{\hat{G}_{i}\left(1-\hat{G}_{i}\right)}} \mathbf{x}_{i}, \quad \frac{\hat{g}_{i}}{\sqrt{\hat{G}_{i}\left(1-\hat{G}_{i}\right)}} \mathbf{z}_{i} . \tag{15.22}
\end{equation*}


The $L M$ statistic is equal to the explained sum of squares from this regression. A test that is asymptotically (but not numerically) equivalent is $N R_{u}^{2}$, where $R_{u}^{2}$ is the uncentered $R$-squared from regression (15.22).

The LM procedure is rather easy to remember. The term $\hat{g}_{i} \mathbf{x}_{i}$ is the gradient of the mean function $G\left(\mathbf{x}_{i} \boldsymbol{\beta}+\mathbf{z}_{i} \gamma\right)$ with respect to $\boldsymbol{\beta}$, evaluated at $\boldsymbol{\beta}=\hat{\boldsymbol{\beta}}$ and $\gamma=\mathbf{0}$. Similarly, $\hat{g}_{i} \mathbf{z}_{i}$ is the gradient of $G\left(\mathbf{x}_{i} \boldsymbol{\beta}+\mathbf{z}_{i} \gamma\right)$ with respect to $\boldsymbol{\gamma}$, again evaluated at $\boldsymbol{\beta}=\hat{\boldsymbol{\beta}}$ and $\gamma=\mathbf{0}$. Finally, under $\mathrm{H}_{0}: \boldsymbol{\gamma}=\mathbf{0}$, the conditional variance of $u_{i}$ given ( $\mathbf{x}_{i}, \mathbf{z}_{i}$ ) is $G\left(\mathbf{x}_{i} \boldsymbol{\beta}\right)\left[1-G\left(\mathbf{x}_{i} \boldsymbol{\beta}\right)\right]$; therefore, $\left[\hat{G}_{i}\left(1-\hat{G}_{i}\right)\right]^{1 / 2}$ is an estimate of the conditional stan-\\
dard deviation of $u_{i}$. The dependent variable in regression (15.22) is often called a standardized residual because it is an estimate of $u_{i} /\left[G_{i}\left(1-G_{i}\right)\right]^{1 / 2}$, which has unit conditional (and unconditional) variance. The regressors are simply the gradient of the conditional mean function with respect to both sets of parameters, evaluated under $\mathrm{H}_{0}$, and weighted by the estimated inverse conditional standard deviation. The first set of regressors in regression (15.22) is $1 \times K$ and the second set is $1 \times Q$.

Under $\mathrm{H}_{0}, L M \sim \chi_{Q}^{2}$. The LM approach can be an attractive alternative to the $L R$ statistic if $\mathbf{z}$ has large dimension, since with many explanatory variables probit can be difficult to estimate.

\subsection*{15.5.2 Testing Nonlinear Hypotheses about $\boldsymbol{\beta}$}
For testing nonlinear restrictions on $\boldsymbol{\beta}$ in equation (15.8), the Wald statistic is computationally the easiest because the unrestricted estimator of $\boldsymbol{\beta}$, which is just probit or logit, is easy to obtain. Actually imposing nonlinear restrictions in estimationwhich is required to apply the score or likelihood ratio methods-can be difficult. However, we must also remember that the Wald statistic for testing nonlinear restrictions is not invariant to reparameterizations, whereas the $L M$ and $L R$ statistics are. (See Sections 12.6 and 13.6; for the $L M$ statistic, we would probably use the expected Hessian.)

Let the restictions on $\boldsymbol{\beta}$ be given by $\mathrm{H}_{0}: \mathbf{c}(\boldsymbol{\beta})=\mathbf{0}$, where $\mathbf{c}(\boldsymbol{\beta})$ is a $Q \times 1$ vector of possibly nonlinear functions satisfying the differentiability and rank requirements from Chapter 13. Then, from the general MLE analysis, the Wald statistic is simply


\begin{equation*}
W=\mathbf{c}(\hat{\boldsymbol{\beta}})^{\prime}\left[\nabla_{\beta} \mathbf{c}(\hat{\boldsymbol{\beta}}) \hat{\mathbf{V}} \nabla_{\beta} \mathbf{c}(\hat{\boldsymbol{\beta}})^{\prime}\right]^{-1} \mathbf{c}(\hat{\boldsymbol{\beta}}) \tag{15.23}
\end{equation*}


where $\hat{\mathbf{V}}$ is given in equation (15.20) and $\nabla_{\beta} \mathbf{c}(\hat{\boldsymbol{\beta}})$ is the $Q \times K$ Jacobian of $\mathbf{c}(\boldsymbol{\beta})$ evaluated at $\hat{\boldsymbol{\beta}}$.

\subsection*{15.5.3 Tests against More General Alternatives}
In addition to testing for omitted variables, sometimes we wish to test the probit or logit model against a more general functional form. When the alternatives are not standard binary response models, the Wald and $L R$ statistics are cumbersome to apply, whereas the LM approach is convenient because it only requires estimation of the null model.

As an example of a more complicated binary choice model, consider the latent variable model (15.9) but assume that $e \mid \mathbf{x} \sim \operatorname{Normal}\left[0, \exp \left(2 \mathbf{x}_{1} \boldsymbol{\delta}\right)\right]$, where $\mathbf{x}_{1}$ is $1 \times K_{1}$ subset of $\mathbf{x}$ that excludes a constant and $\boldsymbol{\delta}$ is a $K_{1} \times 1$ vector of additional parameters. (In many cases we would take $\mathbf{x}_{1}$ to be all nonconstant elements of $\mathbf{x}$.) Therefore, there is heteroskedasticity in the latent variable model, so that $e$ is no\\
longer independent of $\mathbf{x}$. The standard deviation of $e$ given $\mathbf{x}$ is simply $\exp \left(\mathbf{x}_{1} \boldsymbol{\delta}\right)$. Define $r=e / \exp \left(\mathbf{x}_{1} \boldsymbol{\delta}\right)$, so that $r$ is independent of $\mathbf{x}$ with a standard normal distribution. Then


\begin{align*}
\mathrm{P}(y=1 \mid \mathbf{x}) & =\mathrm{P}(e>-\mathbf{x} \boldsymbol{\beta} \mid \mathbf{x})=\mathrm{P}\left[\exp \left(-\mathbf{x}_{1} \boldsymbol{\delta}\right) e>-\exp \left(-\mathbf{x}_{1} \boldsymbol{\delta}\right) \mathbf{x} \boldsymbol{\beta}\right] \\
& =\mathrm{P}\left[r>-\exp \left(-\mathbf{x}_{1} \boldsymbol{\delta}\right) \mathbf{x} \boldsymbol{\beta}\right]=\Phi\left[\exp \left(-\mathbf{x}_{1} \boldsymbol{\delta}\right) \mathbf{x} \boldsymbol{\beta}\right] \tag{15.24}
\end{align*}


The partial effects of $x_{j}$ on $\mathrm{P}(y=1 \mid \mathbf{x})$ are much more complicated in equation (15.24) than in equation (15.8). When $\boldsymbol{\delta}=\mathbf{0}$, we obtain the standard probit model. Therefore, a test of the probit functional form for the response probability is a test of $\mathrm{H}_{0}: \boldsymbol{\delta}=\mathbf{0}$.

To obtain the LM test of $\boldsymbol{\delta}=\mathbf{0}$ in equation (15.24), it is useful to derive the LM test for an index model against a more general alternative. Consider\\
$\mathrm{P}(y=1 \mid \mathbf{x})=m(\mathbf{x} \boldsymbol{\beta}, \mathbf{x}, \boldsymbol{\delta})$,\\
where $\boldsymbol{\delta}$ is a $Q \times 1$ vector of parameters. We wish to test $\mathrm{H}_{0}: \boldsymbol{\delta}=\boldsymbol{\delta}_{0}$, where $\boldsymbol{\delta}_{0}$ is often (but not always) a vector of zeros. We assume that, under the null, we obtain a standard index model (probit or logit, usually):\\
$G(\mathbf{x} \boldsymbol{\beta})=m\left(\mathbf{x} \boldsymbol{\beta}, \mathbf{x}, \boldsymbol{\delta}_{0}\right)$.

In the previous example, $G(\cdot)=\Phi(\cdot), \boldsymbol{\delta}_{0}=\mathbf{0}$, and $m(\mathbf{x} \boldsymbol{\beta}, \mathbf{x}, \boldsymbol{\delta})=\Phi\left[\exp \left(-\mathbf{x}_{1} \boldsymbol{\delta}\right) \mathbf{x} \boldsymbol{\beta}\right]$.\\
Let $\hat{\boldsymbol{\beta}}$ be the probit or logit estimator of $\boldsymbol{\beta}$ obtained under $\boldsymbol{\delta}=\boldsymbol{\delta}_{0}$. Define $\hat{u}_{i} \equiv y_{i}-G\left(\mathbf{x}_{i} \hat{\boldsymbol{\beta}}\right), \hat{G}_{i} \equiv G\left(\mathbf{x}_{i} \hat{\boldsymbol{\beta}}\right)$, and $\hat{g}_{i} \equiv g\left(\mathbf{x}_{i} \hat{\boldsymbol{\beta}}\right)$. The gradient of the mean function $m\left(\mathbf{x}_{i} \boldsymbol{\beta}, \mathbf{x}_{i}, \boldsymbol{\delta}\right)$ with respect to $\boldsymbol{\beta}$, evaluated at $\boldsymbol{\delta}_{0}$, is simply $g\left(\mathbf{x}_{i} \boldsymbol{\beta}\right) \mathbf{x}_{i}$. The only other piece we need is the gradient of $m\left(\mathbf{x}_{i} \boldsymbol{\beta}, \mathbf{x}_{i}, \boldsymbol{\delta}\right)$ with respect to $\boldsymbol{\delta}$, evaluated at $\boldsymbol{\delta}_{0}$. Denote this $1 \times Q$ vector as $\nabla_{\delta} m\left(\mathbf{x}_{i} \boldsymbol{\beta}, \mathbf{x}_{i}, \boldsymbol{\delta}_{0}\right)$. Further, set $\nabla_{\delta} \hat{m}_{i} \equiv \nabla_{\delta} m\left(\mathbf{x}_{i} \hat{\boldsymbol{\beta}}, \mathbf{x}_{i}, \boldsymbol{\delta}_{0}\right)$. The $L M$ statistic can be obtained as the explained sum of squares from the regression\\
$\frac{\hat{u}_{i}}{\sqrt{\hat{G}_{i}\left(1-\hat{G}_{i}\right)}}$ on $\frac{\hat{g}_{i}}{\sqrt{\hat{G}_{i}\left(1-\hat{G}_{i}\right)}} \mathbf{x}_{i}, \quad \frac{\nabla_{\delta} \hat{m}_{i}}{\sqrt{\hat{G}_{i}\left(1-\hat{G}_{i}\right)}}$,\\
which is quite similar to regression (15.22). The null distribution of the $L M$ statistic is $\chi_{Q}^{2}$, where $Q$ is the dimension of $\boldsymbol{\delta}$. An asymptotically equivalent statistic is $N R_{u}^{2}$.

When applying this test to the preceding probit example, we have only $\nabla_{\delta} \hat{m}_{i}$ left to compute. But $m\left(\mathbf{x}_{i} \boldsymbol{\beta}, \mathbf{x}_{i}, \boldsymbol{\delta}\right)=\Phi\left[\exp \left(-\mathbf{x}_{i 1} \boldsymbol{\delta}\right) \mathbf{x}_{i} \boldsymbol{\beta}\right]$, and so\\
$\nabla_{\delta} m\left(\mathbf{x}_{i} \boldsymbol{\beta}, \mathbf{x}_{i}, \boldsymbol{\delta}\right)=-\left(\mathbf{x}_{i} \boldsymbol{\beta}\right) \exp \left(-\mathbf{x}_{i 1} \boldsymbol{\delta}\right) \mathbf{x}_{i 1} \phi\left[\exp \left(-\mathbf{x}_{i 1} \boldsymbol{\delta}\right) \mathbf{x}_{i} \boldsymbol{\beta}\right]$.\\
When evaluated at $\boldsymbol{\beta}=\hat{\boldsymbol{\beta}}$ and $\boldsymbol{\delta}=\mathbf{0}$ (the null value), we get $\nabla_{\delta} \hat{m}_{i}=-\left(\mathbf{x}_{i} \hat{\boldsymbol{\beta}}\right) \phi\left(\mathbf{x}_{i} \hat{\boldsymbol{\beta}}\right) \mathbf{x}_{i 1} \equiv-\left(\mathbf{x}_{i} \hat{\boldsymbol{\beta}}\right) \hat{\phi}_{i} \mathbf{x}_{i 1}$, a $1 \times K_{1}$ vector. Regression (15.27) becomes


\begin{equation*}
\frac{\hat{u}_{i}}{\sqrt{\hat{\Phi}_{i}\left(1-\hat{\Phi}_{i}\right)}} \text { on } \frac{\hat{\phi}_{i}}{\sqrt{\hat{\Phi}_{i}\left(1-\hat{\Phi}_{i}\right)}} \mathbf{x}_{i}, \quad \frac{\left(\mathbf{x}_{i} \hat{\boldsymbol{\beta}}\right) \hat{\phi}_{i}}{\sqrt{\hat{\Phi}_{i}\left(1-\hat{\Phi}_{i}\right)}} \mathbf{x}_{i 1} \tag{15.28}
\end{equation*}


(We drop the minus sign because it does not affect the value of the explained sum of squares or $R_{u}^{2}$.) Under the null hypothesis that the probit model is correctly specified, $L M \sim \chi_{K_{1}}^{2}$. This statistic is easy to compute after estimation by probit.

For a one-degree-of-freedom test regardless of the dimension of $\mathbf{x}_{i}$, replace the last term in regression (15.28) with $\left(\mathbf{x}_{i} \hat{\boldsymbol{\beta}}\right)^{2} \hat{\phi}_{i} / \sqrt{\hat{\Phi}_{i}\left(1-\hat{\Phi}_{i}\right)}$, and then the explained sum of squares is distributed asymptotically as $\chi_{1}^{2}$. See Davidson and MacKinnon (1984) for further examples.

As we discussed briefly in Chapters 12 and 13, variable addition versions of specification tests can be somewhat easier to compute. Rather than running the regression in equation (15.28), we can instead estimate the auxiliary probit model with response probability $\Phi\left[\mathbf{x}_{i} \boldsymbol{\beta}+\left(\mathbf{x}_{i} \hat{\boldsymbol{\beta}}\right) \mathbf{x}_{i 1} \boldsymbol{\delta}\right]$ and test $\mathrm{H}_{0}: \boldsymbol{\delta}=\mathbf{0}$ using a standard Wald test. Again, $\hat{\boldsymbol{\beta}}$ denotes the original probit estimates from probit of $y_{i}$ on $\mathbf{x}_{i}$. The specification test is obtained from probit of $y_{i}$ on $\mathbf{x}_{i},\left(\mathbf{x}_{i} \hat{\boldsymbol{\beta}}\right) \mathbf{x}_{i 1}$ and testing the last $Q$ terms for joint significance. It is easy to see that the variable addition test (VAT) is asymptotically equivalent to the score test, and the VAT circumvents the need to explicitly obtain the weighted residuals and gradients. The VAT approach also makes it clear that the score test for heteroskedasticity in the error from the latent model is indistinguishable from testing for a particular kind of interactive effect in the covariates. We might have arrived at such a test directly, without modeling $\operatorname{Var}(e \mid \mathbf{x})$.

\subsection*{15.6 Reporting the Results for Probit and Logit}
Several statistics should be reported routinely in any probit or logit (or other binary choice) analysis. The $\hat{\beta}_{j}$, their standard errors, and the value of the likelihood function are reported by all software packages that do binary response analysis. The $\hat{\beta}_{j}$ give the signs of the partial effects of each $x_{j}$ on the response probability, and the statistical significance of $x_{j}$ is determined by whether we can reject $\mathrm{H}_{0}: \beta_{j}=0$.

One measure of goodness of fit that is sometimes reported is the percent correctly predicted. The easiest way to describe this statistic is to define a binary predictor of $y_{i}$ to be one if the predicted probability is at least .5 , and zero otherwise. More precisely, define the binary variable $\tilde{y}_{i}=1$ if $G\left(\mathbf{x}_{i} \hat{\boldsymbol{\beta}}\right) \geq .5$ and $\tilde{y}_{i}=0$ if $G\left(\mathbf{x}_{i} \hat{\boldsymbol{\beta}}\right)<.5$. Given $\left\{\tilde{y}_{i}: i=1,2, \ldots, N\right\}$, we can see how well $\tilde{y}_{i}$ predicts $y_{i}$ across all observations. There are four possible outcomes on each pair, $\left(y_{i}, \tilde{y}_{i}\right)$; when both are zero or both are one, we make the correct prediction. In the two cases where one of the pair is zero\\
and the other is one, we make the incorrect prediction. The percent correctly predicted is the percent of times that $\tilde{y}_{i}=y_{i}$. (This goodness-of-fit measure can be computed for the linear probability model, too.)

While the percent correctly predicted is useful as a goodness-of-fit measure, it can be misleading. In particular, it is possible to get rather high percentages correctly predicted even when the least likely outcome is very poorly predicted. For example, suppose that $N=200,160$ observations have $y_{i}=0$, and, out of these 160 observations, 140 of the $\tilde{y}_{i}$ are also zero (so we correctly predict 87.5 percent of the zero outcomes.) Even if none of the predictions is correct when $y_{i}=1$, we still correctly predict $70 \%$ of all outcomes $(140 / 200=.70)$. Often, we hope be able to predict the least likely outcome (such as whether someone is arrested for committing a crime), and we only know how well we do that by obtaining the percent correctly predicted for each outcome. It is easily shown that the overall percent correctly predicted is the weighted average of the percent correctly predicted for $y=0$ and $y=1$, with the weights being the fraction of zero and one outcomes on $y$, respectively.

Some have criticized the prediction rule described above for always using a threshold value of .5, especially when one of the outcomes is unlikely. For example, if $\bar{y}=.08$ (only eight percent "successes" in the sample), it could be that we never predict $y_{i}=1$ because the estimated probability of success is never greater than .5 . One alternative is to use the fraction of successes in the sample as the threshold - .08 in the previous example. In other words, define $\tilde{y}_{i}=1$ when $G\left(\mathbf{x}_{i} \hat{\boldsymbol{\beta}}\right) \geq .08$, and zero otherwise. Using this rule will certainly increase the number of predicted successes, but not without cost: we necessarily make more mistakes-perhaps many more-in predicting the zero outcomes ("failures"). In terms of the overall percent correctly predicted, we may actually do worse than when using the traditional .5 threshold.

A third possibility is to choose the threshold such that the fraction of $\tilde{y}_{i}=1$ in the sample is the same (or very close) to $\bar{y}$. In other words, search over threshold values $\tau, 0<\tau<1$, such that if we define $\tilde{y}_{i}=1$ when $G\left(\mathbf{x}_{i} \hat{\boldsymbol{\beta}}\right) \geq \tau$, then $\sum_{i=1}^{n} \tilde{y}_{i} \approx \sum_{i=1}^{n} y_{i}$. (The trial-and-error effort required to find the desired value of $\tau$ can be tedious, but it is feasible. In some cases, it will not be possible to make the number of predicted successes exactly the same as the number of successes in the sample.) Now, given the set of binary predictors $\tilde{y}_{i}$, we can compute the percent correctly predicted for each of the two outcomes, as well as the overall percent correctly predicted.

Various pseudo-R-squared measures have been proposed for binary response. McFadden (1974) suggests the measure $1-\mathscr{L}_{u r} / \mathscr{L}_{\mathrm{O}}$, where $\mathscr{L}_{u r}$ is the log-likelihood function for the estimated model and $\mathscr{L}_{\mathrm{o}}$ is the log-likelihood function in the model with only an intercept. Because the log likelihood for a binary response model is always negative, $\left|\mathscr{L}_{u r}\right| \leq\left|\mathscr{L}_{\mathrm{o}}\right|$, and so the pseudo- $R$-squared is always between zero\\
and one. Alternatively, we can use a sum of squared residuals measure: $1-\mathrm{SSR}_{u r} / \mathrm{SSR}_{\mathrm{o}}$, where $\mathrm{SSR}_{u r}$ is the sum of squared residuals $\hat{u}_{i}=y_{i}-G\left(\mathbf{x}_{i} \hat{\boldsymbol{\beta}}\right)$ and $\mathrm{SSR}_{\mathrm{o}}$ is the total sum of squares of $y_{i}$. Several other measures have been suggested (see, for example, Maddala, 1983, Chap. 2), but goodness of fit is not as important as statistical and economic significance of the explanatory variables. Estrella (1998) contains a recent comparison of goodness-of-fit measures for binary response.

Usually we want to estimate the effects of the variables $x_{j}$ on the response probabilities $\mathrm{P}(y=1 \mid \mathbf{x})$. If $x_{j}$ is (roughly) continuous, then


\begin{equation*}
\Delta \mathrm{P}(\widehat{y=} 1 \mid \mathbf{x}) \approx\left[g(\mathbf{x} \hat{\boldsymbol{\beta}}) \hat{\beta}_{j}\right] \Delta x_{j} \tag{15.29}
\end{equation*}


for small changes in $x_{j}$. (As usual when using calculus, the notion of "small" here is somewhat vague.) Therefore, the estimated partial effect of a continuous variable on the response probability, evaluated at $\mathbf{x}$, is $g(\mathbf{x} \hat{\boldsymbol{\beta}}) \hat{\beta}_{j}$. Because $g(\mathbf{x} \hat{\boldsymbol{\beta}})$ depends on $\mathbf{x}$, we need to decide which partial effects to report. We could report this scale factor at, say, medians of the explanatory variables, or at different quantiles. For summarizing the magnitudes of the effects, it is useful to have a single scale factor that can be used to multiply the coefficients on (roughly) continuous variables. Often the sample averages of the $x_{j}$ are plugged in to get $g(\overline{\mathbf{x}} \hat{\boldsymbol{\beta}})$, with $\bar{x}_{1} \equiv 1$ because we include a constant. We call the resulting partial effect the partial effect at the average (PEA).

The PEA is easy to compute, but it does have drawbacks. First, it need not represent the partial effect for any particular unit in the population. That is, the average of the explanatory variables may not, in any sensible way, represent the average unit in the population. One issue is that, if $\mathbf{x}$ contains nonlinear functions of underlying variables, such as logarithms, we must decide whether to use the average of the nonlinear function or the nonlinear function of the average. The latter has some appeal-for example, if $\log ($ inc $)$ is an explanatory variable, evaluate $\log ($ inc $)$ at $\overline{\text { inc }}$, rather than use $\overline{\log (\text { inc })}$ - but software packages that support PEA estimation (such as Stata, with its $m f x$, for "marginal effects," command) necessarily use the average of the nonlinear functions (because one must create the nonlinear functions before including them in logit or probit).

If two or more elements of $\mathbf{x}$ are functionally related, such as quadratics or interactions, it is not even clear what the PEAs of individual coefficients mean. For example, suppose $x_{K-1}=$ age and $x_{K}=a g e^{2}$. Then the reported PEAs for age and age $^{2}$ are $g(\overline{\mathbf{x}} \hat{\boldsymbol{\beta}}) \hat{\beta}_{K-1}$ and $g(\overline{\mathbf{x}} \hat{\boldsymbol{\beta}}) \hat{\beta}_{K}$, respectively, where $\overline{\mathbf{x}}=\left(1, \bar{x}_{2}, \ldots, \bar{x}_{K-2}, \overline{\text { age }}, \overline{\text { age }}{ }^{2}\right)$. These PEAs do not tell us what we want to know about the partial effect of age on $\mathrm{P}(y=1 \mid \mathbf{x})$. For any $\mathbf{x}$, the estimated partial effect is $g(\mathbf{x} \hat{\boldsymbol{\beta}})\left(\hat{\beta}_{K-1}+2 \hat{\beta}_{K}\right.$ age $)$. Now, we might be interested in evaluating this partial effect at the mean values, but that would entail using $\overline{a g e}^{2}$, rather than $\overline{a g e^{2}}$, inside $g(\cdot)$. If we are really interested in the\\
effect of age on the response probability, we might want to evaluate the partial effect at several different values of age, perhaps evaluating the other explanatory variables at their means. If $\hat{\beta}_{K-1}$ and $\hat{\beta}_{K}$ have different signs, the estimated turning point in the relationship is just as in models linear in the parameters: $\widehat{\operatorname{age}}^{*}=\left|\hat{\beta}_{K-1} /\left(2 \hat{\beta}_{K}\right)\right|$. The usual turning point calculation works because $g(\mathbf{x} \hat{\boldsymbol{\beta}})>0$ for all $\mathbf{x}$. Similar care must be used for obtaining partial effects when interactions are included in $\mathbf{x}$.

For discrete variables, it is well known that the average need not even be a possible outcome of the variable. If, say, $x_{2}$ is a binary variable, $\bar{x}_{2}$ is the fraction of ones in the sample, and therefore cannot correspond to a particular unit (which must have zero or one). For example, if $x_{2}=$ female is a gender dummy, then the PEA is the partial effect when female is replaced with the fraction of women in the sample. One way to overcome this conceptual problem is to compute partial effects separately for $x_{2}=1$ and $x_{2}=0$, but then we no longer have a single number to report as the partial effect. Similar comments hold for other discrete variables, such as number of children. If the average is 1.5 , plugging this value into $g(\cdot)$ does not produce the partial effect for any particular family.

To obtain standard errors of the partial effects in equation (15.29) we can use the delta method. Consider the case $j=K$ for notational simplicity, and for given $\mathbf{x}$, define $\delta_{K}=\beta_{K} g(\mathbf{x} \boldsymbol{\beta})=\partial \mathbf{P}(y=1 \mid \mathbf{x}) / \partial x_{K}$. Write this relation as $\delta_{K}=h(\boldsymbol{\beta})$ to denote that this is a (nonlinear) function of the vector $\boldsymbol{\beta}$. We assume $x_{1}=1$. The gradient of $h(\boldsymbol{\beta})$ is\\
$\nabla_{\beta} h(\boldsymbol{\beta})=\left[\beta_{K} \frac{\mathrm{~d} g}{\mathrm{~d} z}(\mathbf{x} \boldsymbol{\beta}), \beta_{K} x_{2} \frac{\mathrm{~d} g}{\mathrm{~d} z}(\mathbf{x} \boldsymbol{\beta}), \ldots, \beta_{K} x_{K-1} \frac{\mathrm{~d} g}{\mathrm{~d} z}(\mathbf{x} \boldsymbol{\beta}), \beta_{K} x_{K} \frac{\mathrm{~d} g}{\mathrm{~d} z}(\mathbf{x} \boldsymbol{\beta})+g(\mathbf{x} \boldsymbol{\beta})\right]$,\\
where $\mathrm{d} g / \mathrm{d} z$ is simply the derivative of $g$ with respect to its argument. The delta method implies that the asymptotic variance of $\hat{\delta}_{K}$ is estimated as\\
$\left[\nabla_{\beta} h(\hat{\boldsymbol{\beta}})\right] \hat{\mathbf{V}}\left[\nabla_{\beta} h(\hat{\boldsymbol{\beta}})\right]^{\prime}$,\\
where $\hat{\mathbf{V}}$ is the asymptotic variance estimate of $\hat{\boldsymbol{\beta}}$. The asymptotic standard error of $\hat{\delta}_{K}$ is simply the square root of expression (15.30). This calculation allows us to obtain a large-sample confidence interval for $\hat{\delta}_{K}$. The program Stata does this calculation for logit and probit using the $m f x$ command. Alternatively, we can apply the bootstrap as discussed in Chapter 12.

If $x_{K}$ is a discrete variable, then we can estimate the change in the predicted probability in going from $c_{K}$ to $c_{K}+1$ as


\begin{align*}
\hat{\delta}_{K}= & G\left[\hat{\beta}_{1}+\hat{\beta}_{2} \bar{x}_{2}+\cdots+\hat{\beta}_{K-1} \bar{x}_{K-1}+\hat{\beta}_{K}\left(c_{K}+1\right)\right] \\
& -G\left(\hat{\beta}_{1}+\hat{\beta}_{2} \bar{x}_{2}+\cdots+\hat{\beta}_{K-1} \bar{x}_{K-1}+\hat{\beta}_{K} c_{K}\right) \tag{15.31}
\end{align*}


In particular, when $x_{K}$ is a binary variable, set $c_{K}=0$. Of course, the other $x_{j}$ 's can be evaluated anywhere, but the use of sample averages is typical. The delta method can be used to obtain a standard error of equation (15.31). For probit, Stata does this calculation when $x_{K}$ is a binary variable. Usually the calculations ignore the fact that $\bar{x}_{j}$ is an estimate of $\mathrm{E}\left(x_{j}\right)$ in applying the delta method. If we are truly interested in $\beta_{K} g\left(\boldsymbol{\mu}_{x} \boldsymbol{\beta}\right)$, the estimation error in $\overline{\mathbf{x}}$ can be accounted for, but it makes the calculation more complicated, and it is unlikely to have a large effect. The bootstrap would properly account for the sampling variation in $\overline{\mathbf{x}}$ when $\overline{\mathbf{x}}$ is recomputed for each bootstrap replication.

An alternative way to summarize the estimated marginal effects is to estimate the average value of $\beta_{K} g(\mathbf{x} \boldsymbol{\beta})$ across the population, or $\beta_{K} \mathrm{E}[g(\mathbf{x} \boldsymbol{\beta})]$. This quantity is the average partial effect (APE) that we discussed generally in Chapter 2, and also for linear models with random coefficients in Chapter 4. Here we are averaging across the distribution of all observable covariates. A consistent estimator of the APE is


\begin{equation*}
\hat{\beta}_{K}\left[N^{-1} \sum_{i=1}^{N} g\left(\mathbf{x}_{i} \hat{\boldsymbol{\beta}}\right)\right] \tag{15.32}
\end{equation*}


when $x_{K}$ is continuous or


\begin{equation*}
N^{-1} \sum_{i=1}^{N}\left[G\left(\hat{\beta}_{1}+\hat{\beta}_{2} x_{i 2}+\cdots+\hat{\beta}_{K-1} x_{i, K-1}+\hat{\beta}_{K}\right)-G\left(\hat{\beta}_{1}+\hat{\beta}_{2} x_{i 2}+\cdots+\hat{\beta}_{K-1} x_{i, K-1}\right)\right] \tag{15.33}
\end{equation*}


when $x_{K}$ is binary. Naturally, we can use equation (15.33) to estimate the APE from changing $x_{K}$ (continuous or discrete) from any two values, say $c_{K}^{0}$ to $c_{K}^{1}$. Then $\hat{\beta}_{K}$ is replaced with $\hat{\beta}_{K} c_{K}^{1}$ in the first term and we insert $\hat{\beta}_{K} c_{K}^{0}$ into the second occurrence of $G(\cdot)$. Further, in either equation (15.32) or (15.33), we can fix some of the explanatory variables at specific values and average across the remaining ones. For example, suppose that in a logit or probit model of employment, $x_{K}$ is a job training binary indicator and $x_{K-1}$ is a race indicator-say, one for nonwhite, zero for white. Then, rather than compute (15.33), we can set $x_{K-1}=1$ and compute the APE for nonwhites, and then set $x_{K-1}=0$ to obtain the APE for whites.

As in the case of the PEA, if some elements of $\mathbf{x}$ are functions of each other, obtaining APEs of the form in equation (15.32) is not especially useful. If, say, $x_{K-1}=$ age and $x_{K}=a g e^{2}$, we can estimate the APE of age by averaging the individual partial effects, $\left(\hat{\beta}_{K-1}+2 \hat{\beta}_{K} a g e_{i}\right) \times g\left(\mathbf{x}_{i} \hat{\boldsymbol{\beta}}\right)$, across $i$. Again, it probably makes more sense to evaluate the partial effect at different values of age and then to average these across the other variables, say, $N^{-1} \sum_{i=1}^{N}\left(\hat{\beta}_{K-1}+2 \hat{\beta}_{K} a g e^{o}\right) \times g\left(\hat{\beta}_{1}+\hat{\beta}_{2} x_{i 2}\right. +\cdots+\hat{\beta}_{K-2} x_{i, K-2}+\hat{\beta}_{K-1}$ age $^{o}+\hat{\beta}_{K}\left(\text { age }^{o}\right)^{2}$ ) for a given value age $^{o}$.

Generally, we should not expect the scale factors for the PEAs and APEs to be similar. That is, we should not expect $g(\overline{\mathbf{x}} \hat{\boldsymbol{\beta}})$ to be similar to $N^{-1} \sum_{i=1}^{N} g\left(\mathbf{x}_{i} \hat{\boldsymbol{\beta}}\right)$. The reason is simple: the average of a nonlinear function (the APE scale) is not the nonlinear function of the average (the PEA scale). In the population the scales are not the same either, because the expected value does not pass through nonlinear functions: $g[\mathrm{E}(\mathbf{x} \boldsymbol{\beta})] \neq \mathrm{E}[g(\mathbf{x} \boldsymbol{\beta})]$.

Equation (15.33) has a nice interpretation for policy analysis, where $x_{K}$ is the binary policy indicator (equal to one when the policy is in effect, zero otherwise). We can view each summand in (15.33) as follows. Regardless of whether unit $i$ was subject to the policy, we can obtain the predicted probability in each regime. The term in brackets in equation (15.33) is the difference in the estimated probability that $y_{i}=1$ with and without participation. In other words, we compute the (counterfactual) effect of the policy for each cross section unit $i$, and then average these differences across all $i$. This gives us an estimate of the average effect of the policy. We will have much more to say about a counterfactual framework for policy analysis in Chapter 21 and the estimation of average treatment effects, of which equation (15.33) is an example.

We can obtain standard errors of the estimators in (15.32) and (15.33) by applying the bootstrap to an average, as described in Section 12.8.2, or we can use the delta method; see Problem 15.15.

Unlike the magnitudes of the coefficient estimates, the APEs (and, to a lesser extent, the PEAs) can be compared across models. It is rare to find, say, probit and logit estimates having different signs unless the coefficients are estimated imprecisely. But logit, probit, and the LPM implicitly use different scale factors. The scale for the LPM is unity. For logit and probit, we can look at the functions $g(\cdot)$ at zero to obtain a rough idea of how the $\hat{\beta}_{j}$-at least the slopes on the continuous variables-may differ. For probit, $g(0) \approx .4$ and for logit $g(0)=.25$. Therefore, we expect the logit slope coefficients to have the largest magnitude, followed by the probit estimates, followed by the LPM estimates. In fact, sometimes some (crude) rules of thumb are adopted: multiply the probit coefficients by 1.6 and the LPM coefficients by 4 to make them roughly comparable to the logit estimates. The LPM estimates are multiplied by 2.5 to make them comparable to the probit estimates. If $\overline{\mathbf{x}} \hat{\boldsymbol{\beta}}$ is close to zero, then $g(\overline{\mathbf{x}} \hat{\boldsymbol{\beta}})$ will be close to .25 for logit and close to .4 for probit, and then the rules of thumb for adjusting coefficients can be justified based on the PEAs. (When one of the outcomes on $y$ is much less likely than the other, it is unlikely that $\overline{\mathbf{x}} \hat{\boldsymbol{\beta}}$ will be close to zero.) Generally, it is a good idea to compute the PEAs and APEs, as well as partial effects at other interesting values of $\mathbf{x}$, or at values that allow us to determine policy effects for different groups in the population.

As a general rule, the estimates from an LPM estimation are more comparable to the APEs than to the PEAs. As mentioned above, the PEAs need not come close to representing an individual in the population. On the other hand, the APE is the partial effect averaged across the population and, in one case, we can show that the LPM estimates are consistent for the APEs regardless of the actual function $G(\cdot)$. In fact, using a result of Stoker (1986), we can say much more, for any kind of response variable $y$. Let $m(\mathbf{x})=\mathrm{E}(y \mid \mathbf{x})$ be the conditional mean, so that $\lambda=\mathrm{E}\left[\nabla_{\mathbf{x}} m(\mathbf{x})^{\prime}\right]$ is the vector of APEs, where now $\mathbf{x}$ just includes continuous covariates and no elements are functionally related. Under assumptions about the distribution of $\mathbf{x}$ which are generally satisified if $\mathbf{x}$ has convex, unbounded support-Stoker shows that $\mathrm{E}\left[\nabla_{\mathbf{x}} m(\mathbf{x})^{\prime}\right]=-\mathrm{E}\left[m(\mathbf{x}) \times \nabla_{\mathbf{x}} \log f(\mathbf{x})^{\prime}\right]=-\mathrm{E}\left[y \times \nabla_{\mathbf{x}} \log f(\mathbf{x})^{\prime}\right]$ (where the second equality follows by iterated expectations), and so the vector of APEs is simply $-\mathrm{E}\left[y \times \nabla_{\mathbf{x}} \log f(\mathbf{x})^{\prime}\right]$. Now, if $\mathbf{x} \sim \operatorname{Normal}\left(\boldsymbol{\mu}_{\mathbf{x}}, \boldsymbol{\Sigma}_{\mathbf{x}}\right)$-that is, if $\mathbf{x}$ has a multivariate normal distribution-then it is easily shown that $\nabla_{\mathbf{x}} \log f(\mathbf{x})=-\left(\mathbf{x}-\boldsymbol{\mu}_{\mathbf{x}}\right) \boldsymbol{\Sigma}_{\mathbf{x}}^{-1}$. It follows by substitution that\\
$\lambda=\mathrm{E}\left[y \times \boldsymbol{\Sigma}_{\mathbf{x}}^{-1}\left(\mathbf{x}-\boldsymbol{\mu}_{\mathbf{x}}\right)^{\prime}\right]=\left\{\mathrm{E}\left[\left(\mathbf{x}-\boldsymbol{\mu}_{\mathbf{x}}\right)^{\prime}\left(\mathbf{x}-\boldsymbol{\mu}_{\mathbf{x}}\right)\right]\right\}^{-1} \mathrm{E}\left[\left(\mathbf{x}-\boldsymbol{\mu}_{\mathbf{x}}\right)^{\prime} y\right]$,\\
which is simply the vector of slope coefficients from the linear projection of $y$ on $1, \mathbf{x}$. This calculation demonstrates that, regardless of the conditional mean function $m(\mathbf{x})$, if $\mathbf{x}$ has a multivariate normal distribution, then a linear regression consistently estimates the APEs-or, in Stoker's (1986) terminology, the average derivatives.

The general result applies to binary responses. In fact, there is no reason to even specify the model in the index form $\mathrm{P}(y=1 \mid \mathbf{x})=G(\alpha+\mathbf{x} \boldsymbol{\beta})$. But, in this case-for continuously differentiable functions $G(\cdot)$-the vector of average partial effects is $\mathrm{E}[g(\alpha+\mathbf{x} \boldsymbol{\beta})] \boldsymbol{\beta}$, and we can conclude that $\lambda$ estimates $\boldsymbol{\beta}$ up to the positive scale factor $\mathrm{E}[g(\alpha+\mathbf{x} \boldsymbol{\beta})]$. We discuss results of this sort further in Section 15.7.6.

The result we just derived is of limited practical value because the assumption that $\mathbf{x}$ is multivariate normal is very restrictive. In particular, it rules out discrete variables and continuous variables with asymmetric distributions. In the context of the index model $\mathrm{P}(y=1 \mid \mathbf{x})=G(\alpha+\mathbf{x} \boldsymbol{\beta})$, assuming that all elements of $\mathbf{x}$ are normally distributed is very restrictive. (For example, it rules out some staples of empirical work, such as quadratics and interactions, in addition to discrete explanatory variables.) Nevertheless, the result does show that estimation of an LPM can consistently estimate average partial effects, and it may be a good estimator for moderate deviations of $\mathbf{x}$ from normality.

Example 15.2 (Married Women's Labor Force Participation): We now estimate logit and probit models for women's labor force participation. For comparison, we

\begin{table}[h]
\begin{center}
\captionsetup{labelformat=empty}
\caption{Table 15.1\\
LPM, Logit, and Probit Estimates of Labor Force Participation}
\begin{tabular}{|l|l|l|l|}
\hline
\multicolumn{4}{|l|}{Dependent Variable: inlf} \\
\hline
Independent Variable & LPM (OLS) & Logit (MLE) & Probit (MLE) \\
\hline
nwifeinc & -. 0034 (.0015) & -. 021 (.008) & -. 012 (.005) \\
\hline
educ & . 038 (.007) & . 221 (.043) & . 131 (.025) \\
\hline
exper & . 039 (.006) & . 206 (.032) & . 123 (.019) \\
\hline
exper ${ }^{2}$ & -. 00060 (.00019) & -. 0032 (.0010) & -. 0019 (.0006) \\
\hline
age & -. 016 (.002) & -. 088 (.015) & -. 053 (.008) \\
\hline
kidslt6 & -. 262 (.032) & -1.443 (0.204) & -. 868 (.119) \\
\hline
kidsge6 & . 013 (.013) & . 060 (.075) & . 036 (.043) \\
\hline
constant & . 586 (.151) & . 425 (.860) & . 270 (.509) \\
\hline
Number of observations & 753 & 753 & 753 \\
\hline
Percent correctly predicted & 73.4 & 73.6 & 73.4 \\
\hline
Log-likelihood value & - & -401.77 & -401.30 \\
\hline
Pseudo- $R$-squared & . 264 & . 220 & . 221 \\
\hline
\end{tabular}
\end{center}
\end{table}

report the linear probability estimates. The results, with standard errors in parentheses, are given in Table 15.1 (for the LPM, these are heteroskedasticity-robust).

The estimates from the three models tell a consistent story. The signs of the coefficients are the same across models, and the same variables are statistically significant in each model. The pseudo- $R$-squared for the LPM is just the usual $R$-squared reported for OLS; for logit and probit the pseudo- $R$-squared is the measure based on the log likelihoods described previously. In terms of overall percent correctly predicted, the models do equally well. For the probit model, it correctly predicts "out of the labor force" about 63.1 percent of the time, and it correctly predicts "in the labor force" about 81.3 percent of the time. The LPM has the same overall percent correctly predicted, but there are slight differences within each outcome.

As we emphasized earlier, the magnitudes of the coefficients are not directly comparable across the models, although the ratios of coefficients on the (roughly) continuous explanatory variables are. For example, the ratios of the nwifeinc and educ\\
coefficients are about $-.0895,-.0950$, and -.0916 for the LPM, logit, and probit models, respectively.

If we evaluate the standard normal probability density function, $\phi\left(\hat{\beta}_{0}+\hat{\beta}_{1} x_{1}+\cdots\right. +\hat{\beta}_{k} x_{k}$ ), at the average values of the independent variables in the sample (including the average of exper ${ }^{2}$ ), we obtain about .391 . (This value is close to .4 because the outcomes on inlf are fairly balanced: about $56.7 \%$ of the women report being in the labor force. When one outcome is much more likely than the other, the scale factor tends to be much smaller.) We can multiply the coefficients on the roughly continuous variables-and maybe even the variables measuring the number of children-to get an estimate of the effect of a one-unit increase on the response probability, starting from the mean values. The scale factor for computing the PEAs for logit is about .243 , which is close to .25 , the value of the logistic density at zero. If we use these scale factors to compute the PEA for, say, nwifeinc, we get about -.0051 for logit and about -.0047 for probit. Both are somewhat larger in magnitude than the LPM effect, which is -.0034 .

The scale factors for the APEs are smaller. For probit, the average of $g\left(\hat{\beta}_{0}+\hat{\beta}_{1} x_{i 1}+\cdots+\hat{\beta}_{K} x_{i K}\right)$ across $i$ is about .301 , and for logit it is .179 . When we multiply these scale factors by the nwifeinc coefficients we get -.0038 for logit and -. 0036 for probit. A bootstrap standard error for the probit estimate, using 500 bootstrap samples, is about .0016 . As expected, the APEs are much closer to the LPM effect, and the bootstrap standard error of the APE is very close to the LPM standard error for the nwifeinc coeffcient, .0015 . The same is true of the APEs for the other explanatory variables. Even for a discrete variable such as kidslt6-of which over $96 \%$ of the sample takes on zero or one-the APE for the logit model is about -.258 , and for the probit it is -.261 (bootstrap standard error $=.033$ ). The partial effect in the LPM is -.262 (standard error $=.032$ ).

Potentially the biggest difference between the LPM model, on the one hand, and the logit and probit models on the other is that the LPM implies constant marginal effects for educ, kidlt6, and so on, while the logit and probit models allow for a diminishing effect-for continuous and discrete variables. For example, in the LPM, one more young child, whether going from zero to one or from one to two, is estimated to reduce the labor force participation of a women by .262 , independent of the other income, education, age, and experience of the woman. We just saw that this is a good estimate of the average effect in the sense that it is similar to the estimates for logit and probit (when kidstlt6 is treated as a continuous variable). But the effect might differ across the population, and certainly the effect of having the first young child and the second young child might be different. To get an idea of how much a partial effect might differ from the APE, take a women with nwifeinc $=20.13$,\\
educ $=12.3$, exper $=10.6$, age $=42.5$-which are roughly the sample averagesand kidsge $6=1$. In the probit model, what is the estimated fall in the probability of being in the labor force when going from zero to one small child? We evaluate the standard normal cdf, $\Phi\left(\hat{\beta}_{0}+\hat{\beta}_{1} x_{1}+\cdots+\hat{\beta}_{K} x_{K}\right)$ at kidslt6 $=1$ and kidslt6 $=0$, with the other explanatory variables set at the values just given. We get, roughly, $.373-.707=-.334$, which mean a .334 drop in the probability of being in the labor force. (The scaled coefficient for the PEA is about -.347 , and so it is not much different.) This estimated effect is substantially larger than the constant effect obtained from the LPM. If the woman goes from one young child to two, the probability falls even more, but the marginal effect is not as large: $.117-.373=-.256$.

If we compute the difference in predicted probabilities for each woman at one and zero young children and then average these-that is, if we compute the APE in going from zero to one young child-the estimate is about -.272 , while the APE in going from one to two is about -.220 . The LPM estimate is close to the first estimate, which makes sense when interpreting the LPM as estimates of the average partial effect: less than 4 percent of women have more than one young child, and so the estimated effect of moving from one to two (and certainly from two to three) contributes very little to the APE.

Binary response models apply with little modification to independently pooled cross sections or to other data sets where the observations are independent but not necessarily identically distributed. Often year or other time-period dummy variables are included to account for aggregate time effects. Just as with linear models, probit can be used to evaluate the impact of certain policies in the context of a natural experiment; see Problem 15.13. An application is given in Gruber and Poterba (1994).

\subsection*{15.7 Specification Issues in Binary Response Models}
We now turn to several issues that can arise in applying binary response models to economic data. All of these topics are relevant for general index models, but features of the normal distribution allow us to obtain concrete results in the context of probit models. Therefore, our primary focus is on probit models.

\subsection*{15.7.1 Neglected Heterogeneity}
We begin by studying the consequences of omitting variables when those omitted variables are independent of the included explanatory variables. This is also called the neglected heterogeneity problem. The (structural) model of interest is


\begin{equation*}
\mathbf{P}(y=1 \mid \mathbf{x}, c)=\Phi(\mathbf{x} \boldsymbol{\beta}+\gamma c), \tag{15.34}
\end{equation*}


where $\mathbf{x}$ is $1 \times K$ with $x_{1} \equiv 1$ and $c$ is a scalar. We are interested in the partial effects of the $x_{j}$ on the probability of success, holding $c$ (and the other elements of $\mathbf{x}$ ) fixed. We can write equation (15.34) in latent variable form as $y^{*}=\mathbf{x} \boldsymbol{\beta}+\gamma c+e$, where $y=1\left[y^{*}>0\right]$ and $e \mid \mathbf{x}, c \sim \operatorname{Normal}(0,1)$. Because $x_{1}=1, \mathrm{E}(c)=0$ without loss of generality.

Now suppose that $c$ is independent of $\mathbf{x}$ and $c \sim \operatorname{Normal}\left(0, \tau^{2}\right)$. (Remember, this assumption is much stronger than $\operatorname{Cov}(\mathbf{x}, c)=\mathbf{0}$ or even $\mathrm{E}(c \mid \mathbf{x})=0$ : under independence, the distribution of $c$ given $\mathbf{x}$ does not depend on $\mathbf{x}$.) Given these assumptions, the composite term, $\gamma c+e$, is independent of $\mathbf{x}$ and has a $\operatorname{Normal}\left(0, \gamma^{2} \tau^{2}+1\right)$ distribution. Therefore,\\
$\mathrm{P}(y=1 \mid \mathbf{x})=\mathrm{P}(y c+e>-\mathbf{x} \boldsymbol{\beta} \mid \mathbf{x})=\Phi(\mathbf{x} \boldsymbol{\beta} / \sigma)$,\\
where $\sigma^{2} \equiv \gamma^{2} \tau^{2}+1$. It follows immediately from equation (15.35) that probit of $y$ on $\mathbf{x}$ consistently estimates $\boldsymbol{\beta} / \sigma$. In other words, if $\hat{\boldsymbol{\beta}}$ is the estimator from a probit of $y$ on $\mathbf{x}$, then plim $\hat{\beta}_{j}=\beta_{j} / \sigma$. Because $\sigma=\left(\gamma^{2} \tau^{2}+1\right)^{1 / 2}>1$ (unless $\gamma=0$ or $\tau^{2}=0$ ), $\left|\beta_{j} / \sigma\right|<\left|\beta_{j}\right|$.

The attenuation bias in estimating $\beta_{j}$ in the presence of neglected heterogeneity has prompted statements of the following kind: "In probit analysis, neglected heterogeneity is a much more serious problem than in linear models because, even if the omitted heterogeneity is independent of $\mathbf{x}$, the probit coefficients are inconsistent." We just derived that probit of $y$ on $\mathbf{x}$ consistently estimates $\boldsymbol{\beta} / \sigma$ rather than $\boldsymbol{\beta}$, so the statement is technically correct. However, we should remember that, in nonlinear models, we usually want to estimate partial effects and not just parameters. For the purposes of obtaining the directions of the effects or the relative effects of the continuous explanatory variables, estimating $\boldsymbol{\beta} / \sigma$ is just as good as estimating $\boldsymbol{\beta}$.

To be more precise, the scaled coefficient, $\beta_{j} / \sigma$, has the same sign as $\beta_{j}$, and so we will correctly (with enough data) determine the direction of the partial effect of any variable-discrete, continuous, or some mixture-by estimating the scaled coefficients. Further, the ratio of any two scaled coefficients, $\beta_{j} / \sigma$ and $\beta_{h} / \sigma$, is simply $\beta_{j} / \beta_{h}$. For a continuous $x_{j}$, the partial effect in model (15.34)-which we might call the "structural partial effect"-is\\
$\partial \mathrm{P}(y=1 \mid \mathbf{x}, c) / \partial x_{j}=\beta_{j} \phi(\mathbf{x} \boldsymbol{\beta}+\gamma c)$.

Therefore, the ratio of partial effects for two continuous variables $x_{j}$ and $x_{h}$ is simply $\beta_{j} / \beta_{h}$-the same as the ratio of scaled coefficients.

Is there any quantity of interest we cannot estimate by not being able to estimate $\boldsymbol{\beta}$ ? Yes, although its importance is debatable. Because $c$ is normalized so that $\mathrm{E}(c)=0$, we might be interested in the partial effect in (15.36) evaluated at $c=0$,\\
which is simply $\beta_{j} \phi(\mathbf{x} \boldsymbol{\beta})$. It is clear that we would need to consistently estimate $\boldsymbol{\beta}$ in order to estimate the partial effect at the mean value of heterogeneity (and any value of the covariates). What we consistently estimate from the probit of $y$ on $\mathbf{x}$ is\\
$\left(\beta_{j} / \sigma\right) \phi(\mathbf{x} \boldsymbol{\beta} / \sigma)$.

This expression shows that, if we are interested in the partial effects evaluated at $c=0$, then probit of $y$ on $\mathbf{x}$ does not do the trick. An interesting fact about expression (15.37) is that, even though $\beta_{j} / \sigma$ is closer to zero than $\beta_{j}, \phi(\mathbf{x} \boldsymbol{\beta} / \sigma)$ is larger than $\phi(\mathbf{x} \boldsymbol{\beta})$ because $\phi(z)$ increases as $|z| \rightarrow 0$, and $\sigma>1$. Therefore, for estimating the partial effects in equation (15.36) at $c=0$, it is not clear for what values of $\mathbf{x}$ an attenuation bias exists. Plugging other values of $c$ into equation (15.37) makes little sense, as we cannot identify $\gamma$, and even if we could, we generally know nothing about the distribution of $c$ other than its mean is zero. Indeed, $c$ is typically meant to capture such nebulous concepts as "ability," "health," or "taste for saving," and without further information we have no hope of estimating the distribution of these unobservable attributes.

With $c$ having a normal distribution in the population, the partial effect evaluated at $c=0$ describes only a small fraction of the population. (Technically, $\mathrm{P}(c=0)=0$.) Instead, we can estimate the average partial effect (APE), where we now average out the unobserved heterogeneity and are left with a function of $\mathbf{x}$. In particular, the APE is obtained, for given $\mathbf{x}$, by averaging equation (15.36) across the distribution of $c$ in the population. For emphasis, let $\mathbf{x}^{\circ}$ be a given value of the explanatory variables (which could be, but need not be, the mean value). When we plug $\mathbf{x}^{\circ}$ into equation (15.36) and take the expected value with respect to the distribution of $c$, we get


\begin{equation*}
\mathrm{E}\left[\beta_{j} \phi\left(\mathbf{x}^{\mathrm{o}} \boldsymbol{\beta}+\gamma c\right)\right]=\left(\beta_{j} / \sigma\right) \phi\left(\mathbf{x}^{\mathrm{o}} \boldsymbol{\beta} / \sigma\right) \tag{15.38}
\end{equation*}


In other words, probit of $y$ on $\mathbf{x}$ consistently estimates the average partial effects, which is usually what we want.

The result in equation (15.38) follows from the general treatment of average partial effects in Section 2.2.5. In the current setup, there are no extra conditioning variables, $\mathbf{w}$, and the unobserved heterogeneity is independent of $\mathbf{x}$. It follows from equation (2.35) that the APE with respect to $x_{j}$, evaluated at $\mathbf{x}^{\mathrm{o}}$, is simply $\partial \mathrm{E}\left(y \mid \mathbf{x}^{\mathrm{o}}\right) / \partial x_{j}$. But from the law of iterated expectations, $\mathrm{E}(y \mid \mathbf{x})=\mathrm{E}_{c}[\Phi(\mathbf{x} \boldsymbol{\beta}+\gamma c)]=\Phi(\mathbf{x} \boldsymbol{\beta} / \sigma)$, where $\mathrm{E}_{c}(\cdot)$ denotes the expectation with respect to the distribution of $c$. The derivative of $\Phi(\mathbf{x} \boldsymbol{\beta} / \sigma)$ with respect to $x_{j}$ is $\left(\beta_{j} / \sigma\right) \phi(\mathbf{x} \boldsymbol{\beta} / \sigma)$, which is what we wanted to show.

The bottom line is that omitted heterogeneity in probit models is not a problem when it is independent of $\mathbf{x}$ : ignoring it preserves the signs of all partial effects, gives the same relative effects for continuous explanatory variables, and provides consis-\\
tent estimates of the average partial effects. Of course, the previous arguments hinge on the normality of $c$ and the probit structural equation. If the structural model (15.34) were, say, logit and if $c$ were normally distributed, we would not get a probit or logit for the distribution of $y$ given $\mathbf{x}$; the response probability is more complicated. The lesson from Section 2.2.5 is that we might as well work directly with models for $\mathrm{P}(y=1 \mid \mathbf{x})$ because partial effects of $\mathrm{P}(y=1 \mid \mathbf{x})$ are always the average of the partial effects of $\mathrm{P}(y=1 \mid \mathbf{x}, c)$ over the distribution of $c$.

If $c$ is correlated with $\mathbf{x}$ or is otherwise dependent on $\mathbf{x}$ (for example, if $\operatorname{Var}(c \mid \mathbf{x})$ depends on $\mathbf{x}$ ), then omission of $c$ is serious. In this case we cannot get consistent estimates of the average partial effects. For example, if $c \mid \mathbf{x} \sim \operatorname{Normal}\left(\mathbf{x} \boldsymbol{\delta}, \eta^{2}\right)$, then probit of $y$ on $\mathbf{x}$ gives consistent estimates of $(\boldsymbol{\beta}+\gamma \boldsymbol{\delta}) / \rho$, where $\rho^{2}=\gamma^{2} \eta^{2}+1$. Unless $\gamma=0$ or $\boldsymbol{\delta}=\mathbf{0}$, we do not consistently estimate $\boldsymbol{\beta} / \sigma$. This result is not surprising given what we know from the linear case with omitted variables correlated with the $x_{j}$. We now study what can be done to account for endogenous variables in probit models.

\subsection*{15.7.2 Continuous Endogenous Explanatory Variables}
We now explicitly allow for the case where one of the explanatory variables is correlated with the error term in the latent variable model. One possibility is to estimate an LPM by 2SLS. This procedure is relatively easy and might provide a good estimate of the average effect.

If we want to estimate a probit model with endogenous explanatory variables, we must make some fairly strong assumptions. In this section we consider the case of a continuous endogenous explanatory variable.

Write the model as


\begin{align*}
& y_{1}^{*}=\mathbf{z}_{1} \boldsymbol{\delta}_{1}+\alpha_{1} y_{2}+u_{1}  \tag{15.39}\\
& y_{2}=\mathbf{z}_{1} \boldsymbol{\delta}_{21}+\mathbf{z}_{2} \boldsymbol{\delta}_{22}+v_{2}=\mathbf{z} \boldsymbol{\delta}_{2}+v_{2}  \tag{15.40}\\
& y_{1}=1\left[y_{1}^{*}>0\right], \tag{15.41}
\end{align*}


where ( $u_{1}, v_{2}$ ) has a zero mean, bivariate normal distribution, and is independent of $\mathbf{z}$. Equation (15.39), along with equation (15.41), is the structural equation; equation (15.40) is a reduced form for $y_{2}$, which is endogenous if $u_{1}$ and $v_{2}$ are correlated. If $u_{1}$ and $v_{2}$ are independent, there is no endogeneity problem. Because $v_{2}$ is normally distributed, we are assuming that $y_{2}$ given $\mathbf{z}$ is normal; thus $y_{2}$ should have features of a normal random variable. (For example, $y_{2}$ should not be a discrete variable.)

The model is applicable when $y_{2}$ is correlated with $u_{1}$ because of omitted variables or measurement error. It can also be applied to the case where $y_{2}$ is determined jointly with $y_{1}$, but with a caveat. If $y_{1}$ appears on the right-hand side in a linear structural equation for $y_{2}$, then the reduced form for $y_{2}$ cannot be found with $v_{2}$ having the stated properties. However, if $y_{1}^{*}$ appears in a linear structural equation for $y_{2}$, then $y_{2}$ has the reduced form given by equation (15.40); see Maddala (1983, Chap. 7) for further discussion.

The normalization that gives the parameters in equation (15.39) an average partial effect interpretation, at least in the omitted variable and simultaneity contexts, is $\operatorname{Var}\left(u_{1}\right)=1$, just as in a probit model with all explanatory variables exogenous. To see this point, consider the outcome on $y_{1}$ at two different outcomes of $y_{2}$, say $y_{2}^{\circ}$ and $y_{2}^{\mathrm{o}}+1$. Holding the observed exogenous factors fixed at $\mathbf{z}_{1}^{\mathrm{o}}$, and holding $u_{1}$ fixed, the difference in responses is\\
$1\left[\mathbf{z}_{1}^{\mathrm{o}} \boldsymbol{\delta}_{1}+\alpha_{1}\left(y_{2}^{\mathrm{o}}+1\right)+u_{1} \geq 0\right]-1\left[\mathbf{z}_{1}^{\mathrm{o}} \boldsymbol{\delta}_{1}+\alpha_{1} y_{2}^{\mathrm{o}}+u_{1} \geq 0\right]$.\\
(This difference can take on the values $-1,0$, and 1 .) Because $u_{1}$ is unobserved, we cannot estimate the difference in responses for a given population unit. Nevertheless, if we average across the distribution of $u_{1}$, which is $\operatorname{Normal}(0,1)$, we obtain\\
$\Phi\left[\mathbf{z}_{1}^{\mathrm{o}} \boldsymbol{\delta}_{1}+\alpha_{1}\left(y_{2}^{\mathrm{o}}+1\right)\right]-\Phi\left(\mathbf{z}_{1}^{\mathrm{o}} \boldsymbol{\delta}_{1}+\alpha_{1} y_{2}^{\mathrm{o}}\right)$.\\
Therefore, $\boldsymbol{\delta}_{1}$ and $\alpha_{1}$ are the parameters appearing in the APE. (Alternatively, if we begin by allowing $\sigma_{1}^{2}=\operatorname{Var}\left(u_{1}\right)>0$ to be unrestricted, the APE would depend on $\boldsymbol{\delta}_{1} / \sigma_{1}$ and $\alpha_{1} / \sigma_{1}$, and so we should just rescale $u_{1}$ to have unit variance. The variance and slope parameters are not separately identified, anyway.) The proper normalization for $\operatorname{Var}\left(u_{1}\right)$ should be kept in mind, as two-step procedures, which we cover in the following paragraphs, only consistently estimate $\boldsymbol{\delta}_{1}$ and $\alpha_{1}$ up to scale; we have to do a little more work to obtain estimates of the APE.

The most useful two-step approach is a control function approach due to Rivers and Vuong (1988), as it leads to a simple test for endogeneity of $y_{2}$. To derive the procedure, first note that, under joint normality of $\left(u_{1}, v_{2}\right)$, with $\operatorname{Var}\left(u_{1}\right)=1$, we can write\\
$u_{1}=\theta_{1} v_{2}+e_{1}$,\\
where $\theta_{1}=\eta_{1} / \tau_{2}^{2}, \eta_{1}=\operatorname{Cov}\left(v_{2}, u_{1}\right), \tau_{2}^{2}=\operatorname{Var}\left(v_{2}\right)$, and $e_{1}$ is independent of $\mathbf{z}$ and $v_{2}$ (and therefore of $y_{2}$ ). Because of joint normality of ( $u_{1}, v_{2}$ ), $e_{1}$ is also normally distributed with $\mathrm{E}\left(e_{1}\right)=0$ and $\operatorname{Var}\left(e_{1}\right)=\operatorname{Var}\left(u_{1}\right)-\eta_{1}^{2} / \tau_{2}^{2}=1-\rho_{1}^{2}$, where $\rho_{1}= \operatorname{Corr}\left(v_{2}, u_{1}\right)$. We can now write\\
$y_{1}^{*}=\mathbf{z}_{1} \boldsymbol{\delta}_{1}+\alpha_{1} y_{2}+\theta_{1} v_{2}+e_{1}$,\\
$e_{1} \mid \mathbf{z}, y_{2}, v_{2} \sim \operatorname{Normal}\left(0,1-\rho_{1}^{2}\right)$.

A standard calculation shows that\\
$\mathrm{P}\left(y_{1}=1 \mid \mathbf{z}, y_{2}, v_{2}\right)=\Phi\left[\left(\mathbf{z}_{1} \boldsymbol{\delta}_{1}+\alpha_{1} y_{2}+\theta_{1} v_{2}\right) /\left(1-\rho_{1}^{2}\right)^{1 / 2}\right]$.\\
Assuming for the moment that we observe $v_{2}$, then probit of $y_{1}$ on $\mathbf{z}_{1}, y_{2}$, and $v_{2}$ consistently estimates $\boldsymbol{\delta}_{\rho 1} \equiv \boldsymbol{\delta}_{1} /\left(1-\rho_{1}^{2}\right)^{1 / 2}, \alpha_{\rho 1} \equiv \alpha_{1} /\left(1-\rho_{1}^{2}\right)^{1 / 2}$, and $\theta_{\rho 1} \equiv \theta_{1} /\left(1-\rho_{1}^{2}\right)^{1 / 2}$. Notice that because $\rho_{1}^{2}<1$, each scaled coefficient is greater than its unscaled counterpart unless $y_{2}$ is exogenous ( $\rho_{1}=0$ ).

Since we do not know $\boldsymbol{\delta}_{2}$, we must first estimate it, as in the following procedure:\\
Procedure 15.1: (a) Run the OLS regression $y_{2}$ on $\mathbf{z}$ and save the residuals $\hat{v}_{2}$.\\
(b) Run the probit $y_{1}$ on $\mathbf{z}_{1}, y_{2}, \hat{v}_{2}$ to get consistent estimators of the scaled coefficients $\boldsymbol{\delta}_{\rho 1}, \alpha_{\rho 1}$, and $\theta_{\rho 1}$.

A nice feature of Procedure 15.1 is that the usual probit $t$ statistic on $\hat{v}_{2}$ is a valid test of the null hypothesis that $y_{2}$ is exogenous, that is, $\mathrm{H}_{0}: \theta_{1}=0$. If $\theta_{1} \neq 0$, the usual probit standard errors and test statistics are not strictly valid, and we have only estimated $\boldsymbol{\delta}_{1}$ and $\alpha_{1}$ up to scale. The asymptotic variance of the two-step estimator can be derived using the M-estimator results in Section 12.5.2; see also Rivers and Vuong (1988). Problem 15.15 asks you to obtain the correct variance matrix and to show that it is always greater than the incorrect one, which ignores estimation of $\boldsymbol{\delta}_{2}$, when $\theta_{1} \neq 0$. The bootstrap can be used, but some bootstrap samples may have little variation in $y_{1}$ if one outcome is much more likely.

Under $\mathrm{H}_{0}: \theta_{1}=0, e_{1}=u_{1}$, and so the distribution of $v_{2}$ plays no role under the null. Therefore, the test of exogeneity is valid without assuming normality or homoskedasticity of $v_{2}$, and it can be applied very broadly, even if $y_{2}$ is a binary variable. Unfortunately, if $y_{2}$ and $u_{1}$ are correlated, normality of $v_{2}$ is crucial.

Example 15.3 (Testing Exogeneity of Education in the Women's LFP Model): We test the null hypothesis that educ is exogenous in the married women's labor force participation equation. We first obtain the reduced form residuals, $\hat{v}_{2}$, from regressing educ on all exogenous variables, including motheduc, fatheduc, and huseduc. Then, we add $\hat{v}_{2}$ to the probit from Example 15.2. The $t$ statistic on $\hat{v}_{2}$ is only .867 , which is weak evidence against the null hypothesis that educ is exogenous. As always, this conclusion hinges on the assumption that the instruments for educ are themselves exogenous.

After the two-step estimation, we can easily obtain estimates of the unscaled parameters, $\boldsymbol{\beta}_{1}=\left(\boldsymbol{\delta}_{1}^{\prime}, \alpha_{1}\right)^{\prime}$, which then allows us to estimate partial effects. From the\\
two-step estimation procedure, we have consistent estimators of $\boldsymbol{\delta}_{2}$ and $\tau_{2}^{2}$ from the first-stage regression, and then $\boldsymbol{\delta}_{\rho 1}, \alpha_{\rho 1}$, and $\theta_{\rho 1}$ from the second stage. Straightforward algebra shows that $1+\theta_{\rho 1}^{2} \tau_{2}^{2}=1 /\left(1-\rho_{1}^{2}\right)$. Now, because $\boldsymbol{\delta}_{1}=\left(1-\rho_{1}^{2}\right)^{1 / 2} \boldsymbol{\delta}_{\rho 1}$ and $\alpha_{1}=\left(1-\rho_{1}^{2}\right)^{1 / 2} \alpha_{\rho 1}$, it follows that $\boldsymbol{\beta}_{1}=\boldsymbol{\beta}_{\rho 1} /\left(1+\theta_{\rho 1}^{2} \tau_{2}^{2}\right)^{1 / 2}$, where $\boldsymbol{\beta}_{\rho 1}= \left(\boldsymbol{\delta}_{\rho 1}^{\prime}, \alpha_{\rho 1}\right)^{\prime}$ is the vector of scaled coefficients. Therefore, we can obtain consistent estimators of the original coefficients as\\
$\hat{\boldsymbol{\beta}}_{1}=\hat{\boldsymbol{\beta}}_{\rho 1} /\left(1+\hat{\theta}_{\rho 1}^{2} \hat{\tau}_{2}^{2}\right)^{1 / 2}$,\\
where all quantities on the right-hand side of equation (15.45) are available from the two-step estimation procedure.

Given $\hat{\boldsymbol{\delta}}_{1}$ and $\hat{\alpha}_{1}$, we can compute derivatives and differences in $\Phi\left(\mathbf{z}_{1} \hat{\boldsymbol{\delta}}_{1}+\hat{\alpha}_{1} y_{2}\right)$ at interesting values of $\mathbf{z}_{1}$ and $y_{2}$. Sometimes it is useful to evaluate the partial effects at the mean values. At other times the average partial effects are preferable. For continuous explanatory variables (including $y_{2}$ ), the scale factor for multiplying, say, $\hat{\alpha}_{1}$ is $N^{-1} \sum_{i=1}^{N} \phi\left(\mathbf{z}_{i 1} \hat{\boldsymbol{\delta}}_{1}+\hat{\alpha}_{1} y_{i 2}\right)$. We can use the delta method to obtain standard errors for the APEs (or partial effects at the average), but the required calculations are tedious, given the two-step nature of the estimation. The bootstrap is easily applied (but computationally expensive): for each bootstrap sample, one computes the two steps of the procedure and obtains the scaled and unscaled coefficients, and then computes the APEs (or other partial effects of interest). As described in Section 12.8.2, one computes the standard deviation across the bootstrap replications.

An alternative method of computing the APEs does not exploit the normality assumption for $v_{2}$. We can use the results in Section 2.2.5 by writing $y_{1}=1\left[\mathbf{z}_{1} \boldsymbol{\delta}_{1}+\right. \left.\alpha_{1} y_{2}+u_{1}>0\right]$ and setting $q \equiv u_{1}, \mathbf{x} \equiv\left(\mathbf{z}_{1}, y_{2}\right)$, and $w=v_{2}$ (a scalar in this case). Because $y_{1}$ is a deterministic function of $\left(\mathbf{z}_{1}, y_{2}, u_{1}\right), v_{2}$ is trivially redundant in $E\left(y_{1} \mid \mathbf{z}_{1}, y_{2}, u_{1}, v_{2}\right)$. Further, we have already used that $\mathrm{D}\left(u_{1} \mid \mathbf{z}_{1}, y_{2}, v_{2}\right)=\mathrm{D}\left(u_{1} \mid v_{2}\right)$, and so assumption (2.33) holds as well. It follows that APEs are obtained by taking derivatives or differences of


\begin{equation*}
\mathrm{E}_{v_{2}}\left[\Phi\left(\mathbf{z}_{1} \boldsymbol{\delta}_{\rho 1}+\alpha_{\rho 1} y_{2}+\theta_{\rho 1} v_{2}\right)\right] \tag{15.46}
\end{equation*}


with respect to elements of $\left(\mathbf{z}_{1}, y_{2}\right)$. Using Lemma 12.1, a consistent estimator of (15.46) is given by


\begin{equation*}
N^{-1} \sum_{i=1}^{N} \Phi\left(\mathbf{z}_{1} \hat{\boldsymbol{\delta}}_{\rho 1}+\hat{\alpha}_{\rho 1} y_{2}+\hat{\theta}_{\rho 1} \hat{v}_{i 2}\right) \tag{15.47}
\end{equation*}


where the $\hat{v}_{i 2}$ are the first-stage OLS residuals from regressing $y_{i 2}$ on $\mathbf{z}_{i}, i=1, \ldots, N$. This approach provides a different strategy for estimating APEs: simply compute\\
partial effects with respect to $\mathbf{z}_{1}$ and $y_{2}$ after the second-stage estimation, but then average these across the $\hat{v}_{i 2}$ in the sample. For a continuous explanatory variable, $y_{2}$ often being the most important, the estimated APE is $\hat{\alpha}_{1}\left[N^{-1} \sum_{i=1}^{N} \phi\left(\mathbf{z}_{i 1} \hat{\boldsymbol{\delta}}_{\rho 1}+\hat{\alpha}_{\rho 1} y_{i 2}+\right.\right. \left.\hat{\theta}_{\rho 1} \hat{v}_{i 2}\right)$ ], and a standard error for this APE can be obtained via the delta method or the bootstrap. Again, the bootstrap is much easier to apply: for each bootstrap sample, one applies the two-step estimation method and computes the APEs for that sample. The process is repeated to obtain a bootstrap standard error for the APEs.

The control function (CF) approach has some decided advantages over another two-step approach-one that appears to mimic the 2SLS estimation of the linear model. Rather than conditioning on $v_{2}$ along with $\mathbf{z}$ (and therefore $y_{2}$ ) to obtain $\mathrm{P}\left(y_{1}=1 \mid \mathbf{z}, v_{2}\right)=\mathrm{P}\left(y_{1}=1 \mid \mathbf{z}, y_{2}, v_{2}\right)$, we can obtain $\mathrm{P}\left(y_{1}=1 \mid \mathbf{z}\right)$. To find the latter probability, we plug in the reduced form for $y_{2}$ to get $y_{1}=1\left[\mathbf{z}_{1} \boldsymbol{\delta}_{1}+\alpha_{1}\left(\mathbf{z} \boldsymbol{\delta}_{2}\right)+\alpha_{1} v_{2}+\right. \left.u_{1}>0\right]$. Because $\alpha_{1} v_{2}+u_{1}$ is independent of $\mathbf{z}$ and ( $u_{1}, v_{2}$ ) has a bivariate normal distribution, $\mathbf{P}\left(y_{1}=1 \mid \mathbf{z}\right)=\Phi\left\{\left[\mathbf{z}_{1} \boldsymbol{\delta}_{1}+\alpha_{1}\left(\mathbf{z} \boldsymbol{\delta}_{2}\right)\right] / \omega_{1}\right\}$, where $\omega_{1}^{2} \equiv \operatorname{Var}\left(\alpha_{1} v_{2}+u_{1}\right)= \alpha_{1}^{2} \tau_{2}^{2}+1+2 \alpha_{1} \operatorname{Cov}\left(v_{2}, u_{1}\right)$. (A two-step procedure now proceeds by using the same first-step OLS regression-in this case, to get the fitted values, $\hat{y}_{i 2}=\mathbf{z}_{i} \hat{\boldsymbol{\delta}}_{2}$-now followed by a probit of $y_{i 1}$ on $\mathbf{z}_{i 1}, \hat{y}_{i 2}$. It is easily seen that this method estimates the coefficients up to the common scale factor $1 / \omega_{1}$, which can be any positive value (unlike in the CF case, where we know the scale factor is greater than unity).

As with the CF approach, getting the appropriate standard errors is difficult. A primary drawback of the method that inserts $\hat{y}_{2}$ for $y_{2}$ is that it does not provide a simple test of the null hypothesis that $y_{2}$ is exogenous. Plus, the coefficients cannot be directly compared to the usual probit estimates in a Hausman test because of the different scale factors. Additionally, while the APEs can be recovered in much the same way as they can be for the CF approach, equation (15.47) is not available for the fitted values approach. Finally, plugging in fitted values is strictly limited to the structural equation (15.39); adding other functions of $y_{2}$ is very cumbersome, and one is prone to making mistakes. See Problem 15.14 and the discussion at the end of this subsection.

Example 15.3 (Endogeneity of Nonwife Income in the Women's LFP Model): We use the data in MROZ.RAW to test the null hypothesis that nwifeinc is exogenous in the probit model estimated in Table 15.1. We use as an instrument for nwifeinc husband's years of schooling, huseduc. Therefore, the identification assumption (perhaps a tenuous one) is that husband's schooling is unrelated to factors that affect a married woman's labor force decision once nwifeinc and the other variables (including the woman's education) are accounted for. In the first-stage regression of nwifeinc on huseduc and the other explanatory variables listed in Table 15.1, the fully robust $t$\\
statistic on huseduc is about 6.92, which is hardly surprising, because nwifeinc is pretty highly correlated with husband's labor earnings, which in turn depends on husband's education. When the reduced form residuals, $\hat{v}_{2}$, are added to the probit, its coefficient is about .027 with a $t$ statistic of about 1.41 -only moderate evidence that nwifeinc is endogenous. The coefficient on nwifeinc becomes about -.037 , which is certainly larger in magnitude than the probit estimate in Table 15.1. But we must compare partial effects. If we average the scale factor $\phi\left(\mathbf{z}_{i 1} \hat{\boldsymbol{\delta}}_{\rho 1}+\hat{\alpha}_{\rho 1} y_{i 2}+\hat{\theta}_{\rho 1} \hat{v}_{i 2}\right)$ across all $i$ (so that we are averaging across the distribution of $\left(\mathbf{z}_{1}, y_{2}\right)$ in addition to averaging out $v_{2}$ ), we obtain about .300 . Therefore, the APE for nwifeinc is $.300(-.037) \approx-.011$ with a boostrap standard error, based on 500 bootstrap samples, of about .0058 . The estimate, which is marginally statistically significant based on the bootstrap standard error, is about three times larger than the probit estimate that treats nwifeinc as exogenous, which was -.0036 . If we use equation (15.46) to recover the estimates of the original coefficients, $\hat{\alpha}_{1}=-.0355$, and this agrees with the joint MLE (as it should because the model is just identified). The scale factor based on $\phi\left(\mathbf{z}_{i 1} \hat{\boldsymbol{\delta}}_{1}+\hat{\alpha}_{1} y_{i 2}\right)$ is about .297, which is very close to the scale factor based on (15.47). For comparision, we estimate a linear probability model by 2SLS. The 2SLS coefficient on nwifeinc is about -.012 (robust standard error $=.0059$ ), which, for practical purposes, is the same as the APE for the CF estimate in the probit model.

In the end, we are left in a difficult situation: the evidence against exogeneity of nwifeinc is not particularly strong, yet the APE that treats nwifeinc as endogenous is quite a bit larger than the APE that treats nwifeinc as exogenous.

In the previous example the model is just identified because we have one instrument, huseduc, for the endogenous variable nwifeinc. If we have overidentifying restrictions, these are easily tested using the CF approach. The key restriction is that $\mathrm{D}\left(u_{1} \mid y_{2}, \mathbf{z}\right)=\mathrm{D}\left(u_{1} \mid v_{2}\right)$; that is, the conditional distribution depends only on the linear combination $y_{2}-\mathbf{z} \boldsymbol{\delta}_{2}$. If $\mathbf{z}_{2}$ is the $1 \times L_{2}$ vector of exogenous variables excluded from the structural equation, then this restriction on $\mathrm{D}\left(u_{1} \mid y_{2}, \mathbf{z}\right)$ imposes $L_{2}-1$ overidentifying restrictions. We can test these by including any $1 \times\left(L_{2}-1\right)$ subvector of $\mathbf{z}_{2}$, say $\mathbf{h}_{2}$, as additional explanatory variables in part (b) of Procedure 15.1, and conduct a joint significance test. In effect, we are testing $\mathrm{E}\left(\mathbf{h}_{2}^{\prime} e_{1}\right)=\mathbf{0}$, where $u_{1}=\rho_{1} v_{2}+e_{1}$. (Naturally, if we are allowing $y_{2}$ to be endogenous under the null-as we should - then the first-step estimation of $\boldsymbol{\delta}_{2}$ should be accounted for, either via the delta method or the bootstrap.) It can be shown that the test is invariate to which subset of $\mathbf{z}_{2}$ we choose for $\mathbf{h}_{2}$, provided $\mathbf{h}_{2}$ has $L_{2}-1$ elements.

Rather than use a two-step procedure, we can estimate equations (15.39)-(15.41) by conditional maximum likelihood estimation (CMLE). To obtain the joint distribution of $\left(y_{1}, y_{2}\right)$, conditional on $\mathbf{z}$, recall that


\begin{equation*}
f\left(y_{1}, y_{2} \mid \mathbf{z}\right)=f\left(y_{1} \mid y_{2}, \mathbf{z}\right) f\left(y_{2} \mid \mathbf{z}\right) \tag{15.48}
\end{equation*}


(see Property CD. 2 in Appendix 13A). Since $y_{2} \mid \mathbf{z} \sim \operatorname{Normal}\left(\mathbf{z} \boldsymbol{\delta}_{2}, \tau_{2}^{2}\right)$, the density $f\left(y_{2} \mid \mathbf{z}\right)$ is easy to write down. We can also derive the conditional density of $y_{1}$ given $\left(y_{2}, \mathbf{z}\right)$. Since $v_{2}=y_{2}-\mathbf{z} \boldsymbol{\delta}_{2}$ and $y_{1}=1\left[y_{1}^{*}>0\right]$,


\begin{equation*}
\mathrm{P}\left(y_{1}=1 \mid y_{2}, \mathbf{z}\right)=\Phi\left[\frac{\mathbf{z}_{1} \boldsymbol{\delta}_{1}+\alpha_{1} y_{2}+\left(\rho_{1} / \tau_{2}\right)\left(y_{2}-\mathbf{z} \boldsymbol{\delta}_{2}\right)}{\left(1-\rho_{1}^{2}\right)^{1 / 2}}\right], \tag{15.49}
\end{equation*}


where we have used the fact that $\theta_{1}=\rho_{1} / \tau_{2}$.\\
Let $w$ denote the term in inside $\Phi(\cdot)$ in equation (15.49). Then we have derived\\
$f\left(y_{1}, y_{2} \mid \mathbf{z}\right)=\{\Phi(w)\}^{y_{1}}\{1-\Phi(w)\}^{1-y_{1}}\left(1 / \tau_{2}\right) \phi\left[\left(y_{2}-\mathbf{z} \boldsymbol{\delta}_{2}\right) / \tau_{2}\right]$,\\
and so the log likelihood for observation $i$ (apart from terms not depending on the parameters) is\\
$y_{i 1} \log \Phi\left(w_{i}\right)+\left(1-y_{i 1}\right) \log \left[1-\Phi\left(w_{i}\right)\right]-\frac{1}{2} \log \left(\tau_{2}^{2}\right)-\frac{1}{2}\left(y_{i 2}-\mathbf{z}_{i} \boldsymbol{\delta}_{2}\right)^{2} / \tau_{2}^{2}$,\\
where we understand that $w_{i}$ depends on the parameters $\left(\boldsymbol{\delta}_{1}, \alpha_{1}, \rho_{1}, \boldsymbol{\delta}_{2}, \tau_{2}\right)$ :\\
$w_{i} \equiv\left[\mathbf{z}_{i 1} \boldsymbol{\delta}_{1}+\alpha_{1} y_{i 2}+\left(\rho_{1} / \tau_{2}\right)\left(y_{i 2}-\mathbf{z}_{i} \boldsymbol{\delta}_{2}\right)\right] /\left(1-\rho_{1}^{2}\right)^{1 / 2}$.\\
Summing expression (15.50) across all $i$ and maximizing with respect to all parameters gives the MLEs of $\boldsymbol{\delta}_{1}, \alpha_{1}, \rho_{1}, \boldsymbol{\delta}_{2}, \tau_{2}^{2}$. The general theory of conditional MLE applies, and so standard errors can be obtained using the estimated Hessian, the estimated expected Hessian, or the outer product of the score. MLE applied to this model sometimes goes by the name of instrumental variables probit, or IV probit.

Maximum likelihood estimation has some decided advantages over two-step procedures. First, MLE is more efficient than any two-step procedure. Second, we get direct estimates of $\boldsymbol{\delta}_{1}$ and $\alpha_{1}$, the parameters of interest for computing partial effects. Evans, Oates, and Schwab (1992) study peer effects on teenage behavior using the full MLE. Of course, obtaining standard errors for the APEs still requires using the delta method or the bootstrap.

Testing that $y_{2}$ is exogenous is easy once the MLE has been obtained: just test $\mathrm{H}_{0}: \rho_{1}=0$ using an asymptotic $t$ test. We could also use a likelihood ratio test.

The drawback to the MLE is computational. Sometimes it can be difficult to get the iterations to converge, as $\hat{\rho}_{1}$ sometimes tends toward 1 or -1 .

Comparing the Rivers-Vuong approach to the MLE shows that the former is a limited information procedure. Essentially, Rivers and Vuong focus on $f\left(y_{1} \mid y_{2}, \mathbf{z}\right)$, where they replace the unknown $\boldsymbol{\delta}_{2}$ with the OLS estimator $\hat{\boldsymbol{\delta}}_{2}$ (and they ignore the\\
rescaling problem by taking $e_{1}$ in equation (15.43) to have unit variance). MLE estimates the parameters using the information in $f\left(y_{1} \mid y_{2}, \mathbf{z}\right)$ and $f\left(y_{2} \mid \mathbf{z}\right)$ simultaneously. For the initial test of whether $y_{2}$ is exogenous, the Rivers-Vuong approach has significant computational advantages. If exogeneity is rejected, it might be worth doing MLE. As shown by Rivers and Vuong (1988), the two-step estimators and MLEs are identical when the model is just identified.

Another benefit of the MLE approach for this and related problems is that it forces discipline on us in coming up with consistent estimation procedures and correct standard errors. It is easy to abuse two-step procedures if we are not careful in deriving estimating equations. With MLE, although it can be difficult to derive joint distributions of the endogenous variables given the exogenous variables, we know that, if the underlying distributional assumptions hold, consistent and efficient estimators are obtained.

The CF approach has a somewhat subtle robustness property. It is easy to see that the CF approach is consistent if we assume that $\mathrm{D}\left(u_{1} \mid \mathbf{z}, v_{2}\right)=\mathrm{D}\left(u_{1} \mid v_{2}\right)$ and that $\mathrm{D}\left(u_{1} \mid v_{2}\right)$ is normal with mean linear in $v_{2}$ and constant variance. Independence between ( $u_{1}, v_{2}$ ) and $\mathbf{z}$ along with bivariate normality of ( $u_{1}, v_{2}$ ) is sufficient but not necessary. It is certainly possible for $\mathrm{D}\left(u_{1} \mid v_{2}\right)$ to be normal without $v_{2}$ having a normal distribution. In fact, we could even allow the mean $\mathrm{E}\left(u_{1} \mid v_{2}\right)$ to be a known, nonlinear function of $v_{2}$-say, a quadratic-and then add these functions to the probit in the second-step estimation.

For notational and interpretational simplicity, our previous analysis has assumed that a single endogenous explanatory variable appears additively inside the index function. But a moment's thought shows that the previous analysis goes through much more generally. For example, suppose we specify\\
$y_{1}=1\left[\mathbf{g}_{1}\left(\mathbf{z}_{1}, y_{2}\right) \boldsymbol{\beta}_{1}+u_{1}>0\right]$\\
$y_{2}=\mathbf{g}_{2}(\mathbf{z}) \boldsymbol{\delta}_{2}+v_{2}$,\\
where $\mathbf{g}_{1}(\cdot)$ and $\mathbf{g}_{2}(\cdot)$ are known vector functions. As we discussed in Section 6.2, CF approaches offer considerable flexibility in specifying functional form-provided the underlying assumptions hold. Because $y_{2}$ is a function of ( $\mathbf{z}, v_{2}$ ), if we maintain independence between ( $u_{1}, v_{2}$ ) and $\mathbf{z}$, then $\mathrm{D}\left(u_{1} \mid \mathbf{z}, y_{2}, v_{2}\right)=\mathrm{D}\left(u_{1} \mid \mathbf{z}, v_{2}\right)=\mathrm{D}\left(u_{1} \mid v_{2}\right)$, and the previous analysis, whether it is the Rivers-Vuong CF approach or MLEgoes through by replacing $\mathbf{x}_{1} \equiv\left(\mathbf{z}_{1}, y_{2}\right)$ with $\mathbf{x}_{1} \equiv \mathbf{g}_{1}\left(\mathbf{z}_{1}, y_{2}\right)$. For example, we can choose $\mathbf{g}_{1}\left(\mathbf{z}_{1}, y_{2}\right)=\left(\mathbf{z}_{1}, y_{2} \mathbf{z}_{1}\right)$ to allow full interaction effects (in addition to a main effect, since the first element of $\mathbf{z}_{1}$ should be unity). Adding a quadratic or other polynomials in $y_{2}$ causes no difficulties. Partial effect and APE calculations become\\
more complicated only insofar as derivatives of the function $\mathbf{g}_{1}\left(\mathbf{z}_{1}, y_{2}\right)$ must be obtained. So, for example, in a model with $y_{1}=1\left[\mathbf{z}_{1} \boldsymbol{\delta}_{1}+y_{2} \mathbf{z}_{1} \boldsymbol{\alpha}_{1}+y_{1} y_{2}^{2}+u_{1}>0\right]$, the partial effect of $y_{2}$ is $\left(\mathbf{z}_{1} \boldsymbol{\alpha}_{1}+2 \gamma_{1} y_{2}\right) \phi\left(\mathbf{z}_{1} \boldsymbol{\delta}_{1}+y_{2} \mathbf{z}_{1} \boldsymbol{\alpha}_{1}+\gamma_{1} y_{2}^{2}\right)$. All of the parameters can be estimated using the two-step CF estimator or the MLE. Because of our assumptions, testing for exogeneity proceeds exactly as before: we can use a simple $t$ statistic on $\hat{v}_{i 2}$ in the second-step probit (or directly test for zero correlation between $u_{1}$ and $v_{2}$ in the context of MLE). (Standard software packages that have an "instrumental variables probit" command are easily tricked into estimating general models, provided $y_{2}$ appears on its own as an explanatory variable-such as models with quadratics and interactions in $y_{2}$. One simply specifies $y_{2}$ as the lone endogenous explanatory variable with other functions of $y_{2}$ listed as if they were exogenous; the likelihood function will be properly computed.)

Another nice feature of the CF approach is that we can allow for a strictly monotonic transformation of the endogenous explanatory variable in the reduced form. For example, suppose $y_{2}=\log$ (income). Then, of course, income $=\exp \left(y_{2}\right)$, and so any function of income-say, its level, quadratic, or interaction terms-can appear in the structural equation if these are deemed better functional forms than $\log$ (income) in the probit. This flexibility is handy because sometimes a variable needs to be transformed before it is reasonable to assume that it has a reduced form with an additive error that is independent of $\mathbf{z}$. As another example, if $w_{2}$ is an endogenous variable that is strictly in the unit interval, we might choose $y_{2}= \log \left[w_{2} /\left(1-w_{2}\right)\right]$ as the variable that has a reduced form linear in parameters with an additive error. Then $w_{2}$ itself, and any function of it, can appear in the probit model for $y_{1}$ because $w_{2}$ is a well-defined function of $y_{2}$.

That including the reduced form residuals in the CF approach accounts for endogeneity of $y_{2}$, even if we have general functions of $y_{2}$ and $\mathbf{z}_{1}$ in the model, hinges crucially on independence between ( $u_{1}, v_{2}$ ) and $\mathbf{z}$. Because of the additivity of $v_{2}$ in the reduced form, independence between $v_{2}$ and $\mathbf{z}$ pretty much rules out any discreteness in $y_{2}$. And, we are assuming at least normality of $\mathrm{D}\left(u_{1} \mid v_{2}\right)$ (although that can be relaxed, as we briefly discuss in Section 15.7.5).

If we have multiple exogenous explanatory variables, say a $1 \times G_{1}$ row vector $\mathbf{y}_{2}$, we can solve the endogeneity problem in a very similar way. Maximum likelihood estimation becomes much more difficult computationally, but the CF approach is straightforward. The key assumption for the CF approach is $\mathbf{y}_{2}=\mathbf{z} \boldsymbol{\Delta}_{2}+\mathbf{v}_{2}$, where $\mathrm{D}\left(u_{1} \mid \mathbf{z}, \mathbf{v}_{2}\right)=\mathrm{D}\left(u_{1} \mid \mathbf{v}_{2}\right)$ and this latter distribution is normal with constant variance and mean linear in $\mathbf{v}_{2}$. Calculation of partial effects requires averaging out $\mathbf{v}_{2}$, either by computing the unscaled coefficients or using the extension of equation (15.47).

Maximum likelihood estimation is also feasible when ( $u_{1}, \mathbf{v}_{2}$ ) is jointly normal and independent of $\mathbf{z}$.

Finally, one should not minimize the usefulness of estimating a linear model for $y_{1}$ using standard IV estimation, probably 2SLS. As we discussed in Section 15.6 with exogenous explanatory variables, the OLS estimates of an LPM can provide good estimates of the APEs. The same is true when one (or more) explanatory variable is endogenous. At a minimum, it makes sense to compare 2SLS estimates of an LPM with the APEs obtained from the probit model with endogenous $y_{2}$.

\subsection*{15.7.3 Binary Endogenous Explanatory Variable}
We now consider the case where the probit model contains a binary explanatory variable that is endogenous. We begin with the simplest model,


\begin{equation*}
y_{1}=1\left[\mathbf{z}_{1} \boldsymbol{\delta}_{1}+\alpha_{1} y_{2}+u_{1}>0\right] \tag{15.51}
\end{equation*}


$y_{2}=1\left[\mathbf{z} \boldsymbol{\delta}_{2}+v_{2}>0\right]$,\\
where ( $u_{1}, v_{2}$ ) is independent of $\mathbf{z}$ and distributed as bivariate normal with mean zero, each has unit variance, and $\rho_{1}=\operatorname{Corr}\left(u_{1}, v_{2}\right)$. If $\rho_{1} \neq 0$, then $u_{1}$ and $y_{2}$ are correlated, and probit estimation of equation (15.51) is inconsistent for $\boldsymbol{\delta}_{1}$ and $\alpha_{1}$.

Model (15.51) and (15.52) applies primarily to omitted variables situations. In particular, we could not obtain the reduced form (15.52) if the structural model has $y_{1}$ as a determinant of $y_{2}$. Measurement error in binary responses does not lead to this model, either.

As discussed in Section 15.7.2, the normalization $\operatorname{Var}\left(u_{1}\right)=1$ is the proper one for computing partial effects. Often, the effect of $y_{2}$ is of primary interest, especially when $y_{2}$ indicates participation in some sort of program, such as job training, and the binary outcome $y_{1}$ might denote employment status. The average treatment effect (for a given value of $\mathbf{z}_{1}$ ) is $\Phi\left(\mathbf{z}_{1} \boldsymbol{\delta}_{1}+\alpha_{1}\right)-\Phi\left(\mathbf{z}_{1} \boldsymbol{\delta}_{1}\right)$. This effect can be computed for different subgroups or averaged across the distribution of $\mathbf{z}_{1}$.

To derive the likelihood function, we again need the joint distribution of $\left(y_{1}, y_{2}\right)$ given $\mathbf{z}$, which we obtain from equation (15.48). To obtain $\mathrm{P}\left(y_{1}=1 \mid y_{2}, \mathbf{z}\right)$, first note that\\
$\mathrm{P}\left(y_{1}=1 \mid v_{2}, \mathbf{z}\right)=\Phi\left[\left(\mathbf{z}_{1} \boldsymbol{\delta}_{1}+\alpha_{1} y_{2}+\rho_{1} v_{2}\right) /\left(1-\rho_{1}^{2}\right)^{1 / 2}\right]$.

Since $y_{2}=1$ if and only if $v_{2}>-\mathbf{z} \boldsymbol{\delta}_{2}$, we need a basic fact about truncated normal distributions: If $v_{2}$ has a standard normal distribution and is independent of $\mathbf{z}$, then the density of $v_{2}$ given $v_{2}>-\mathbf{z} \boldsymbol{\delta}_{2}$ is\\
$\phi\left(v_{2}\right) / \mathrm{P}\left(v_{2}>-\mathbf{z} \boldsymbol{\delta}_{2}\right)=\phi\left(v_{2}\right) / \Phi\left(\mathbf{z} \boldsymbol{\delta}_{2}\right)$.

Therefore,


\begin{align*}
\mathrm{P}\left(y_{1}=1 \mid y_{2}=1, \mathbf{z}\right) & =\mathrm{E}\left[\mathrm{P}\left(y_{1}=1 \mid v_{2}, \mathbf{z}\right) \mid y_{2}=1, \mathbf{z}\right] \\
& =\mathrm{E}\left\{\Phi\left[\left(\mathbf{z}_{1} \boldsymbol{\delta}_{1}+\alpha_{1} y_{2}+\rho_{1} v_{2}\right) /\left(1-\rho_{1}^{2}\right)^{1 / 2}\right] \mid y_{2}=1, \mathbf{z}\right\} \\
& =\frac{1}{\Phi\left(\mathbf{z} \boldsymbol{\delta}_{2}\right)} \int_{-\mathbf{z} \boldsymbol{\delta}_{2}}^{\infty} \Phi\left[\left(\mathbf{z}_{1} \boldsymbol{\delta}_{1}+\alpha_{1} y_{2}+\rho_{1} v_{2}\right) /\left(1-\rho_{1}^{2}\right)^{1 / 2}\right] \phi\left(v_{2}\right) \mathrm{d} v_{2} \tag{15.55}
\end{align*}


where $v_{2}$ in the integral is a dummy argument of integration. Of course, $\mathbf{P}\left(y_{1}=\right. 0 \mid y_{2}=1, \mathbf{z}$ ) is just one minus equation (15.55).

Similarly, $\mathrm{P}\left(y_{1}=1 \mid y_{2}=0, \mathbf{z}\right)$ is


\begin{equation*}
\frac{1}{1-\Phi\left(\mathbf{z} \boldsymbol{\delta}_{2}\right)} \int_{-\infty}^{-\mathbf{z} \boldsymbol{\delta}_{2}} \Phi\left[\left(\mathbf{z}_{1} \boldsymbol{\delta}_{1}+\alpha_{1} y_{2}+\rho_{1} v_{2}\right) /\left(1-\rho_{1}^{2}\right)^{1 / 2}\right] \phi\left(v_{2}\right) \mathrm{d} v_{2} \tag{15.56}
\end{equation*}


Combining the four possible outcomes of $\left(y_{1}, y_{2}\right)$, along with the probit model for $y_{2}$, and taking the log gives the log-likelihood function for maximum likelihood analysis. Several authors, for example Greene (2003, Section 21.6), have noted a useful computational feature of the model in equations (15.51) and (15.52). To describe this feature-and to explain the additional flexibility it affords for extending the basic model-we introduce the bivariate probit model, typically specified for two binary responses as\\
$y_{1}=1\left[\mathbf{x}_{1} \boldsymbol{\beta}_{1}+e_{1}>0\right]$\\
$y_{2}=1\left[\mathbf{x}_{2} \boldsymbol{\beta}_{2}+e_{2}>0\right]$,\\
where $\mathbf{x}_{1}$ is $1 \times K_{1}$ and $\mathbf{x}_{2}$ is $1 \times K_{2}$. In the traditional formulation of bivariate probit, the error term, $\mathbf{e} \equiv\left(e_{1}, e_{2}\right)$, is assumed to be independent of $\left(\mathbf{x}_{1}, \mathbf{x}_{2}\right)$ with a bivariate normal distribution. In particular, $\mathbf{e} \mid \mathbf{x} \sim \operatorname{Normal}(\mathbf{0}, \boldsymbol{\Omega})$, where $\mathbf{x}$ consists of all exogenous variables and $\boldsymbol{\Omega}$ is the $2 \times 2$ matrix with ones down its diagonal and offdiagonal element $\rho=\operatorname{Corr}\left(e_{1}, e_{2}\right)$. These assumptions imply that $y_{1}$ and $y_{2}$ each follow probit models conditional on $\mathbf{x}$. Therefore, $\boldsymbol{\beta}_{1}$ and $\boldsymbol{\beta}_{2}$ can be consistently estimated by estimating separate probit models. Not surprisingly, if $e_{1}$ and $e_{2}$ are correlated, a joint maximum likelihood procedure is more efficient than the separate probits. With exogenous explanatory variables, increased efficiency in estimating $\boldsymbol{\beta}_{1}$ and $\boldsymbol{\beta}_{2}$ is the main reason for a joint estimation procedure. (Incidentally, when $\mathbf{x}_{1}= \mathbf{x}_{2}=\mathbf{x}$, there are generally efficiency gains from joint MLE over probit on each equation. By contrast, as we saw in Chapter 7 for the linear model, OLS on each equation and feasible GLS are identical.)

A simple way to obtain the log-likelihood function is to construct the joint density as $f\left(y_{1} \mid y_{2}, \mathbf{x}\right) f\left(y_{2} \mid \mathbf{x}\right)$, and it is here that a useful feature of the bivariate probit as it relates to probit with a binary endogenous variable emerges: the form of the conditional density, $f\left(y_{1} \mid y_{2}, \mathbf{x}\right)$, is the same even if $\mathbf{x}_{1}$ includes $y_{2}$. In fact, $\mathbf{x}_{1}$ can be any function of $\left(\mathbf{z}_{1}, y_{2}\right)$. In other words, in expressions (15.55) and (15.56) we can replace $\mathbf{z}_{1} \boldsymbol{\delta}_{1}+\alpha_{1} y_{2}$ with $\mathbf{x}_{1} \boldsymbol{\beta}_{1}=\mathbf{g}_{1}\left(\mathbf{z}_{1}, y_{2}\right) \boldsymbol{\beta}_{1}$. The reason we can allow this generality is simple: when we obtain the density of $y_{1}$ given $\left(y_{2}, \mathbf{x}\right)$, we are already conditioning on $y_{2}$, so the form of the density is the same whether or not $y_{2}$ is in $\mathbf{x}_{1}$. We discussed a similar feature of the log-likelihood function for estimating the model in Section 15.7.2.

The practical implication of our being able to use the bivariate probit loglikelihood function for endogenous explanatory variables is computational: if an econometrics package estimates a bivariate probit, then it can be used directly to estimate the parameters in (15.51) and (15.52). Further, we can use exactly the same routine to estimate a model with interactions in the structural equation, $y_{1}=1\left[\mathbf{z}_{1} \boldsymbol{\delta}_{1}+\right. \left.y_{2} \mathbf{z}_{1} \boldsymbol{\alpha}_{1}+u_{1}>0\right]$, with no more work than defining the interactions and including them in the appropriate command (specifying $y_{2}$ as the only endogenous variable but including the interactions with $y_{2}$ as additional explanatory variables).

Evans and Schwab (1995) use model (15.51) and (15.52) (and linear probability models) to estimate the causal effect of attending a Catholic high school ( $y_{2}$ ) on the probability of attending college ( $y_{1}$ ), allowing the Catholic high school indicator to be correlated with unobserved factors that also affect college attendance. As an IV for $y_{2}$ they use a binary indicator of whether a student is Catholic. Recently, Altonji, Elder, and Taber (2005) (AET for short) have revisited this question, using affiliation with the Catholic Church and geographic proximity to Catholic schools as instruments. They compare 2SLS estimates of a linear model - which, as always, may provide good approximations to the average partial effects - to those from the bivariate probit. In addition to finding the instruments to be suspect, AET conclude that identification of the parameters can be driven largely by the nonlinearity in the bivariate probit model. Unlike the model in Section 15.7.2, where an exclusion restriction in the structural equation is needed for identification, that is not the case in (15.51) and (15.52). Typically, one would be suspicious if an exclusion restriction is not available. AET's results suggest that even if such restrictions are available, they may have little impact on the estimated average treament effect (the APE of the Catholic high school dummy).

Because computing the MLE requires nonlinear optimization, it is tempting to use some seemingly "obvious" two-step procedures. As an example, we might try to inappropriately mimic 2SLS. Since $\mathrm{E}\left(y_{2} \mid \mathbf{z}\right)=\Phi\left(\mathbf{z} \boldsymbol{\delta}_{2}\right)$ and $\boldsymbol{\delta}_{2}$ is consistently esti-\\
mated by probit of $y_{2}$ on $\mathbf{z}$, it is tempting to estimate $\boldsymbol{\delta}_{1}$ and $\alpha_{1}$ from the probit of $y_{1}$ on $\mathbf{z}, \hat{\Phi}_{2}$, where $\hat{\Phi}_{2} \equiv \Phi\left(\mathbf{z} \hat{\boldsymbol{\delta}}_{2}\right)$. This approach does not produce consistent parameter estimators, for the same reasons the forbidden regression discussed in Section 9.5 for nonlinear simultaneous equations models does not. For this two-step procedure to work, we would have to have $\mathrm{P}\left(y_{1}=1 \mid \mathbf{z}\right)=\Phi\left[\mathbf{z}_{1} \boldsymbol{\delta}_{1}+\alpha_{1} \Phi\left(\mathbf{z} \boldsymbol{\delta}_{2}\right)\right]$. But $\mathrm{P}\left(y_{1}=1 \mid \mathbf{z}\right)= \mathrm{E}\left(y_{1} \mid \mathbf{z}\right)=\mathrm{E}\left(1\left[\mathbf{z}_{1} \boldsymbol{\delta}_{1}+\alpha_{1} y_{2}+u_{1}>0\right] \mid \mathbf{z}\right)$, and since the indicator function $1[\cdot]$ is nonlinear, we cannot pass the expected value through.

Greene (1998), in commenting on Burnett (1997), asserted that the two-step procedure that uses first-stage probit fitted values in place of $y_{2}$ in a second-stage probit is consistent but inefficient. As a way to obtain more efficient estimators, he proposed using the full MLE. But it is important to understand that the problem with the twostep procedure is not just a matter of inefficiency and adjusting for the two-step estimation; the procedure does not produce consistent estimators of the parameters, and probably not the APEs either (although a study of this issue would be informative). Burnett (1997) suggested adding the probit residuals, $\hat{r}_{i 2} \equiv y_{i 2}-\Phi\left(\mathbf{z}_{i} \hat{\boldsymbol{\delta}}_{2}\right)$, in an attempt to mimic Rivers and Vuong's (1988) method for a continuous $y_{2}$, although she applied nonlinear least squares, not MLE, in the second step. Interestingly, the Burnett approach does provide a valid test of the null hypothesis that $y_{2}$ is exogenous, although it makes more sense to use probit in the second stage, too. So, do a probit of $y_{i 1}$ on $\mathbf{z}_{i 1}, y_{i 2}, \hat{r}_{i 2}$ and obtain the usual asymptotic $t$ statistic on $\hat{r}_{i 2}$. Under the null hypothesis that $y_{2}$ is exogenous, $P\left(y_{1}=1 \mid \mathbf{z}, y_{2}\right)=P\left(y_{1}=1 \mid \mathbf{z}_{1}, y_{2}\right)= \Phi\left(\mathbf{z}_{1} \boldsymbol{\delta}_{1}+\alpha_{1} y_{2}\right)$, and so no additional functions of ( $y_{i 2}, \mathbf{z}_{i}$ ) should appear. Rather than use the residuals from the linear reduced form for $y_{2}$ (which is what the RiversVuong test does), we can use the residuals from a probit. Unfortunately, if $y_{2}$ is endogenous, Burnett's suggestion does not appear to consistently estimate the parameters or the APEs.

As mentioned in the previous subsection, we can use the Rivers-Vuong approach to test for exogeneity of $y_{2}$. This has the virtue of being simple, and, if the test fails to reject, we may not need to compute the MLE. A more efficient test is the score test of $\mathrm{H}_{0}: \rho_{1}=0$, and this does not require estimation of the full MLE. Of course, if one has computed the MLE, a $t$ test or $L R$ test can be used.

Example 15.4 (Women's Labor Force Participation and Having More than Two Children): We use a data set from Angrist and Evans (1998), in LABSUP.RAW, to study the effects of having more than two children on women's labor force participation decisions. The population is married women in the United States who have at least two children. The endogenous explanatory variable is $y_{2}=$ morekids, which is unity if a woman has three or more children (about 49 percent of the sample). The response variable is $y_{1}=$ worked; roughly 59 percent of the women report being in

Dependent Variable: worked

\begin{table}[h]
\begin{center}
\captionsetup{labelformat=empty}
\caption{Table 15.2\\
Estimated Effect of Having Three or More Children on Women's Labor Force Participation}
\begin{tabular}{|l|l|l|l|l|l|}
\hline
 & (1) & (2) & (3) & (4) & (5) \\
\hline
Model & LPM & Probit & LPM & Bivariate probit & Bivariate probit \\
\hline
Estimation method & OLS & MLE & 2SLS: samesex as IV & MLE: samesex as IV & MLE: no IV \\
\hline
Coefficient on morekids & -. 109 (.006) & -. 299 (.015) & -. 201 (.096) & -. 703 (.204) & -. 966 (.243) \\
\hline
APE for morekids & -. 109 (.006) & -. 109 (.006) & -. 201 (.096) & -. 256 (.072) & -. 349 (*) \\
\hline
$\hat{\rho}$ & - & - & - & . 254 (.131) & . 426 (.162) \\
\hline
Number of observations & 31,857 & 31,857 & 31,857 & 31,857 & 31,857 \\
\hline
\end{tabular}
\end{center}
\end{table}

Standard errors for estimated coefficients and APEs are given in parentheses next to coefficients. For the nonlinear models, the APE standard errors were obtained from 500 bootstrap samplings.\\
A bootstrap standard error could not be obtained for column (5) due to computational problems for some bootstrap samples.\\
the labor force at the time of the survey. We also include the variables nonmomi ("non-mom" income), educ (years of schooling), age, age ${ }^{2}$, and the race indicators black and hispan; these are all treated as exogenous. As an instrumental variable for morekids we use samesex, which is a binary variable equal to one if the first two children are of the same sex. (While the outcome of samesex is legitimately treated as random, it is not necessarily exogenous to the labor force participation decision: having two children of the same sex can shift a family's budget constraint because, for example, a bedroom is more easily shared and clothes and toys are more easily handed down.) Table 15.2 contains the estimation results from five different approaches: OLS estimation of a linear probability model, probit treating morekids as exogenous, 2SLS estimation of an LPM, bivariate probit allowing morekids to be endogenous, and bivariate probit that drops samesex from the probit for morekids. We include the latter primarily to illustrate that the nonlinearity in the model identifies the parameters. For brevity, we only report the coefficients and average partial effects for morekids. For the LPMs, the standard errors are robust to arbitrary heteroskedasticity.

When morekids is assumed exogenous, the LPM and probit models give the same estimated average partial effect of having more than two children to three decimal places: on average, women with more than two children are about .11 less likely to be in the labor force than women with two children. In addition, the standard errors are the same to three decimal places (and the standard error for probit was obtained from bootstrapping). When samesex is used as an IV in the LPM estimation, the effect almost doubles. Column 4 contains the estimates from the bivariate probit\\
where morekids is treated as endogenous and samesex is excluded from the structural equation but appears in the probit model for morekids-the standard case where we impose an exclusion restriction. Now, the nonlinear probit model gives a larger estimated APE than the linear model: -.256 versus -.201 . Interestingly, according to the APE standard error for the bivariate probit, obtained using bootstrapping, the APE is more precisely estimated using the nonlinear model.

If we drop samesex from the probit model for morekids-which, in a linear context, would lead to a lack of identification-we estimate an even larger effect: the APE is -.349 . The estimated value of $\rho$ increases substantially when we drop samesex from the model for morekids, suggesting that including samesex in the model for morekids helps reduce the correlation between unobservables that affect both morekids and worked. The large increase in the magnitude of the APE suggests that the nonlinearity in the bivariate probit plays a critical role in determining the APE, a point made by Altonji, Elder, and Taber (2005) when studying the effects of attending a Catholic high school on outcomes such as attending college. Generally, it is dangerous to rely on nonlinearities to identify parameters and partial effects, and one should avoid doing so in bivariate probit models. When bootstrapping was applied to obtain a standard error for the APE in column (5), for some bootstrap samples the computational algorithm would not converge, suggesting that the model without an exclusion restriction is ill specified. In this application, the real issue should be whether the difference in APEs between the linear model estimated by 2SLS and the bivariate probit that uses the same instrument is important.

We can also implement the inconsistent two-step estimation approach where the probit fitted values, say morekids, are plugged in for morekids in the second-stage probit. The coefficient on morekids is about about -.843 , which is quite a bit larger in magnitude than the full MLE estimate when samesex is used as an IV, -.703 . The estimated APE for the two-step procedure is about -.305 , which again is larger in magnitude than the full MLE estimate, -.256. The difference in APEs is less than the difference in APEs between the linear 2SLS estimate and full MLE. It is intriguing to think that, somehow, the inconsistent two-step procedure might generally deliver reasonable estimates of the APEs, but there is no evidence on this point.

\subsection*{15.7.4 Heteroskedasticity and Nonnormality in the Latent Variable Model}
In applying the probit model it is easy to become confused about the problems of heteroskedasticity and nonnormality. The confusion stems from a failure to distinguish between the underlying latent variable formulation, as in $y^{*}=\mathbf{x} \boldsymbol{\beta}+e$, and the response probability in equation (15.8). As we have emphasized throughout\\
this chapter, for most purposes we want to estimate $\mathrm{P}(y=1 \mid \mathbf{x})$. The latent variable formulation is convenient for certain manipulations, but we are rarely interested in $\mathrm{E}\left(y^{*} \mid \mathbf{x}\right)$. (Data censoring cases, where we are only interested in the parameters of an underlying linear model, are treated in Chapter 19.)

Once we focus on the response probability, we can easily see why comparing nonnormality and heteroskedasticity in the latent variable model with the same problems in a linear (or nonlinear) regression requires considerable care. First consider the problem of nonnormality of $e$ in a probit model. If $e$ is independent of $\mathbf{x}$, we can write $\mathrm{P}(y=1 \mid \mathbf{x})=G(\mathbf{x} \boldsymbol{\beta}) \neq \Phi(\mathbf{x} \boldsymbol{\beta})$, where, generally, $G(z)=1-F(-z)$ and $F(\cdot)$ is the cdf of $e$. One often hears statements such as "nonnormality in a probit model causes bias and inconsistency in the parameter estimates." Technically, this statement is correct, but it largely misses the point. In Section 15.6 we noted that when $\mathbf{x}$ has a multivariate normal distribution, the LPM consistently estimates the APEs for any smooth function $G(\cdot)$. Using a linear model when the response probability is nonlinear is a fairly serious form of misspecification, yet we consistently estimate perhaps the most useful quantities: the APEs. The probit model is likely to do a reasonable job of approximating the APES in lots of cases where $G \neq \Phi$. In fact, the situation that emerges in Example 15.2 is quite common in applications: probit and logit give very similar estimated partial effects. That the logit parameter estimates are larger than the probit estimates by roughly a factor of 1.6 is a consequence of the different implicit scale factors, but the different scalings are easily accounted for by computing partial effects.

Worrying about the choice of $G$ when the response probability is of interest is no different from worrying about the functional form for $\mathrm{E}(y \mid \mathbf{x})$ in a regression context. For example, if $y \geq 0$ and we are considering two models for $\mathrm{E}(y \mid \mathbf{x})$, a linear model $\mathrm{E}(y \mid \mathbf{x})=\mathbf{x} \boldsymbol{\beta}$ and an exponential model $\mathrm{E}(y \mid \mathbf{x})=\exp (\mathbf{x} \boldsymbol{\beta})$, we do not couch the choice between them in the language of "biased" or "inconsistent" parameter estimation. That OLS estimation of a linear model will not consistently estimate the parameters in $\exp (\mathbf{x} \boldsymbol{\beta})$ is both obvious and not particularly relevant. The issue is whether a linear or exponential function form provides the best fit, and whether the estimated partial effects $\mathrm{E}(y \mid \mathbf{x})$ are very different across the models. In some cases, the linear model does provide good estimates of the APEs. We should have the same discussion when deciding on $G$ in the index model $\mathrm{P}(y=1 \mid \mathbf{x})=\mathrm{E}(y \mid \mathbf{x})=G(\mathbf{x} \boldsymbol{\beta})$ : our choice of $G$ is a functional form issue. A nonnormal distribution of $u$ in the linear regression model $y=\mathbf{x} \boldsymbol{\beta}+u, \mathrm{E}(u \mid \mathbf{x})=0$, does not change the functional form of $\mathrm{E}(y \mid \mathbf{x})$, and that is why nonnormality is harmless (provided one can rely on asymptotic theory). Nonnormality in $e$ in $y=1[\mathbf{x} \boldsymbol{\beta}+e>0]$ means that the probit\\
response probability is incorrect. Even if we could estimate $\boldsymbol{\beta}$ consistently we could not obtain magnitudes of the partial effects. So our focus should be on how well different methods approximate partial effects, and not on whether they estimate parameters consistently. (In the next subsection, we show that relative partial effects of continuous variables can be identified without specifying $G(\cdot)$, and sometimes even if we allow $e$ and $\mathbf{x}$ to have some dependence.)

Once we view nonnormality of $e$ as a functional form problem for $p(\mathbf{x})$, we can evaluate more general parametric models in a sensible way. It can be a good idea to replace $\Phi(\mathbf{x} \boldsymbol{\beta})$ with a function such as $G(\mathbf{x} \boldsymbol{\beta}, \gamma)$, where $\gamma$ is an extra set of parameters, especially if $G(\mathbf{x} \boldsymbol{\beta}, \boldsymbol{\gamma})$ is chosen to nest the probit model. (Moon (1988) covers some interesting possibilities in the context of logit models, including asymmetric distributions. See also Problem 15.16.) But generalizing functional form by making the distribution of $e$ more general is not necessarily better than just specifying more flexible models for the response probability directly, as in McDonald (1996). And we definitely should not reject the standard models just because the estimates of $\boldsymbol{\beta}$ seem to change a lot: the basis for comparison should be partial effects at various values of $\mathbf{x}$, APEs, and goodness-of-fit measures such as the values of the log-likelihood functions and the percent correctly predicted.

Similar comments can be made about heteroskedasticity in the latent error $e$. Traditionally, discussions of its implications tend to suffer from a failure to specify what it is we hope to learn when we estimate binary response models. For example, one often sees reference to Yatchew and Griliches (1984) concerning the inconsistency of the probit MLE when $\operatorname{Var}(e \mid \mathbf{x})$ depends on $\mathbf{x}$, even if $\mathrm{D}(e \mid \mathbf{x})$ is normal. If $\mathrm{D}(e \mid \mathbf{x})=\operatorname{Normal}(0, h(\mathbf{x}))$, the response probability takes the form $\mathrm{P}(y=1 \mid \mathbf{x})= \Phi\left[\mathbf{x} \boldsymbol{\beta} / h(\mathbf{x})^{1 / 2}\right]$, and so it is pretty obvious that probit of $y$ on $\mathbf{x}$ could not consistently estimate $\boldsymbol{\beta}$. Some authors have noted this finding is particularly troubling because "heteroskedasticity is prevalent with microeconomic data." But what does this mean? Even if the usual probit model is correct-so that $\operatorname{Var}(e \mid \mathbf{x})=1$-the variance of the response variable, $y$, is heteroskedastic: $\operatorname{Var}(y \mid \mathbf{x})=\Phi(\mathbf{x} \boldsymbol{\beta})[1-\Phi(\mathbf{x} \boldsymbol{\beta})]$. In fact, most observed discrete random variables $y$ will have conditional distributions such that $\operatorname{Var}(y \mid \mathbf{x})$ is not constant, but that is due to the discreteness in $y$, not necessarily because of heteroskedasticity in an underlying latent error.

Is there any reason we should consider possible heteroskedasticity in $\operatorname{Var}\left(y^{*} \mid \mathbf{x}\right)$ when $y^{*}$ is an unknowable latent variable? There are two reasons. The simplest is for generalizing the functional form of the response probability-just as when we allow for nonnormality in $\mathrm{D}\left(y^{*} \mid \mathbf{x}\right)$. In Section 15.5.3 we discussed testing the probit model against an alternative where $\operatorname{Var}(e \mid \mathbf{x})=\exp \left(2 \mathbf{x}_{1} \boldsymbol{\delta}\right)$ for $\mathbf{x}_{1}$ a subset of $\mathbf{x}$ (which, at a\\
minimum, excludes a constant). Then $\mathrm{P}(y=1 \mid \mathbf{x})=\Phi\left[\exp \left(-\mathbf{x}_{1} \boldsymbol{\delta}\right) \mathbf{x} \boldsymbol{\beta}\right]$, and so we have potentially allowed a much more flexible response probability. But there is a cost to the more general model: the partial effects, $\partial \mathbf{P}(y=1 \mid \mathbf{x}) / \partial x_{j}$ are more complicated. If $\delta_{j}$ is the coefficient on $x_{j}$ in the vector $\mathbf{x}_{1}$, then\\
$\frac{\partial \partial \mathrm{P}(y=1 \mid \mathbf{x})}{\partial x_{j}}=\phi\left[\exp \left(-\mathbf{x}_{1} \boldsymbol{\delta}\right) \mathbf{x} \boldsymbol{\beta}\right] \exp \left(-\mathbf{x}_{1} \boldsymbol{\delta}\right)\left[\boldsymbol{\beta}_{j}-\delta_{j} \cdot(\mathbf{x} \boldsymbol{\beta})\right]$,\\
and so the sign of the partial effect at $\mathbf{x}$ depends on the sign of $\boldsymbol{\beta}_{j}-\delta_{j} \cdot(\mathbf{x} \boldsymbol{\beta})$; it need not be the same sign as $\beta_{j}$. (We can actually find models where the partial effect is always the opposite sign from the coefficient. If $y=1\left[\beta_{0}+\beta_{1} x_{1}+e>0\right]$ and $\mathrm{D}(e \mid \mathbf{x})=\operatorname{Normal}\left(0, x_{1}^{2}\right)$, then $\mathrm{P}\left(y=1 \mid x_{1}\right)=\Phi\left(\beta_{0} / x_{1}+\beta_{1}\right)$, and so $\partial \mathrm{P}\left(y=1 \mid x_{1}\right) / \partial x_{1}=-\left(\beta_{0} / x_{1}^{2}\right) \phi\left(\beta_{0} / x_{1}+\beta_{1}\right)$. If $\beta_{0}>0$ and $\beta_{1}>0$, the partial effect of $x_{1}$ on $\mathrm{P}\left(y=1 \mid x_{1}\right)$ is negative while $\beta_{1}$, the partial effect of $x_{1}$ on $\mathrm{E}\left(y^{*} \mid x_{1}\right)$, is positive.) These days, estimating a so-called heteroskedastic probit model (which usually means an exponential conditional variance for $e$ ) is fairly straightforward. The challenge is in intelligently using the resulting estimates. One possibility is to compute partial effects from equation (15.57) and average them across the sample. In any case, it makes no sense to compare the estimates of $\boldsymbol{\beta}$ from a heteroskedastic probit with those from a standard probit. In each case the MLE estimates adjust to fit the data, and this often means that different methods can provide similar partial effects, at least for some values of $\mathbf{x}$.

There is another, more subtle reason to entertain the notion of heteroskedasticity in the latent variable model. As discussed in Wooldridge (2005c), the partial effects in the model $y_{i}=1\left[\mathbf{x}_{i} \boldsymbol{\beta}+e_{i}>0\right]$, averaged across the distribution of $e_{i}$, are the same sign as the corresponding $\beta_{j}$. In fact, the ratios of APEs for the continuous variables are equal to the ratios of the coefficients, even when $e_{i}$ contains heteroskedasticity. To see why, it is easiest to work with the average structural function (ASF) defined by Blundell and Powell (2004); see also Section 2.2.5. For binary response, the ASF as a function of $\mathbf{x}$ is


\begin{equation*}
\operatorname{ASF}(\mathbf{x})=\mathrm{E}_{e_{i}}\left\{1\left[\mathbf{x} \boldsymbol{\beta}+e_{i}>0\right]\right\}=\mathrm{P}\left(e_{i}>-\mathbf{x} \boldsymbol{\beta}\right)=1-F(-\mathbf{x} \boldsymbol{\beta}), \tag{15.58}
\end{equation*}


where $F(\cdot)$ is the cdf of $e_{i}$ (which may not be symmetrically distributed about zero), and we put an $i$ subscript on $e_{i}$ to emphasize it is the random variable that is averaged out. From expression (15.58), we can see that the APE of $x_{j}$ has the same sign as $\beta_{j}$, and, for a continuous $x_{j}$, the APE is simply $\beta_{j} f(-\mathbf{x} \boldsymbol{\beta})$, assuming that $F$ is continuously differentiable with density $f$. The relative APEs for two continuous variables $x_{j}$ and $x_{h}$ is $\beta_{j} / \beta_{h}$-the same conclusion in Section 15.7.1 when we introduced heterogeneity additively inside the probit function that was independent of $\mathbf{x}$.

That the sign of the APE for $x_{j}$ is given by the sign of $\beta_{j}$, and that the relative APEs for continuous covariates are given by the ratios of elements of $\boldsymbol{\beta}$-whether or not $e$ is independent of $\mathbf{x}$-suggest that $\boldsymbol{\beta}$ is still of some interest. In the next subsection we will discuss estimation of $\boldsymbol{\beta}$ under weaker assumptions (up to a scale factor). But what implications does this discussion about APEs have for heteroskedastic probit? Interestingly, in the heteroskedastic probit model, we can easily recover the APEs. The easiest way is through the ASF. Now, rather than using the unconditional distribution of $e_{i}$, we use iterated expectations, because what we have modeled is $\mathrm{D}\left(e_{i} \mid \mathbf{x}_{i}\right)=\mathrm{D}\left(e_{i} \mid \mathbf{x}_{i 1}\right)$ :


\begin{equation*}
\operatorname{ASF}(\mathbf{x})=\mathrm{E}_{\mathbf{x}_{i 1}}\left(\mathrm{E}\left\{1\left[\mathbf{x} \boldsymbol{\beta}+e_{i}>0\right] \mid \mathbf{x}_{i 1}\right\}\right)=\mathrm{E}_{\mathbf{x}_{i 1}}\left\{\Phi\left[\exp \left(-\mathbf{x}_{i 1} \boldsymbol{\delta}\right) \mathbf{x} \boldsymbol{\beta}\right]\right\}, \tag{15.59}
\end{equation*}


where $\mathrm{E}\left\{1\left[\mathbf{x} \boldsymbol{\beta}+e_{i}>0\right] \mid \mathbf{x}_{i 1}\right\}=\Phi\left[\exp \left(-\mathbf{x}_{i 1} \boldsymbol{\delta}\right) \mathbf{x} \boldsymbol{\beta}\right]$ follows from $1\left[\mathbf{x} \boldsymbol{\beta}+e_{i}>0\right]= 1\left[\exp \left(-\mathbf{x}_{i 1} \boldsymbol{\delta}\right) e_{i}>-\exp \left(-\mathbf{x}_{i 1} \boldsymbol{\delta}\right) \mathbf{x} \boldsymbol{\beta}\right]$ along with $\mathrm{D}\left(\exp \left(-\mathbf{x}_{i 1} \boldsymbol{\delta}\right) e_{i} \mid \mathbf{x}_{i 1}\right)=\operatorname{Normal}(0,1)$. It is important to see that $\mathbf{x}$ is just a fixed argument in (15.59), whereas $\mathbf{x}_{i 1}$ denotes a random vector that we average out. For a continuous covariate $x_{j}$, the partial effect on the ASF is


\begin{equation*}
\mathrm{E}_{\mathbf{x}_{i l}}\left\{\phi\left[\exp \left(-\mathbf{x}_{i l} \boldsymbol{\delta}\right) \mathbf{x} \boldsymbol{\beta}\right]\right\} \beta_{j} . \tag{15.60}
\end{equation*}


Given the maximum likelihood estimators $\hat{\boldsymbol{\beta}}$ and $\hat{\boldsymbol{\delta}}$ from heteroskedastic probit, a consistent estimator of (15.60) is $\left(N^{-1} \sum_{i=1}^{N} \phi\left[\exp \left(-\mathbf{x}_{i 1} \hat{\boldsymbol{\delta}}\right) \mathbf{x} \hat{\boldsymbol{\beta}}\right]\right) \hat{\beta}_{j}$, and then we can insert interesting values of $\mathbf{x}$ or further average out across $\mathbf{x}_{i}$. Bootstrapping would be a sensible method for obtaining valid standard errors of the APEs.

Now we seem to be in a quandary. We have two ways of computing partial effects, and they can give conflicting answers, not just on magnitude of the effects but also on the direction of the effects. Initially, we discussed how the partial effects based on $\mathrm{P}(y=1 \mid \mathbf{x})$-given for continuous $x_{j}$ in equation (15.57)-need not have the same signs as the $\beta_{j}$. But equation (15.60) shows that the APEs obtained from averaging out $e_{i}$ are, in fact, proportional to the $\beta_{j}$. (Of course, to obtain the partial effects in either case we cannot ignore $\boldsymbol{\delta}$, as it appears directly in both formulas.) Which partial effect is "correct"? Unfortunately, there is a fundamental lack of identification that does not allow us to choose between the two. Equation (15.60) was obtained from $y_{i}=1\left[\mathbf{x}_{i} \boldsymbol{\beta}+e_{i}>0\right]$ where $e_{i}$ given $\mathbf{x}_{i}$ has a heteroskedastic normal distribution. But suppose we start with the model $y_{i}=1\left[\mathbf{x}_{i} \boldsymbol{\beta}+\exp \left(\mathbf{x}_{i 1} \boldsymbol{\delta}\right) a_{i}>0\right]$, where $a_{i}$ is independent of $\mathbf{x}_{i}$ with a standard normal distribution. Now, it is easily shown that the ASF is just the same as the response probability evaluated at $\mathbf{x}$, namely, $\Phi\left[\exp \left(-\mathbf{x}_{1} \boldsymbol{\delta}\right) \mathbf{x} \boldsymbol{\beta}\right]$, which takes us back to the partial effects in equation (15.57). In fact, the APEs based on (15.57) are commonly reported for heteroskedastic probit, whereas we have just as good a case for using (15.60).

As discussed by Wooldridge (2005a), on aesthetic grounds we might prefer the usual index structure with $\operatorname{Var}(e \mid \mathbf{x})$ heteroskedastic, especially because the APEs are easy to obtain and the $\beta_{j}$ give directions of effects and relative effects. But both models nest the usual probit model, and aesthetic considerations cannot change the fact that two observationally equivalent models lead to possibly quite different APEs. Plus, there is nothing unappealing about index models of the form $y_{i}=1\left(\mathbf{x}_{i} \boldsymbol{\beta}+a_{i}+\right. a_{i} \mathbf{x}_{i 1} \boldsymbol{\delta}>0$ ) where $\mathrm{D}\left(a_{i} \mid \mathbf{x}_{i}\right)=\operatorname{Normal}(0,1)$; in fact, interacting the unobservable $a_{i}$ with a subset of $\mathbf{x}_{i}$ makes perfectly good sense, and we studied linear models with this feature in Chapters 4 and 6. Because $a_{i}$ is independent of $\mathbf{x}_{i}$, the partial effects defined in terms of the ASF are, like those in equation (15.57), the same as the partial effects on the response probability $p(\mathbf{x})$. The uncomfortable conclusion is that we have no convincing way of choosing between equations (15.57) and (15.60).

\subsection*{15.7.5 Estimation under Weaker Assumptions}
Probit, logit, and the extensions of these mentioned in the previous subsection are all parametric models: $\mathrm{P}(y=1 \mid \mathbf{x})$ depends on a finite number of parameters. There have been many advances in estimation of binary response models that relax parametric assumptions on $\mathrm{P}(y=1 \mid \mathbf{x})$. We briefly discuss some of those here.

If we are interested in estimating the directions and relative sizes of the partial effects, and not the response probabilities, several results are known. In Section 15.6 we noted a special case- $\mathbf{x}$ is multivariate normal-where the slope coefficients in the linear projection of $y$ on $1, \mathbf{x}$, say $\boldsymbol{\lambda}$, are the partial effects of $p(\mathbf{x})$ averaged across the distribution of $\mathbf{x}$-regardless of the true response probability! Chung and Goldberger (1984) obtain conditions under which the linear projection identifies the parameters up to scale. In fact, in the index formulation $y=1\left[y^{*}>0\right]$ with $y^{*}=\alpha+\mathbf{x} \boldsymbol{\beta}+e$, Chung and Goldberger simply take $\alpha$ and $\boldsymbol{\beta}$ to be the parameters in the linear projection of $y^{*}$ on $1, \mathbf{x}$. Their main result is that if $\mathrm{E}\left(\mathbf{x} \mid y^{*}\right)$ is linear in $y^{*}$-sufficient but not necessary is that $\left(\mathbf{x}, y^{*}\right)$ is multivariate normal-then the slope parameters in the linear projection of $y$ on $1, \mathbf{x}$ are proportional to $\boldsymbol{\beta}: \boldsymbol{\lambda}=\tau \boldsymbol{\beta}$, where $\tau$ is the slope in the linear projection of $y$ on $y^{*}$. (For binary response, it can be shown that $\tau>0$.) In the standard index model $p(\mathbf{x})=G(\alpha+\mathbf{x} \boldsymbol{\beta})$, the partial effect for continuous covariates are proportional to their betas, and so the linear projection identifies the relative partial effects of continuous explanatory variables. Unfortunately, we cannot conclude that the $\lambda$ are themselves the APEs unless $\mathbf{x}$ is multivariate normal.

Ruud (1983) obtains similar results for the index models estimated by maximum likelihood, but where we misspecify $G(\cdot)$-thereby employing quasi-MLE. Ruud obtains a result that the slopes in the misspecified MLE are consistent for $\tau \boldsymbol{\beta}$ for an unknown scale factor $\tau$. Ruud's key condition is linearity of the conditional means\\
$\mathrm{E}(\mathbf{x} \mid \mathbf{x} \boldsymbol{\beta})$; multivariate normality of $\mathbf{x}$ is sufficient but not necessary. Of course, this condition is unlikely to be generally satisfied. Ruud (1986) shows how to exploit these results to consistently estimate the slope parameters up to scale fairly generally.

An alternative approach is to explicitly recognize that we do not know the function $G(\cdot)$, but the response probability has the index form in equation (15.8). This arises from the latent variable formulation (15.9) when $e$ is independent of $\mathbf{x}$ but the distribution of $e$ is not known. There are several semiparametric estimators of the slope parameters, up to scale, that do not require knowledge of $G$. Under certain restrictions on the function $G$ and the distribution of $\mathbf{x}$, the semiparametric estimators are consistent and $\sqrt{N}$-asymptotically normal. See, for example, Stoker (1986), Powell, Stock, and Stoker (1989), Ichimura (1993), Klein and Spady (1993), and Ai (1997). Powell (1994) contains a survey of these methods.

Once $\hat{\boldsymbol{\beta}}$ is obtained, the function $G$ can be consistently estimated (in a sense we cannot make precise here, as $G$ is part of an infinite dimensional space). Thus, the response probabilities, as well as the partial effects on these probabilities, can be consistently estimated for unknown $G$. Obtaining $\hat{G}$ requires nonparametric regression of $y_{i}$ on $\mathbf{x}_{i} \hat{\boldsymbol{\beta}}$, where $\hat{\boldsymbol{\beta}}$ are the scaled slope estimators. Accessible treatments of the methods used are contained in Stoker (1992), Powell (1994), and Härdle and Linton (1994).

Remarkably, it is possible to estimate $\boldsymbol{\beta}$ up to scale without assuming that $e$ and $\mathbf{x}$ are independent in the model (15.9). In the specification $y=1[\mathbf{x} \boldsymbol{\beta}+e>0]$, Manski $(1975,1988)$ shows how to consistently estimate $\boldsymbol{\beta}$, subject to a scaling, under the assumption that the median of $e$ given $\mathbf{x}$ is zero. Some mild restrictions are needed on the distribution of $\mathbf{x}$; the most important of these is that at least one element of $\mathbf{x}$ with nonzero coefficient is essentially continuous. This allows $e$ to have any distribution, and $e$ and $\mathbf{x}$ can be dependent; for example, $\operatorname{Var}(e \mid \mathbf{x})$ is unrestricted. Manski's estimator, called the maximum score estimator, is a least absolute deviations estimator. Since the median of $y$ given $\mathbf{x}$ is $1[\mathbf{x} \boldsymbol{\beta}>0]$, the maximum score estimator solves

$$
\min _{\boldsymbol{\beta}} \sum_{i=1}^{N}\left|y_{i}-1\left[\mathbf{x}_{i} \boldsymbol{\beta}>0\right]\right|
$$

over all $\boldsymbol{\beta}$ with, say, $\boldsymbol{\beta}^{\prime} \boldsymbol{\beta}=1$, or with some element of $\boldsymbol{\beta}$ fixed at unity if the corresponding $x_{j}$ is known to appear in $\operatorname{Med}(y \mid \mathbf{x})$. (A normalization is needed because if $\operatorname{Med}(y \mid \mathbf{x})=1[\mathbf{x} \boldsymbol{\beta}>0]$ then $\operatorname{Med}(y \mid \mathbf{x})=1[\mathbf{x}(\tau \boldsymbol{\beta})>0]$ for any $\tau>0$.) The resulting estimator is consistent-for a recent proof, see Newey and McFadden (1994)-but its limiting distribution is nonnormal. In fact, it converges to its limiting distribution at rate $N^{1 / 3}$. Horowitz (1992) proposes a smoothed version of the maximum score estimator that converges at a rate close to $\sqrt{N}$.

The maximum score estimator's strength is that it consistently estimates $\boldsymbol{\beta}$ up to scale in cases where the index model (15.8) does not hold. As we saw in the previous subsection, when $y_{i}=1\left[\mathbf{x}_{i} \boldsymbol{\beta}+e_{i}>0\right]$, the APE for $x_{j}$ has same sign as $\beta_{j}$ and the relative APEs for continuous variables are given by $\beta_{j} / \beta_{h}$. Thus, there is some value in estimating the parameters up to a common scale factor. But maximum score estimation does not allow estimation of the APEs for either continuous or discrete covariates because the unconditional distribution of $e_{i}$ is not identified and, without making further assumptions, we cannot find $\mathrm{P}\left(y_{i}=1 \mid \mathbf{x}_{i}\right)$. We should not be too surprised by this state of affairs: if our assumptions are weak, we might not be able to learn everything we would like to know about how $\mathbf{x}$ affects $y$.

Lewbel (2000) offers a different approach to estimating the coefficients up to scale when $e$ is not independent of $\mathbf{x}$ but $e$ and $\mathbf{x}$ are uncorrelated. The identification results and estimation methods can be found in Lewbel's paper; here we discuss the key assumptions and their applicability. Lewbel's key assumption is the existence of a continuous variable, say $x_{K}$, with $\beta_{K} \neq 0$, such that $\mathrm{D}(e \mid \mathbf{x})=\mathrm{D}\left(e \mid x_{1}, x_{2}, \ldots, x_{K-1}\right)$. In other words, conditional on $\left(x_{1}, \ldots, x_{K-1}\right), e$ and $x_{K}$ are independent. Applied to the case of heteroskedasticity in the latent variable model, practically speaking the conditional independence assumption means $\mathrm{E}\left(y^{*} \mid \mathbf{x}\right)$ must depend on $x_{K}$ but $\operatorname{Var}\left(y^{*} \mid \mathbf{x}\right)$ must not depend on $x_{K}$. Therefore, one must be willing to assume that a variable definitely affects the conditional mean of $y^{*}$ but not its conditional variance. If we parametrically model $\mathrm{D}(e \mid \mathbf{x})$, we need not impose such a restriction, and so the semiparametric approach does not uniformly improve on parametric approaches. (And, as we already know, we can only learn about relative sizes of the coefficients using the semiparametric approach, whereas the parametric approach also delivers APEs.)

As pointed out by Lewbel, one case where his assumption holds is when the latent variable model is a random coefficient model of the form $y_{i}^{*}=b_{i 1} x_{i 1}+\cdots+ b_{i, K-1} x_{i, K-1}+\beta_{K} x_{i K}+a_{i}$, where $b_{i 1}, \ldots, b_{i, K-1}$ are random slopes and $\beta_{K} \neq 0$. If the vector ( $a_{i}, b_{i 1}, \ldots, b_{i, K-1}$ ) is independent of $\mathbf{x}_{i}$, then Lewbel's key assumption holds. The problem is that we have arbitrarily restricted $x_{i K}$ to have a constant, nonzero coefficient. Even if $x_{i K}$ is generated randomly and independently of the other covariates and the unobserved heterogeneity, there is no reason to think that its coefficient is constant while other regressors might have heterogeneous coefficients; it is simply an arbitrary restriction. (For example, if $y^{*}$ is a latent variable measuring propensity to be employed and $x_{K}$ is a job training indicator, random assignment of $x_{K}$ has nothing to do with whether job training has differential effects across individuals on the propensity to be employed.) If we focus on partial effects, we can allow the entire vector $\mathbf{b}_{i}=\left(b_{i 1}, \ldots, b_{i K}\right)^{\prime}$ to vary along with $a_{i}$. Then, we know that the APEs on\\
$\mathrm{P}\left(y_{i}=1 \mid \mathbf{x}_{i}, a_{i}, \mathbf{b}_{i}\right)$ are easily obtained from $\mathrm{P}\left(y_{i}=1 \mid \mathbf{x}_{i}\right)$; see Section 2.2.5. Therefore, we can directly specify a flexible model for $\mathrm{P}\left(y_{i}=1 \mid \mathbf{x}_{i}\right)$ or derive the response probability from $\mathrm{P}\left(y_{i}=1 \mid \mathbf{x}_{i}, a_{i}, \mathbf{b}_{i}\right)$, such as probit, and a distribution for $\left(a_{i}, \mathbf{b}_{i}^{\prime}\right)^{\prime}$, such as multivariate normal.

Some progress has been made in estimating parameters up to scale in the model $y_{1}=1\left[\mathbf{z}_{1} \boldsymbol{\delta}_{1}+\alpha_{1} y_{2}+u_{1}>0\right]$, where $y_{2}$ might be correlated with $u_{1}$ and $\mathbf{z}_{1}$ is a $1 \times L_{1}$ vector of exogenous variables. Lewbel's (2000) general approach applies to this situation as well. Let $\mathbf{z}$ be the vector of all exogenous variables uncorrelated with $u_{1}$. Then Lewbel requires a continuous element of $\mathbf{z}_{1}$ with nonzero coefficient-say, the last element, $z_{L_{1}}$-that does not appear in $\mathrm{D}\left(u_{1} \mid y_{2}, \mathbf{z}\right)$. (Clearly, $y_{2}$ cannot play the role of the variable excluded from $\mathrm{D}\left(u_{1} \mid y_{2}, \mathbf{z}\right)$ if $y_{2}$ is thought to be endogenous.) When might Lewbel's exclusion restriction hold? Sufficient is $y_{2}=g_{2}\left(\mathbf{z}_{2}\right)+v_{2}$, where ( $u_{1}, v_{2}$ ) is independent of $\mathbf{z}$ and $\mathbf{z}_{2}$ does not contain $z_{L_{1}}$. But this means that we have imposed an exclusion restriction on the reduced form of $y_{2}$, something usually discouraged in parametric contexts. (Putting restrictions on reduced forms is much different, and less defensible, than putting retrictions on structural equations.) Randomization of $z_{L_{1}}$ does not make its exclusion from the reduced form of $y_{2}$ legitimate. In fact, one often hopes that an instrument for $y_{2}$ is effectively randomized, which means that $z_{L_{1}}$ does not appear in the structural equation but does appear in the reduced form of $y_{2}$-the opposite of Lewbel's assumption. (As a generic example, eligibility is often randomized but participation, $y_{2}$, is not, and then one hopes to use an eligibility dummy as an IV for the participation dummy.)

As we have discussed at several points, when we are actually interested in partial effects on response probabilities, estimating coefficients up to an unknown scale is usually unsatisfying. Recently, Blundell and Powell (2004) have shown how to estimate an average structural function-and therefore the APEs-when $y_{2}$ is a continuous endogenous explanatory variable. We can think of their method as a nonparametric extension of the model in Section 15.7.2.

Blundell and Powell allow complete flexibility of functional forms, but we can illustrate their general approach using an index model $y_{1}=1\left[\mathbf{x}_{1} \boldsymbol{\beta}_{1}+u_{1}>0\right]$, where $\mathbf{x}_{1}$ can be any function of $\left(\mathbf{z}_{1}, y_{2}\right)$. The key assumption is that $y_{2}$ can be written as $y_{2}=g_{2}(\mathbf{z})+v_{2}$, where ( $u_{1}, v_{2}$ ) is independent of $\mathbf{z}$. The independence of the additive error $v_{2}$ and $\mathbf{z}$ pretty much rules out discreteness in $y_{2}$, even though $g_{2}(\cdot)$ can be left unspecified. Under the independence assumption,\\
$\mathrm{P}\left(y_{1}=1 \mid \mathbf{z}, v_{2}\right)=\mathrm{E}\left(y_{1} \mid \mathbf{z}, v_{2}\right)=H\left(\mathbf{x}_{1} \boldsymbol{\beta}_{1}, v_{2}\right)$\\
for some (generally unknown) function $H(\cdot, \cdot)$. The average structural function is just $\operatorname{ASF}\left(\mathbf{z}_{1}, y_{2}\right)=\mathrm{E}_{v_{i 2}}\left[H\left(\mathbf{x}_{1} \boldsymbol{\beta}_{1}, v_{i 2}\right)\right]$. We can estimate $H$ and $\boldsymbol{\beta}_{1}$ quite generally by first\\
estimating the function $g_{2}(\cdot)$ and then obtaining residuals $\hat{v}_{i 2}=y_{i 2}-\hat{g}_{2}\left(\mathbf{z}_{i}\right)$. Then, $H$ and $\boldsymbol{\beta}_{1}$ can be estimated in a second step by a semiparametric procedure-that is, one that does not assume $H$ is in a parametric family. Then the ASF is estimated by averaging out the reduced form residuals,


\begin{equation*}
\widehat{\operatorname{ASF}}\left(\mathbf{z}_{1}, y_{2}\right)=N^{-1} \sum_{i=1}^{N} \hat{H}\left(\mathbf{x}_{1} \hat{\boldsymbol{\beta}}_{1}, \hat{v}_{i 2}\right) ; \tag{15.61}
\end{equation*}


derivatives and changes can be computed with respect to elements of $\left(\mathbf{z}_{1}, y_{2}\right)$.\\
Blundell and Powell (2004) actually allow $P\left(y_{1}=1 \mid \mathbf{z}, y_{2}\right)$ to have the general form $H\left(\mathbf{z}_{1}, y_{2}, v_{2}\right)$, and then the second-step estimation is entirely nonparametric. They also allow $\hat{g}_{2}(\cdot)$ to be fully nonparametric. But parametric approximations in each stage might produce good estimates of the APEs. For example, $y_{i 2}$ can be regressed on flexible functions of $\mathbf{z}_{i}$ to obtain $\hat{v}_{i 2}$. Then, one can estimate probit or logit models in the second stage that include functions of $\mathbf{z}_{1}, y_{2}$, and $\hat{v}_{2}$ in a flexible way-for example, with levels, quadratics, interactions, and maybe even higherorder polynomials of each. Then, one simply averages out $\hat{v}_{i 2}$, as in equation (15.61). Valid standard errors and test statistics can be obtained by bootstrapping or by using the delta method.

Obtaining flexible methods that allow for general discrete $y_{2}$, which would extend the parametric approach of Section 15.7.3, is an important task for future research.

\subsection*{15.8 Binary Response Models for Panel Data}
When analyzing binary responses in the context of panel data, it is often useful to begin with a linear model with an additive, unobserved effect, and then, just as in Chapters 10 and 11, use the within transformation or first differencing to remove the unobserved effect. A linear probability model for binary outcomes has the same problems as in the cross section case. In fact, it is probably less appealing for unobserved effects models, as it implies the unnatural restrictions $\mathbf{x}_{i t} \boldsymbol{\beta} \leq c_{i} \leq 1-\mathbf{x}_{i t} \boldsymbol{\beta}, t=1, \ldots, T$, on the unobserved effects. To see this, note that the response probability for an unobserved effects LPM is $\mathrm{P}\left(y_{i t}=1 \mid \mathbf{x}_{i}, c_{i}\right)=\mathrm{P}\left(y_{i t}=1 \mid \mathbf{x}_{i t}, c_{i}\right)= \mathbf{x}_{i t} \boldsymbol{\beta}+c_{i}$. As in the pure cross section case, the linear functional form is almost certainly false, except in special cases. But FE or FD estimation of the LPM might provide reasonable estimates of APEs. Plus, the LPM has the advantage of not requiring a distributional assumption on $\mathrm{D}\left(c_{i} \mid \mathbf{x}_{i}\right)$ (which means we do not need to be concerned that $c_{i}$ is bounded by linear combinations of $\mathbf{x}_{i t}$ ); nor do we have to assume independence of the responses $\left\{y_{i 1}, \ldots, y_{i T}\right\}$ conditional on $\left(\mathbf{x}_{i}, c_{i}\right)$ (provided\\
we make our inference robust to serial dependence, as well as heteroskedasticity). The nonlinear methods we cover in this section require one or both of these assumptions. First, we briefly cover estimation when the model is not specified with unobserved heterogeneity and has explanatory variables that may not be strictly exogenous.

\subsection*{15.8.1 Pooled Probit and Logit}
In Section 13.8 we used a probit model to illustrate partial likelihood methods with panel data. Naturally, we can use logit or any other binary response function as well. Suppose the model is


\begin{equation*}
\mathrm{P}\left(y_{i t}=1 \mid \mathbf{x}_{i t}\right)=G\left(\mathbf{x}_{i t} \boldsymbol{\beta}\right), \quad t=1,2, \ldots, T \tag{15.62}
\end{equation*}


where $G(\cdot)$ is a known function taking on values in the open unit interval. As we discussed in Chapter 13, $\mathbf{x}_{i t}$ can contain a variety of factors, including time dummies, interactions of time dummies with time-constant or time-varying variables, and lagged dependent variables.

In specifying the model (15.62) we have not assumed nearly enough to obtain the distribution of $\mathbf{y}_{i} \equiv\left(y_{i 1}, \ldots, y_{i T}\right)$ given $\mathbf{x}_{i}=\left(\mathbf{x}_{i 1}, \ldots, \mathbf{x}_{i T}\right)$, for two reasons. First, we have not assumed $\mathrm{D}\left(y_{i t} \mid \mathbf{x}_{i 1}, \ldots, \mathbf{x}_{i T}\right)=\mathrm{D}\left(y_{i t} \mid \mathbf{x}_{i t}\right)$, so that $\left\{\mathbf{x}_{i t}: t=1, \ldots, T\right\}$ is not necessarily strictly exogenous. Second, even if we assume strict exogeneity, we have not restricted the dependence in $\left\{y_{i t}: t=1, \ldots, T\right\}$ conditional on $\mathbf{x}_{i}$. As with many pooled MLE problems, we have only specified a model for $\mathrm{D}\left(y_{i t} \mid \mathbf{x}_{i t}\right)$. If that model is correctly specified, we can obtain a $\sqrt{N}$-consistent, asymptotically normal estimator by maximizing the partial (pooled) log-likelihood function. For binary response, we maximize the partial log likelihood\\
$\sum_{i=1}^{N} \sum_{t=1}^{T}\left\{y_{i t} \log G\left(\mathbf{x}_{i t} \boldsymbol{\beta}\right)+\left(1-y_{i t}\right) \log \left[1-G\left(\mathbf{x}_{i t} \boldsymbol{\beta}\right)\right]\right\}$,\\
which is simply an exercise in pooled estimation. Without further assumptions, a robust variance matrix estimator is needed to account for serial correlation in the scores across $t$; see equation (13.53) with $\hat{\boldsymbol{\beta}}$ in place of $\hat{\boldsymbol{\theta}}$ and $G$ in place of $\Phi$. Wald and score statistics can be computed as in Chapter 12.

In the case that the model (15.62) is dynamically complete, that is,


\begin{equation*}
\mathbf{P}\left(y_{i t}=1 \mid \mathbf{x}_{i t}, y_{i, t-1}, \mathbf{x}_{i, t-1}, \ldots\right)=\mathbf{P}\left(y_{i t}=1 \mid \mathbf{x}_{i t}\right), \tag{15.63}
\end{equation*}


or\\
$\mathrm{D}\left(y_{i t} \mid \mathbf{x}_{i t}, y_{i, t-1}, \mathbf{x}_{i, t-1}, \ldots\right)=\mathrm{D}\left(y_{i t} \mid \mathbf{x}_{i t}\right)$,\\
inference is considerably easier: all the usual statistics from a probit or logit that pools observations and treats the sample as a long independent cross section of size $N T$ are valid, including likelihood ratio statistics. Remember, we are definitely not assuming independence across $t$ (for example, $\mathbf{x}_{i t}$ can contain lagged dependent variables). Dynamic completeness implies that the scores are serially uncorrelated across $t$, which is the key condition for the standard inference procedures to be valid. (See the general treatment in Section 13.8.)

To test for dynamic completeness, we can always add a lagged dependent variable and possibly lagged explanatory variables. As an alternative, we can derive a simple one-degree-of-freedom test that works regardless of what is in $\mathbf{x}_{i t}$. For concreteness, we focus on the probit case; other index models are handled similarly. Define $u_{i t} \equiv y_{i t}-\Phi\left(\mathbf{x}_{i t} \boldsymbol{\beta}\right)$, so that, under assumption (15.63), $\mathrm{E}\left(u_{i t} \mid \mathbf{x}_{i t}, y_{i, t-1}, \mathbf{x}_{i, t-1}, \ldots\right)=0$, all $t$. It follows that $u_{i t}$ is uncorrelated with any function of the variables $\left(\mathbf{x}_{i t}, y_{i, t-1}\right.$, $\mathbf{x}_{i, t-1}, \ldots$ ), including $u_{i, t-1}$. By studying equation (13.53), we can see that it is serial correlation in the $u_{i t}$ that makes the usual inference procedures invalid. Let $\hat{u}_{i t}= y_{i t}-\Phi\left(\mathbf{x}_{i t} \hat{\boldsymbol{\beta}}\right)$. Then a simple test is available by using pooled probit to estimate the artificial model


\begin{equation*}
" \mathrm{P}\left(y_{i t}=1 \mid \mathbf{x}_{i t}, \hat{u}_{i, t-1}\right)=\Phi\left(\mathbf{x}_{i t} \boldsymbol{\beta}+\gamma_{1} \hat{u}_{i, t-1}\right) " \tag{15.64}
\end{equation*}


using time periods $t=2, \ldots, T$. The null hypothesis is $\mathrm{H}_{0}: \gamma_{1}=0$. If $\mathrm{H}_{0}$ is rejected, then so is assumption (15.63). This is a case where under the null hypothesis, the estimation of $\boldsymbol{\beta}$ required to obtain $\hat{u}_{i, t-1}$ does not affect the limiting distribution of any of the usual test statistics, Wald, $L R$, or $L M$, of $\mathrm{H}_{0}: \gamma_{1}=0$. The Wald statistic, that is, the $t$ statistic on $\hat{\gamma}_{1}$, is the easiest to obtain. For the $L M$ and $L R$ statistics we must be sure to drop the first time period in estimating the restricted model ( $\gamma_{1}=0$ ).

\subsection*{15.8.2 Unobserved Effects Probit Models under Strict Exogeneity}
A popular model for binary outcomes with panel data is the unobserved effects probit model. A key assumption for many of the estimators for unobserved effects probit models, as well as other nonlinear panel data models, is that the observed covariates are strictly exogenous conditional on the unobserved effect. Now, we need this assumption in terms of conditional distributions, which we can state generically as


\begin{equation*}
\mathrm{D}\left(y_{i t} \mid \mathbf{x}_{i}, c_{i}\right) \equiv \mathrm{D}\left(y_{i t} \mid \mathbf{x}_{i 1}, \mathbf{x}_{i 2}, \ldots, \mathbf{x}_{i T}, c_{i}\right)=\mathrm{D}\left(y_{i t} \mid \mathbf{x}_{i t}, c_{i}\right), \quad t=1, \ldots, T \tag{15.65}
\end{equation*}


where $c_{i}$ is the unobserved effect and $\mathbf{x}_{i} \equiv\left(\mathbf{x}_{i 1}, \mathbf{x}_{i 2}, \ldots, \mathbf{x}_{i T}\right)$. In stating this assumption, we are not restricting in any way the joint distribution conditional on $\left(\mathbf{x}_{i}, c_{i}\right)$, which we write as $\mathrm{D}\left(\mathbf{y}_{i} \mid \mathbf{x}_{i}, c_{i}\right)$ for $\mathbf{y}_{i}=\left(y_{i 1}, \ldots, y_{i T}\right)$. As in previous chapters, where we stated the strict exogeneity assumption conditional on $c_{i}$ in terms of the condi-\\
tional expectation, assumption (15.65) rules out models with lagged dependent variables in $\mathbf{x}_{i t}$ and also other situations where one or more elements of $\mathbf{x}_{i t}$ may react in the future to idiosyncratic changes in $y_{i t}$. Assumption (15.65) also requires that $\mathbf{x}_{i t}$ includes enough lags of underlying explanatory variables if distributed lag dynamics are present.

For the unobserved effects probit model, we specify a response probability that fully determines the conditional distribution $\mathbf{D}\left(y_{i t} \mid \mathbf{x}_{i t}, c_{i}\right)$ :


\begin{equation*}
\mathrm{P}\left(y_{i t}=1 \mid \mathbf{x}_{i t}, c_{i}\right)=\Phi\left(\mathbf{x}_{i t} \boldsymbol{\beta}+c_{i}\right), \quad t=1, \ldots, T . \tag{15.66}
\end{equation*}


Most analyses are not convincing unless $\mathbf{x}_{i t}$ contains a full set of time dummies. In what follows, we usually leave time dummies implicit to simplify the notation.

Our interest is in the response probability (15.66), and we would like to be able to estimate the parameters and partial effects just by specifying equation (15.66). But we already saw in Chapters 10 and 11 that just specifying a model with contemporaneous conditioning variables is not sufficient for estimating the parameters, and the same is true here: without more assumptions, $\boldsymbol{\beta}$ is not identified. Even adding the strict exogeneity assumption (15.65) is not enough. Unlike with linear models, we face a further difficulty, which has implications both for estimating $\boldsymbol{\beta}$ and for estimating partial effects of interest. Namely, we must specify how $c_{i}$ relates to the covariates.

What happens if we try to proceed without placing restrictions on $\mathrm{D}\left(c_{i} \mid \mathbf{x}_{i}\right)$ ? One possibility is to add, in addition to assumptions (15.65) and (15.66), the assumption that the responses are independent conditional on $\left(\mathbf{x}_{i}, c_{i}\right)$ :\\
$y_{i 1}, \ldots, y_{i T}$ are independent conditional on $\left(\mathbf{x}_{i}, c_{i}\right)$.

Because of the presence of $c_{i}$, the $y_{i t}$ are dependent across $t$ conditional only on the observables, $\mathbf{x}_{i}$. (Assumption (15.67) is analogous to the linear model assumption that, conditional on $\left(\mathbf{x}_{i}, c_{i}\right)$, the $y_{i t}$ are serially uncorrelated; see Assumption FE. 3 in Chapter 10.) Under assumptions (15.65), (15.66), and (15.67), we can derive the density of $\left(y_{i 1}, \ldots, y_{i T}\right)$ conditional on $\left(\mathbf{x}_{i}, c_{i}\right)$ :


\begin{equation*}
f\left(y_{1}, \ldots, y_{T} \mid \mathbf{x}_{i}, c_{i} ; \boldsymbol{\beta}\right)=\prod_{t=1}^{T} f\left(y_{t} \mid \mathbf{x}_{i t}, c_{i} ; \boldsymbol{\beta}\right), \tag{15.68}
\end{equation*}


where $f\left(y_{t} \mid \mathbf{x}_{t}, c ; \boldsymbol{\beta}\right)=\Phi\left(\mathbf{x}_{t} \boldsymbol{\beta}+c\right)^{y_{t}}\left[1-\Phi\left(\mathbf{x}_{t} \boldsymbol{\beta}+c\right)\right]^{1-y_{t}}$. Ideally, we could estimate the quantities of interest without restricting the relationship between $c_{i}$ and the $\mathbf{x}_{i t}$. In this spirit, we might view the $c_{i}$ as parameters to be estimated along with $\boldsymbol{\beta}$, as this treatment obviates the need to make assumptions about the distribution of $c_{i}$ given\\
$\mathbf{x}_{i}$. The log-likelihood function is $\sum_{i=1}^{N} \ell_{i}\left(c_{i}, \boldsymbol{\beta}\right)$, where $\ell_{i}\left(c_{i}, \boldsymbol{\beta}\right)$ is the log of equation (15.68) evaluated at the $y_{i t}$. Unfortunately, in addition to being computationally difficult, estimation of the $c_{i}$ along with $\boldsymbol{\beta}$ introduces an incidental parameters problem. Unlike in the linear case, where estimating the $c_{i}$ along with $\boldsymbol{\beta}$ leads to the $\sqrt{N}$ consistent FE estimator of $\boldsymbol{\beta}$, in the present case estimating the $c_{i}$ ( $N$ of them) along with $\boldsymbol{\beta}$ leads to inconsistent estimation of $\boldsymbol{\beta}$ with $T$ fixed and $N \rightarrow \infty$. We discuss the incidental parameters problem in more detail for the unobserved effects logit model in Section 15.8.3.

The estimator of $\boldsymbol{\beta}$ obtained by treating the $c_{i}$ as parameters has been called the "fixed effects probit" estimator, an unfortunate name. As we saw with linear models, and as we will see with the logit model in the next subsection and with count data models in Chapter 18, in some cases we can consistently estimate the parameters $\boldsymbol{\beta}$ without specifying a distribution for $c_{i}$ given $\mathbf{x}_{i}$. This feature is the hallmark of an FE analysis for most microeconometric applications. By contrast, treating the $c_{i}$ as parameters to estimate can lead to potentially serious biases, as in the probit case.

Here we follow the same approach adopted for linear models: we always treat $c_{i}$ as an unobservable random variable drawn along with $\left(\mathbf{x}_{i}, \mathbf{y}_{i}\right)$. The question is, under what additional assumptions can we consistently estimate parameters, as well as interesting partial effects? Unfortunately, for the unobserved effects probit model, we must make an assumption about the relationship between $c_{i}$ and $\mathbf{x}_{i}$. The traditional random effects probit model adds, to assumptions (15.65), (15.66), and (15.67), the assumption


\begin{equation*}
c_{i} \mid \mathbf{x}_{i} \sim \operatorname{Normal}\left(0, \sigma_{c}^{2}\right) \tag{15.69}
\end{equation*}


This is a strong assumption, as it implies that $c_{i}$ and $\mathbf{x}_{i}$ are independent and that $c_{i}$ has a normal distribution. It is not enough to assume that $c_{i}$ and $\mathbf{x}_{i}$ are uncorrelated, or even that $\mathrm{E}\left(c_{i} \mid \mathbf{x}_{i}\right)=0$. The assumption $\mathrm{E}\left(c_{i}\right)=0$ is without loss of generality provided $\mathbf{x}_{i t}$ contains an intercept, as it always should.

Before we discuss estimation of the random effects probit model, we should be sure we know what we want to estimate. As in Section 15.7.1, consistent estimation of $\boldsymbol{\beta}$ means that we can consistently estimate the partial effects of the elements of $\mathbf{x}_{t}$ on the response probability $\mathrm{P}\left(y_{t}=1 \mid \mathbf{x}_{t}, c\right)$ at the average value of $c$ in the population, $c=0$. (We can also estimate the relative effects of any two elements of $\mathbf{x}_{t}$ for any value of $c$, as the relative effects do not depend on $c$.) For the reasons discussed in Section 15.7.1, APEs are at least as useful. Since $c_{i} \sim \operatorname{Normal}\left(0, \sigma_{c}^{2}\right)$, the APE for a continuous $x_{t j}$ is $\left[\beta_{j} /\left(1+\sigma_{c}^{2}\right)^{1 / 2}\right] \phi\left[\mathbf{x}_{t} \boldsymbol{\beta} /\left(1+\sigma_{c}^{2}\right)^{1 / 2}\right]$, just as in equation (15.38). Therefore, we only need to estimate $\boldsymbol{\beta}_{c} \equiv \boldsymbol{\beta} /\left(1+\sigma_{c}^{2}\right)^{1 / 2}$ to estimate the APEs, for\\
either continuous or discrete explanatory variables. (In other branches of applied statistics, such as biostatistics and education, the coefficients indexing the APEs- $\boldsymbol{\beta}_{c}$ in our notation-are called the population-averaged parameters.)

Under assumptions (15.65), (15.66), (15.67), and (15.69), a conditional maximum likelihood approach is available for estimating $\boldsymbol{\beta}$ and $\sigma_{c}^{2}$. This is a special case of the approach in Section 13.9. Because the $c_{i}$ are not observed, they cannot appear in the likelihood function. Instead, we find the joint distribution of $\left(y_{i 1}, \ldots, y_{i T}\right)$ conditional on $\mathbf{x}_{i}$, a step that requires us to integrate out $c_{i}$. Since $c_{i}$ has a $\operatorname{Normal}\left(0, \sigma_{c}^{2}\right)$ distribution,


\begin{equation*}
f\left(y_{1}, \ldots, y_{T} \mid \mathbf{x}_{i} ; \boldsymbol{\theta}\right)=\int_{-\infty}^{\infty}\left[\prod_{t=1}^{T} f\left(y_{t} \mid \mathbf{x}_{i t}, c ; \boldsymbol{\beta}\right)\right]\left(1 / \sigma_{c}\right) \phi\left(c / \sigma_{c}\right) \mathrm{d} c, \tag{15.70}
\end{equation*}


where $f\left(y_{t} \mid \mathbf{x}_{t}, c ; \boldsymbol{\beta}\right)=\Phi\left(\mathbf{x}_{t} \boldsymbol{\beta}+c\right)^{y_{t}}\left[1-\Phi\left(\mathbf{x}_{t} \boldsymbol{\beta}+c\right)\right]^{1-y_{t}}$ and $\boldsymbol{\theta}$ contains $\boldsymbol{\beta}$ and $\sigma_{c}^{2}$. Plugging in $y_{i t}$ for all $t$ and taking the log of equation (15.70) gives the conditional $\log$ likelihood $\ell_{i}(\boldsymbol{\theta})$ for each $i$. The log-likelihood function for the entire sample of size $N$ can be maximized with respect to $\boldsymbol{\beta}$ and $\sigma_{c}^{2}$ (or $\boldsymbol{\beta}$ and $\sigma_{c}$ ) to obtain $\sqrt{N}$-consistent asymptotically normal estimators; Butler and Moffitt (1982) describe a procedure for approximating the integral in equation (15.70). The conditional MLE in this context is typically called the random effects probit estimator, and the theory in Section 13.9 can be applied directly to obtain asymptotic standard errors and test statistics. Since $\boldsymbol{\beta}$ and $\sigma_{c}^{2}$ can be estimated, the partial effects at $c=0$ as well as the APEs can be estimated. Since the variance of the idiosyncratic error in the latent variable model is unity, the relative importance of the unobserved effect is measured as $\rho= \sigma_{c}^{2} /\left(\sigma_{c}^{2}+1\right)$, which is also the correlation between the composite latent error, say, $c_{i}+e_{i t}$, across any two time periods. Many random effects probit routines report $\hat{\rho}$ and its standard error; these statistics lead to an easy test for the presence of the unobserved effect.

Assumptions (15.67) and (15.69) are fairly strong, and it is possible to relax them. First consider relaxing assumption (15.67). One useful observation is that, under assumptions (15.66) and (15.69) only (even without assumption (15.65)),


\begin{equation*}
\mathbf{P}\left(y_{i t}=1 \mid \mathbf{x}_{i t}\right)=\Phi\left(\mathbf{x}_{i t} \boldsymbol{\beta}_{c}\right), \tag{15.71}
\end{equation*}


where $\boldsymbol{\beta}_{c}=\boldsymbol{\beta} /\left(1+\sigma_{c}^{2}\right)^{1 / 2}$. Therefore, just as in Section 15.8.1, we can estimate $\boldsymbol{\beta}_{c}$ from pooled probit of $y_{i t}$ on $\mathbf{x}_{i t}, t=1, \ldots, T, i=1, \ldots, N$, meaning that we directly estimate the APEs. If $c_{i}$ is truly present, $\left\{y_{i t}: t=1, \ldots, T\right\}$ will not be independent conditional on $\mathbf{x}_{i}$. Robust inference is needed to account for the serial dependence, as discussed in Section 15.8.1.

Rather than simply calculating robust standard errors for $\hat{\boldsymbol{\beta}}_{c}$ after pooled probit estimation, or using the full random effects (RE) assumptions and obtaining the MLE, we can apply the generalized estimating equation (GEE) approach that we introduced in Section 12.9. As we discussed there, GEE is simply multivariate weighted nonlinear least squares with a correctly specified model for $\mathrm{E}\left(\mathbf{y}_{i} \mid \mathbf{x}_{i}\right)$ but a generally misspecified model for the conditional variance matrix $\operatorname{Var}\left(\mathbf{y}_{i} \mid \mathbf{x}_{i}\right)$, where $\mathbf{y}_{i}$ is the $T$-vector of responses for unit $i$. In the present application, we have correctly specified each element of $\mathrm{E}\left(\mathbf{y}_{i} \mid \mathbf{x}_{i}\right)$ as $\mathrm{E}\left(y_{i t} \mid \mathbf{x}_{i}\right)=\Phi\left(\mathbf{x}_{i t} \boldsymbol{\beta}_{c}\right)$. Plus, each conditional variance, $\operatorname{Var}\left(y_{i t} \mid \mathbf{x}_{i}\right)$, is necessarily equal to $\operatorname{Var}\left(y_{i t} \mid \mathbf{x}_{i}\right)=\Phi\left(\mathbf{x}_{i t} \boldsymbol{\beta}_{c}\right)\left[1-\Phi\left(\mathbf{x}_{i t} \boldsymbol{\beta}_{c}\right)\right]$. What is difficult to obtain, even under the full set of RE probit assumptions, are the conditional correlations, $\operatorname{Corr}\left(y_{i t}, y_{i s} \mid \mathbf{x}_{i}\right)$. (In fact, there are no closed-form expressions for these correlations.) The most straightforward GEE method is to specify a constant and exchangeable "working" correlation matrix, which means in using WMNLS we act as if $\operatorname{Corr}\left(y_{i t}, y_{i s} \mid \mathbf{x}_{i}\right)=\rho$ for $|\rho|<1$. This assumption is almost certainly false; the actual conditional correlations likely depend on $\mathbf{x}_{i}$ in a complicated fashion. Nevertheless, as we explained in Section 12.9, the hope is that allowing a nonzero correlation in the working variance-covariance matrix will produce an estimator asymptotically more efficient than just using pooled probit. At the same time, such a WMNLS estimator will be consistent under assumptions (15.65), (15.66), and (15.69), the same assumptions we use for consistency of pooled probit. Practically, we can think of GEE for RE probit models as a method that has the same robustness as pooled probit yet, ideally, gets back some of the efficiency lost by ignoring the serial dependence in estimation.

The terminology used in the GEE literature when applied to nonlinear unobserved effects models, such as RE probit, can be elusive. We can minimize our confusion by keeping straight the distinctions among a model, the quantities of interest, and an estimation method. The response probability conditional only on $\mathbf{x}_{i t}, \mathrm{P}\left(y_{i t}=1 \mid \mathbf{x}_{i t}\right) =\Phi\left(\mathbf{x}_{i t} \boldsymbol{\beta}_{c}\right)$-for which we only need assumptions (15.66) and (15.69) to derive-is called the population-averaged (PA) model in the GEE literature because the parameters, $\boldsymbol{\beta}_{c}$, index the population-averaged effects. Therefore, using the GEE terminology, the WMNLS estimator is designed to estimate the parameters in the PA model, just as is pooled probit. The model under the full set of RE probit assumptions is called the subject-specific (SS) model, primarily because the response probability in equation (15.66), which is conditional on $\mathbf{x}_{i t}$ and $c_{i}$, allows us, in principle, to compute partial effects for different individuals as described by the heterogeneity, $c_{i}$. In actuality, because $c_{i}$ is not observed (and cannot be estimated with small $T$ ), all we can really do is estimate the partial effects on $\mathrm{P}\left(y_{i t}=1 \mid \mathbf{x}_{i t}=\mathbf{x}_{t}, c_{i}=c\right)$ for different values of $c$. These values include the mean (median) value, $c=0$, but also percentiles\\
in the $\operatorname{Normal}\left(0, \sigma_{c}^{2}\right)$ distribution (because we can estimate $\sigma_{c}^{2}$ along with $\boldsymbol{\beta}$ ), but the implication that we can estimate partial effects for specific individuals under the full RE probit assumptions is somewhat misleading.

Rather than act as if there are different models, in the context of unobserved effects probit models it seems better to organize the discussion around quantities that can be estimated under different assumptions on a common model. Much of the time it makes sense to start with equation (15.66) as the model of interest. Then we specify partial effects at different values of $c$ or APEs as the quantities of interest. Finally, we recognize that under (15.66) and (15.69) only, we can estimate the APEs (or PA effects) by pooled probit. If we add the strict exogeneity assumption (15.65), then we can use GEE. If we use the full set of RE probit assumptions and apply MLE, then we can separately estimate $\boldsymbol{\beta}$ and $\sigma_{c}^{2}$, and therefore the APEs and partial effects at interesting values of $c$.

A different way to relax assumption (15.67) is to assume a particular correlation structure and then use full CMLE. For example, for each $t$ write the latent variable model as


\begin{equation*}
y_{i t}^{*}=\mathbf{x}_{i t} \boldsymbol{\beta}+c_{i}+e_{i t}, \quad y_{i t}=1\left[y_{i t}^{*}>0\right] \tag{15.72}
\end{equation*}


and assume that the $T \times 1$ vector $\mathbf{e}_{i}$ is multivariate normal, with unit variances, but unrestricted correlation matrix. This assumption, along with assumptions (15.65), (15.66), and (15.69), fully characterizes the distribution of $\mathbf{y}_{i}$ given $\mathbf{x}_{i}$. However, even for moderate $T$, computation of the CMLE can be very difficult. Recent advances in simulation methods of estimation make it possible to estimate such models for fairly large $T$; see, for example, Keane (1993) and Geweke and Keane (2001). The pooled probit and GEE procedures that we have described are valid for estimating $\boldsymbol{\beta}_{c}$, the vector that indexes APEs, regardless of the serial dependence in $\left\{e_{i t}\right\}$, but they are inefficient relative to the full CMLE.

As in the linear case, in many applications the point of introducing the unobserved effect, $c_{i}$, is to explicitly allow unobservables to be correlated with some elements of $\mathbf{x}_{i t}$. Chamberlain (1980) allowed for correlation between $c_{i}$ and $\mathbf{x}_{i}$ by assuming a conditional normal distribution with linear expectation and constant variance. A Mundlak (1978) version of Chamberlain's assumption is


\begin{equation*}
c_{i} \mid \mathbf{x}_{i} \sim \operatorname{Normal}\left(\psi+\overline{\mathbf{x}}_{i} \boldsymbol{\xi}, \sigma_{a}^{2}\right), \tag{15.73}
\end{equation*}


where $\overline{\mathbf{x}}_{i}$ is the average of $\mathbf{x}_{i t}, t=1, \ldots, T$ and $\sigma_{a}^{2}$ is the variance of $a_{i}$ in the equation $c_{i}=\psi+\overline{\mathbf{x}}_{i} \boldsymbol{\xi}+a_{i}$. (In other words, $\sigma_{a}^{2}$ is the conditional variance of $c_{i}$, which is assumed not to depend on $\mathbf{x}_{i}$.) Chamberlain (1980) allowed more generality by having $\mathbf{x}_{i}$, the vector of all explanatory variables across all time periods, in place of\\
$\overline{\mathbf{x}}_{i}$. We will work with assumption (15.73), as it conserves on parameters; the more general model requires a simple notational change. Chamberlain (1980) called model (15.66) under assumptions (15.65) and (15.73) a random effects probit model, so we refer to the model as Chamberlain's correlated random effects probit model. While assumption (15.73) is restrictive in that it specifies a distribution for $c_{i}$ given $\mathbf{x}_{i}$, it at least allows for some dependence between $c_{i}$ and $\mathbf{x}_{i}$.

As in the linear case, we can only estimate the effects of time-varying elements in $\mathbf{x}_{i t}$. In particular, $\mathbf{x}_{i t}$ should no longer contain a constant, as that would be indistinguishable from $\psi$ in assumption (15.73). If our original model contains a time-constant explanatory variable, say $w_{i}$, it can be included among the explanatory variables, but we cannot distinguish its effect from $c_{i}$ unless we assume that the coefficient for $w_{i}$ in $\boldsymbol{\xi}$ is zero. (That is, unless we assume that $c_{i}$ is partially uncorrelated with $w_{i}$.) Time dummies, which do not vary across $i$, are omitted from $\overline{\mathbf{x}}_{i}$.

If assumptions (15.65), (15.66), (15.67), and (15.73) hold, estimation of $\boldsymbol{\beta}, \boldsymbol{\psi}, \boldsymbol{\xi}$, and $\sigma_{a}^{2}$ is straightforward because we can write the latent variable as $y_{i t}^{*}=\psi+\mathbf{x}_{i t} \boldsymbol{\beta}+ \overline{\mathbf{x}}_{i} \boldsymbol{\xi}+a_{i}+e_{i t}$, where the $e_{i t}$ are independent $\operatorname{Normal}(0,1)$ variates (conditional on $\left(\mathbf{x}_{i}, a_{i}\right)$ ), and $a_{i} \mid \mathbf{x}_{i} \sim \operatorname{Normal}\left(0, \sigma_{a}^{2}\right)$. In other words, by adding $\overline{\mathbf{x}}_{i}$ to the equation for each time period, we arrive at a traditional RE probit model. (The variance we estimate is $\sigma_{a}^{2}$ rather than $\sigma_{c}^{2}$, but, as we will see, this suits our purposes nicely.) Adding $\overline{\mathbf{x}}_{i}$ as a set of controls for unobserved heterogeneity is very intuitive: we are estimating the effect of changing $x_{i t j}$ but holding the time average fixed. A test of the usual RE probit model is easily obtained as a test of $\mathrm{H}_{0}: \boldsymbol{\xi}=\mathbf{0}$. Estimation can be carried out using standard RE probit software. Given estimates of $\psi$ and $\boldsymbol{\xi}$, we can estimate $\mathrm{E}\left(c_{i}\right)=\psi+\mathrm{E}\left(\overline{\mathbf{x}}_{i}\right) \boldsymbol{\xi}$ by $\hat{\mu}_{c} \equiv \hat{\psi}+\overline{\mathbf{x}} \hat{\boldsymbol{\xi}}$, where $\overline{\mathbf{x}}$ is the sample average of $\overline{\mathbf{x}}_{i}$. Therefore, for any vector $\mathbf{x}_{t}$, we can estimate the response probability at $\mathrm{E}\left(c_{i}\right)$ as $\Phi\left(\hat{\mu}_{c}+\mathbf{x}_{t} \hat{\boldsymbol{\beta}}\right)$. Taking differences or derivatives (with respect to the elements of $\mathbf{x}_{t}$ ) allows us to estimate the partial effects on the response probabilities for any value of $\mathbf{x}_{t}$. Further, we can estimate $\sigma_{c}^{2}$ by using the relationship $\sigma_{c}^{2}=\xi^{\prime} \operatorname{Var}\left(\overline{\mathbf{x}}_{i}\right) \boldsymbol{\xi}+\sigma_{a}^{2}$, and so $\hat{\sigma}_{c}^{2} \equiv \hat{\boldsymbol{\xi}}^{\prime}\left[N^{-1} \sum_{i=1}^{N}\left(\overline{\mathbf{x}}_{i}-\overline{\mathbf{x}}\right)^{\prime}\left(\overline{\mathbf{x}}_{i}-\overline{\mathbf{x}}\right)\right] \hat{\boldsymbol{\xi}}+\hat{\sigma}_{a}^{2}$ is consistent for $\sigma_{c}^{2}$ as $N \rightarrow \infty$. Therefore, we can plug in values of $c$ that are a certain number of estimated standard deviations from $\hat{\mu}_{c}$, say $\hat{\mu}_{c} \pm \hat{\sigma}_{c}$. But there is a subtle point at work here: if $\boldsymbol{\xi} \neq \mathbf{0}, c_{i}$ does not generally have an unconditional normal distribution unless $\overline{\mathbf{x}}_{i} \boldsymbol{\xi}$ is normally distributed; sufficient is that $\overline{\mathbf{x}}_{i}$ is multivariate normal. With large $T$ and weak dependence in the time series dimension, it might be that $\overline{\mathbf{x}}_{i} \boldsymbol{\xi}$ has an approximate normal distribution (because it is an average across $t$ ). But weak dependence is not necessarily a good assumption, and even if it were, if we had a large $T$ we might proceed by treating the $c_{i}$ as parameters to estimate. Generally, we should not expect the distribution of $c_{i}$ to be normal.

If we drop assumption (15.67), we can still estimate scaled versions of $\boldsymbol{\psi}, \boldsymbol{\beta}$, and $\boldsymbol{\xi}$. Under assumptions (15.65), (15.66), and (15.73) we have


\begin{align*}
\mathrm{P}\left(y_{i t}=1 \mid \mathbf{x}_{i}\right) & =\Phi\left[\left(\psi+\mathbf{x}_{i t} \boldsymbol{\beta}+\overline{\mathbf{x}}_{i} \boldsymbol{\xi}\right) \cdot\left(1+\sigma_{a}^{2}\right)^{-1 / 2}\right] \\
& \equiv \Phi\left(\psi_{a}+\mathbf{x}_{i t} \boldsymbol{\beta}_{a}+\overline{\mathbf{x}}_{i} \boldsymbol{\xi}_{a}\right) \tag{15.74}
\end{align*}


where the $a$ subscript means that a parameter vector has been multiplied by $\left(1+\sigma_{a}^{2}\right)^{-1 / 2}$. It follows immediately that $\psi_{a}, \boldsymbol{\beta}_{a}$, and $\boldsymbol{\xi}_{a}$ can be consistently estimated using a pooled probit analysis of $y_{i t}$ on $1, \mathbf{x}_{i t}, \overline{\mathbf{x}}_{i}, t=1, \ldots, T, i=1, \ldots, N$. Because the $y_{i t}$ will be dependent condition on $\mathbf{x}_{i}$, inference that is robust to arbitrary time dependence is required. We can also use GEE.

Conveniently, once we have estimated $\psi_{a}, \boldsymbol{\beta}_{a}$, and $\boldsymbol{\xi}_{a}$, we can estimate the APEs. (We could apply the results from Section 2.2.5 here, but a direct argument is instructive.) To see how, we need to average $\mathrm{P}\left(y_{t}=1 \mid \mathbf{x}_{t}=\mathbf{x}^{\mathrm{o}}, c_{i}\right)$ across the distribution of $c_{i}$; that is, we need to find $\mathrm{E}\left[\mathrm{P}\left(y_{t}=1 \mid \mathbf{x}_{t}=\mathbf{x}^{\mathrm{o}}, c_{i}\right)\right]=\mathrm{E}\left[\Phi\left(\mathbf{x}^{\mathrm{o}} \boldsymbol{\beta}+c_{i}\right)\right]$ for any given value $\mathbf{x}^{\mathrm{o}}$ of the explanatory variables. (In what follows, $\mathbf{x}^{\mathrm{o}}$ is a nonrandom vector of numbers that we choose as interesting values of the explanatory variables. For emphasis, we include an $i$ subscript on the random variables appearing in the expectations.) Writing $c_{i}=\psi+\overline{\mathbf{x}}_{i} \boldsymbol{\xi}+a_{i}$ and using iterated expectations, we have $\mathrm{E}\left[\Phi\left(\psi+\mathbf{x}^{\mathrm{o}} \boldsymbol{\beta}+\overline{\mathbf{x}}_{i} \boldsymbol{\xi}+a_{i}\right)\right]=\mathrm{E}\left[\mathrm{E}\left\{\Phi\left(\psi+\mathbf{x}^{\mathrm{o}} \boldsymbol{\beta}+\overline{\mathbf{x}}_{i} \boldsymbol{\xi}+a_{i}\right) \mid \mathbf{x}_{i}\right\}\right]$ [where the first expectation is with respect to $\left.\left(\mathbf{x}_{i}, a_{i}\right)\right]$. Using the same argument from Section 15.7.1, $\mathrm{E}\left[\Phi\left(\psi+\mathbf{x}^{\mathrm{o}} \boldsymbol{\beta}+\overline{\mathbf{x}}_{i} \boldsymbol{\xi}+a_{i}\right) \mid \mathbf{x}_{i}\right]=\Phi\left[\left\{\psi+\mathbf{x}^{\mathrm{o}} \boldsymbol{\beta}+\overline{\mathbf{x}}_{i} \boldsymbol{\xi}\right\} \cdot\left(1+\sigma_{a}^{2}\right)^{-1 / 2}\right]=\Phi\left(\psi_{a}+\mathbf{x}^{\mathrm{o}} \boldsymbol{\beta}_{a}+\overline{\mathbf{x}}_{i} \boldsymbol{\xi}_{a}\right)$, and so\\
$\mathrm{E}\left[\Phi\left(\mathbf{x}^{\mathrm{o}} \boldsymbol{\beta}+c_{i}\right)\right]=\mathrm{E}\left[\Phi\left(\psi_{a}+\mathbf{x}^{\mathrm{o}} \boldsymbol{\beta}_{a}+\overline{\mathbf{x}}_{i} \boldsymbol{\xi}_{a}\right)\right]$.

Because the only random variable in the right-hand-side expectation is $\overline{\mathbf{x}}_{i}$, a consistent estimator of the right-hand side of equation (15.75) is simply


\begin{equation*}
N^{-1} \sum_{i=1}^{N} \Phi\left(\hat{\psi}_{a}+\mathbf{x}^{\mathrm{o}} \hat{\boldsymbol{\beta}}_{a}+\overline{\mathbf{x}}_{i} \hat{\boldsymbol{\xi}}_{a}\right) \tag{15.76}
\end{equation*}


APEs can be estimated by evaluating expression (15.76) at two different values for $\mathbf{x}^{\circ}$ and forming the difference, or, for continuous variable $x_{j}$, by using the average across $i$ of $\hat{\beta}_{a j} \phi\left(\hat{\psi}_{a}+\mathbf{x}^{\mathrm{o}} \hat{\boldsymbol{\beta}}_{a}+\overline{\mathbf{x}}_{i} \hat{\boldsymbol{\xi}}_{a}\right)$ to get the approximate APE of a one-unit increase in $x_{j}$. See also Chamberlain (1984, equation (3.4)). If we use Chamberlain's more general version of assumption (15.73), $\mathbf{x}_{i}$ replaces $\overline{\mathbf{x}}_{i}$ everywhere.

Our focus on the APEs raises an interesting question: How does treating the $c_{i}$ 's as parameters to estimate-in a "fixed effects probit" analysis-affect estimation of the APEs? Given $\hat{c}_{i}, i=1, \ldots, N$ and $\hat{\boldsymbol{\beta}}$, the APEs could be based on\\
$N^{-1} \sum_{i=1}^{N} \Phi\left(\hat{c}_{i}+\mathbf{x}^{\mathrm{o}} \hat{\boldsymbol{\beta}}\right)$. Even though $\hat{\boldsymbol{\beta}}$ does not consistently estimate $\boldsymbol{\beta}$ and the $\hat{c}_{i}$ are estimates of the incidental parameters, it could be that the resulting estimates of the APEs have reasonable properties. In fact, recent progress has been made on this question. Fernández-Val (2009) contains theoretical results showing that the inconsistency in APEs constructed using the $\hat{c}_{i}$ and the accompanying $\hat{\boldsymbol{\beta}}$ is on the order of $T^{-1}$ (and, if there is no heterogeneity, the inconsistency is only $O\left(T^{-2}\right)$ ). In simulations, Hahn and Newey (2004) found small biases in the usual probit APEs; they suggested bias-corrected estimators that have inconsistency on the order of $T^{-2}$. In related work, Greene (2004), using simulations, found that for even moderate $T$ (say, $T \geq 5$ ), the bias in the partial effects evaluated at the average (as opposed to the average partial effect) can be quite small. All of these papers maintain the conditional independence assumption (15.67) and, in addition, the assumption that the covariates are independent, identically distributed across $t$. Hopefully, the findings carry over to more realistic scenarios that allow time dependence in the covariates as well as violations of (15.67).

Under assumptions (15.65), (15.66), and the more general version of assumption (15.73), Chamberlain (1980) suggested a minimum distance approach analogous to the linear case (see Section 14.6.2). Namely, obtain $\hat{\boldsymbol{\pi}}_{t}$ for each $t$ by running a crosssectional probit of $y_{i t}$ on $1, \mathbf{x}_{i}, i=1, \ldots, N$. The mapping from the structural parameters $\boldsymbol{\theta}_{a} \equiv\left(\psi_{a}, \boldsymbol{\beta}_{a}^{\prime}, \boldsymbol{\xi}_{a}^{\prime}\right)^{\prime}$ to the vector $\boldsymbol{\pi}$ is exactly as in the linear case (see Section 14.6.2). The variance matrix estimator for $\hat{\boldsymbol{\pi}}$ is obtained by pooling all $T$ probits and computing the robust variance matrix estimator in equation (13.53), with $\hat{\boldsymbol{\theta}}$ replaced by $\hat{\boldsymbol{\pi}}$; see also Chamberlain (1984, Section 4.5). The minimum distance approach leads to a straightforward test of $\mathrm{H}_{0}: \boldsymbol{\xi}_{a}=\mathbf{0}$, which is a test of assumption (15.69) that does not impose assumption (15.67).

Strict exogeneity of the covariates conditional on $c_{i}$ is critical for the previous analysis. As mentioned earlier, this assumption rules out lagged dependent variables, a case we consider explicitly in Section 15.8.4. But there are other cases where strict exogeneity is questionable. For example, suppose that $y_{i t}$ is an employment indicator for person $i$ in period $t$ and $w_{i t}$ is measure of recent arrests. It is possible that whether someone is employed in this period has an effect on future arrests. If so, then shocks that affect employment status could be correlated with future arrests, and such correlation would violate strict exogeneity. Whether this situation is empirically important is largely unknown.

On the one hand, correcting for an explanatory variable that is not strictly exogenous is quite difficult in nonlinear models; Wooldridge (2000) suggests one possible approach. On the other hand, obtaining a test of strict exogeneity is fairly easy. Let $\mathbf{w}_{i t}$ denote a $1 \times G$ subset of $\mathbf{x}_{i t}$ that we suspect of failing the strict exogeneity re-\\
quirement. Then a simple test is to add $\mathbf{w}_{i, t+1}$ as an additional set of covariates; under the null hypothesis, $\mathbf{w}_{i, t+1}$ should be insignificant. In implementing this test, we can use either RE probit or pooled probit, where, in either case, we lose the last time period. (In the pooled probit case, we should use a fully robust Wald or score test.) We should still obtain $\overline{\mathbf{x}}_{i}$ from all $T$ time periods, as the test is either based on the distribution of $\left(y_{i 1}, \ldots, y_{i, T-1}\right)$ given $\left(\mathbf{x}_{i 1}, \ldots, \mathbf{x}_{i T}\right)$ (RE probit) or on the marginal distributions of $y_{i t}$ given $\left(\mathbf{x}_{i 1}, \ldots, \mathbf{x}_{i T}\right), t=1, \ldots, T-1$ (pooled probit). If the test does not reject, it provides at least some justification for the strict exogeneity assumption.

\subsection*{15.8.3 Unobserved Effects Logit Models under Strict Exogeneity}
The unobserved effects probit model of the previous subsection has a logit counterpart. In the leading case, the heterogeneity, $c_{i}$, is still assumed to satisfy (15.69), but we replace assumption (15.66) with the response probability


\begin{equation*}
\mathrm{P}\left(y_{i t}=1 \mid \mathbf{x}_{i t}, c_{i}\right)=\Lambda\left(\mathbf{x}_{i t} \boldsymbol{\beta}+c_{i}\right), \quad t=1, \ldots, T, \tag{15.77}
\end{equation*}


where $\Lambda(\cdot)$ is the logistic function. If we maintain the full set of assumptions, given by the strict exogeneity condition (15.65), the conditional independence assumption (15.67), the normality assumption (15.69), and the model in equation (15.77), we arrive at the random effects logit model. In the GEE literature, the RE logit model is an example of a subject-specific model (like the RE probit model).

From a computational standpoint, the RE logit model is less desirable than the RE probit model. Even if we focus on a pooled method under assumptions (15.69) and (15.77), estimation is complicated because the response probability obtained by integrating out $c_{i}$,


\begin{equation*}
\mathrm{P}\left(y_{i t}=1 \mid \mathbf{x}_{i t}\right)=\int_{-\infty}^{\infty} \Lambda\left(\mathbf{x}_{i t} \boldsymbol{\beta}+c\right)\left(1 / \sigma_{c}\right) \phi\left(c / \sigma_{c}\right) \mathrm{d} c \tag{15.78}
\end{equation*}


does not have a closed form. In fact, one rarely sees equation (15.78) used as the sole basis for estimating $\boldsymbol{\beta}$ and $\sigma_{c}^{2}$, whether using a pooled binary response method or a GEE method under assumption (15.65). (Unlike in the probit model, $\boldsymbol{\beta}$ and $\sigma_{c}^{2}$ are separately identified in (15.78), although that does not mean they are easy to estimate.) Interestingly, although it seems natural to define (15.78), along with the strict exogeneity assumption, as constituting the PA version of the SS model, that is not the designation in the GEE literature. Instead, the PA model is specified as


\begin{equation*}
\mathrm{P}\left(y_{i t}=1 \mid \mathbf{x}_{i}\right)=\mathrm{P}\left(y_{i t}=1 \mid \mathbf{x}_{i t}\right)=\Lambda\left(\mathbf{x}_{i t} \boldsymbol{\beta}_{c}\right), \tag{15.79}
\end{equation*}


where we use the " $c$ " subscript on $\boldsymbol{\beta}_{c}$ to emphasize that the beta vector in equation (15.79) cannot be the same as $\boldsymbol{\beta}$ in equation (15.77) (and to draw an analogy with\\
probit). In fact, (15.79) is incompatible with (15.77); in specifying (15.79), we are abandoning the more structural model in (15.77) for the sake of expediency. In the GEE literature more generally, a PA model is usually specified to have a convenient functional form for the conditional mean, $\mathrm{E}\left(y_{i t} \mid \mathbf{x}_{i t}\right)$ without worrying about whether it can be derived from an underlying unobserved effects model.

If we estimate the model in (15.77) under the full set of RE logit assumptions (or even if we just use equation (15.78)), the partial effects at interesting values of $c$ are easier to obtain than the APEs: we just use differences or derivatives of $\Lambda\left(\mathbf{x}_{t} \boldsymbol{\beta}+c\right)$ with respect to elements of $\mathbf{x}_{t}$ and plug in interesting values of $c$, which can be determined because $c_{i}$ is distributed as $\operatorname{Normal}\left(0, \sigma_{c}^{2}\right)$. Further, we can apply the Mundlak-Chamberlain device just as in assumption (15.73), and then we can estimate the unconditional mean and variance of $c_{i}$ as in the probit case. If we want to allow $\mathrm{D}\left(c_{i} \mid \mathbf{x}_{i}\right)$ to depend on $\overline{\mathbf{x}}_{i}$ but do not want to use RE logit as the underlying model, we can adopt a strategy that extends the standard PA model strategy. Namely, we just specify\\
$\mathrm{P}\left(y_{i t}=1 \mid \mathbf{x}_{i}\right)=\Lambda\left(\psi_{a}+\mathbf{x}_{i t} \boldsymbol{\beta}_{a}+\overline{\mathbf{x}}_{i} \boldsymbol{\xi}_{a}\right), \quad t=1, \ldots, T$,\\
where we use the " $a$ " subscript to indicate a new set of parameters (while again drawing an analogy to the probit case). Given equation (15.80), we can estimate the parameters by pooled logit or GEE, and then, to obtain APEs, compute the ASF $N^{-1} \sum_{i=1}^{N} \Lambda\left(\hat{\psi}_{a}+\mathbf{x}_{t} \hat{\boldsymbol{\beta}}_{a}+\overline{\mathbf{x}}_{i} \hat{\boldsymbol{\xi}}_{a}\right)$ for given $\mathbf{x}_{t}$. Although equation (15.80) cannot be derived from (15.77), it might provide a good approximation to the APEs.

The real advantage of specifying equation (15.77) in place of (15.66) is that under the logit specification, we can obtain a $\sqrt{N}$-consistent, asymptotically normal estimator of $\boldsymbol{\beta}$ without any assumptions on $\mathrm{D}\left(c_{i} \mid \mathbf{x}_{i}\right)$, provided, of course, that each element of $\mathbf{x}_{i t}$ is time varying. In addition to assumption (15.77), we maintain the strict exogeneity assumption (15.65) and the conditional independence assumption (15.67).

How can we allow $c_{i}$ and $\mathbf{x}_{i}$ to be arbitrarily related in the unobserved effects logit model? In the linear case we used the FE or FD transformation to eliminate $c_{i}$ from the estimating equation. It turns out that a similar strategy works in the logit case, although the argument is more subtle. What we do is find the joint distribution of $\mathbf{y}_{i} \equiv\left(y_{i 1}, \ldots, y_{i T}\right)^{\prime}$ conditional on $\mathbf{x}_{i}, c_{i}$ and $n_{i} \equiv \sum_{t=1}^{T} y_{i t}$. It turns out that this conditional distribution does not depend on $c_{i}$, so that it is also the distribution of $\mathbf{y}_{i}$ given $\mathbf{x}_{i}$ and $n_{i}$. Therefore, we can use standard CMLE methods to estimate $\boldsymbol{\beta}$. (The fact that we can find a conditional distribution that does not depend on the $c_{i}$ is a feature of the logit functional form. Unfortunately, the same argument does not work for the unobserved effects probit model.)

First consider the $T=2$ case, where $n_{i}$ takes a value in $\{0,1,2\}$. Intuitively, the conditional distribution of $\left(y_{i 1}, y_{i 2}\right)^{\prime}$ given $n_{i}$ cannot be informative for $\boldsymbol{\beta}$ when $n_{i}=0$ or $n_{i}=2$ because these values completely determine the outcome on $\mathbf{y}_{i}$. However, for $n_{i}=1$,

$$
\begin{aligned}
\mathrm{P}\left(y_{i 2}=\right. & \left.1 \mid \mathbf{x}_{i}, c_{i}, n_{i}=1\right)=\mathrm{P}\left(y_{i 2}=1, n_{i}=1 \mid \mathbf{x}_{i}, c_{i}\right) / \mathrm{P}\left(n_{i}=1 \mid \mathbf{x}_{i}, c_{i}\right) \\
= & \mathrm{P}\left(y_{i 2}=1 \mid \mathbf{x}_{i}, c_{i}\right) \mathrm{P}\left(y_{i 1}=0 \mid \mathbf{x}_{i}, c_{i}\right) /\left\{\mathrm{P}\left(y_{i 1}=0, y_{i 2}=1 \mid \mathbf{x}_{i}, c_{i}\right)\right. \\
& \left.+\mathrm{P}\left(y_{i 1}=1, y_{i 2}=0 \mid \mathbf{x}_{i}, c_{i}\right)\right\} \\
= & \Lambda\left(\mathbf{x}_{i 2} \boldsymbol{\beta}+c_{i}\right)\left[1-\Lambda\left(\mathbf{x}_{i 1} \boldsymbol{\beta}+c_{i}\right)\right] /\left\{\left[1-\Lambda\left(\mathbf{x}_{i 1} \boldsymbol{\beta}+c_{i}\right)\right] \Lambda\left(\mathbf{x}_{i 2} \boldsymbol{\beta}+c_{i}\right)\right. \\
& \left.+\Lambda\left(\mathbf{x}_{i 1} \boldsymbol{\beta}+c_{i}\right)\left[1-\Lambda\left(\mathbf{x}_{i 2} \boldsymbol{\beta}+c_{i}\right)\right]\right\}=\Lambda\left[\left(\mathbf{x}_{i 2}-\mathbf{x}_{i 1}\right) \boldsymbol{\beta}\right] .
\end{aligned}
$$

Similarly, $\quad \mathbf{P}\left(y_{i 1}=1 \mid \mathbf{x}_{i}, c_{i}, n_{i}=1\right)=\Lambda\left[-\left(\mathbf{x}_{i 2}-\mathbf{x}_{i 1}\right) \boldsymbol{\beta}\right]=1-\Lambda\left[\left(\mathbf{x}_{i 2}-\mathbf{x}_{i 1}\right) \boldsymbol{\beta}\right]$. The conditional $\log$ likelihood for observation $i$ is\\
$\ell_{i}(\boldsymbol{\beta})=1\left[n_{i}=1\right]\left(w_{i} \log \Lambda\left[\left(\mathbf{x}_{i 2}-\mathbf{x}_{i 1}\right) \boldsymbol{\beta}\right]+\left(1-w_{i}\right) \log \left\{1-\Lambda\left[\left(\mathbf{x}_{i 2}-\mathbf{x}_{i 1}\right) \boldsymbol{\beta}\right]\right\}\right)$,\\
where $w_{i}=1$ if $\left(y_{i 1}=0, y_{i 2}=1\right)$ and $w_{i}=0$ if $\left(y_{i 1}=1, y_{i 2}=0\right)$. The CMLE is obtained by maximizing the sum of the $\ell_{i}(\boldsymbol{\beta})$ across $i$. The indicator function $1\left[n_{i}=1\right]$ selects out the observations for which $n_{i}=1$; as stated earlier, observations for which $n_{i}=0$ or $n_{i}=2$ do not contribute to the log likelihood. Interestingly, equation (15.81) is just a standard cross-sectional logit of $w_{i}$ on $\left(\mathbf{x}_{i 2}-\mathbf{x}_{i 1}\right)$ using the observations for which $n_{i}=1$. (This approach is analogous to differencing in the linear case with $T=2$.)

The CMLE from equation (15.81) is often called the fixed effects logit estimator and sometimes called the conditional logit estimator. We must emphasize that the FE logit estimator does not arise by treating the $c_{i}$ as parameters to be estimated along with $\boldsymbol{\beta}$. (This convention is confusing, as the FE probit estimator does estimate the $c_{i}$ along with $\boldsymbol{\beta}$.) As shown recently by Abrevaya (1997), the MLE of $\boldsymbol{\beta}$ that is obtained by maximizing the $\log$ likelihood over $\boldsymbol{\beta}$, and the $c_{i}$ has probability limit $2 \boldsymbol{\beta}$. (This finding extends a simple example due to Andersen, 1970; see also Hsiao, 1986, Section 7.3.)

Sometimes the CMLE is described as "conditioning on the unobserved effects in the sample." This description is misleading. What we have done is found a conditional density-which describes the subpopulation with $n_{i}=1$-that depends only on observable data and the parameter $\boldsymbol{\beta}$. One should not expend energy worrying about the loss of observations for which $y_{i 1}=y_{i 2}$. Because $\mathrm{D}\left(c_{i} \mid \mathbf{x}_{i}\right)$ is unrestricted, $c_{i}$ is free to vary as much as required to make $\mathrm{P}\left(y_{i t}=1 \mid \mathbf{x}_{i t}, c_{i}\right)$ arbitarily close to one (if\\
$y_{i t}=1, t=1,2$ ) or arbitrarily close to zero (if $y_{i t}=0, t=1,2$ ). When $y_{i t}$ does not change across $t, c_{i}$ can adjust to make both observed outcomes occur with certainty. Clearly, such data points contain no information for estimating $\boldsymbol{\beta}$, and so they should drop out of the estimation.

For general $T$ the log likelihood is more complicated, but it is tractable. First,


\begin{align*}
\mathrm{P}\left(y_{i 1}\right. & \left.=y_{1}, \ldots, y_{i T}=y_{T} \mid \mathbf{x}_{i}, c_{i}, n_{i}=n\right) \\
& =\mathrm{P}\left(y_{i 1}=y_{1}, \ldots, y_{i T}=y_{T} \mid \mathbf{x}_{i}, c_{i}\right) / \mathrm{P}\left(n_{i}=n \mid \mathbf{x}_{i}, c_{i}\right), \tag{15.82}
\end{align*}


and the numerator factors as $\mathrm{P}\left(y_{i 1}=y_{1} \mid \mathbf{x}_{i}, c_{i}\right) \cdots \mathrm{P}\left(y_{i T}=y_{T} \mid \mathbf{x}_{i}, c_{i}\right)$ by the conditional independence assumption. The denominator is the complicated part, but it is easy to describe: $\mathrm{P}\left(n_{i}=n \mid \mathbf{x}_{i}, c_{i}\right)$ is the sum of the probabilities of all possible outcomes of $\mathbf{y}_{i}$ such that $n_{i}=n$. Using the specific form of the logit function we can write


\begin{equation*}
\ell_{i}(\boldsymbol{\beta})=\log \left\{\exp \left(\sum_{t=1}^{T} y_{i t} \mathbf{x}_{i t} \boldsymbol{\beta}\right)\left[\sum_{\mathbf{a} \in R_{i}} \exp \left(\sum_{t=1}^{T} a_{t} \mathbf{x}_{i t} \boldsymbol{\beta}\right)\right]^{-1}\right\} \tag{15.83}
\end{equation*}


where $R_{i}$ is the subset of $\mathbb{R}^{T}$ defined as $\left\{\mathbf{a} \in \mathbb{R}^{T}: a_{t} \in\{0,1\}\right.$ and $\left.\sum_{t=1}^{T} a_{t}=n_{i}\right\}$. The log likelihood summed across $i$ can be used to obtain a $\sqrt{N}$-asymptotically normal estimator of $\boldsymbol{\beta}$, and all inference follows from conditional MLE theory. Observations for which equation (15.82) is zero or unity-and which therefore do not depend on $\boldsymbol{\beta}$-drop out of $\mathscr{L}(\boldsymbol{\beta})$. See Chamberlain (1984).

The FE logit estimator $\hat{\boldsymbol{\beta}}$ immediately gives us the effect of each element of $\mathbf{x}_{t}$ on the log-odds ratio, $\log \left\{\Lambda\left(\mathbf{x}_{t} \boldsymbol{\beta}+c\right) /\left[1-\Lambda\left(\mathbf{x}_{t} \boldsymbol{\beta}+c\right)\right]\right\}=\mathbf{x}_{t} \boldsymbol{\beta}+c$. Unfortunately, we cannot estimate the partial effects on the response probabilities unless we plug in a value for $c$. Because the distribution of $c_{i}$ is unrestricted-in particular, $\mathrm{E}\left(c_{i}\right)$ is not necessarily zero-it is hard to know what to plug in for $c$. In addition, we cannot estimate APEs, as doing so would require finding $\mathrm{E}\left[\Lambda\left(\mathbf{x}_{t} \boldsymbol{\beta}+c_{i}\right)\right]$, a task that apparently requires specifying a distribution for $c_{i}$.

The conditional logit approach also has the drawback of apparently requiring the conditional independence assumption (15.67) for consistency. As we saw in Section 15.8.2, if we are willing to make the normality assumption (15.73), the probit approach allows unrestricted serial dependence in $y_{i t}$ even after conditioning on $\mathbf{x}_{i}$ and $c_{i}$. This possibility may be especially important when several time periods are available.

Now that we have covered all of the leading estimation methods for unobserved effects models with strictly exogenous explanatory variables, we can provide an example comparing the different approaches. We use a panel data set on the labor force participation of married women from Chay and Hyslop (2001).

\begin{table}[h]
\begin{center}
\captionsetup{labelformat=empty}
\caption{Table 15.3\\
Panel Data Models for Married Women's Labor Force Participation}
\begin{tabular}{|l|l|l|l|l|l|l|l|l|}
\hline
Model & (1) Linear & \multicolumn{2}{|l|}{(2) Probit} & \multicolumn{2}{|l|}{(3) Chamberlain's RE Probit} & \multicolumn{2}{|l|}{(4) Chamberlain's RE Probit} & (5) FE Logit \\
\hline
\multirow[b]{2}{*}{Estimation Method} & Fixed Effects & \multicolumn{2}{|l|}{Pooled MLE} & \multicolumn{2}{|l|}{Pooled MLE} & \multicolumn{2}{|l|}{MLE} & MLE \\
\hline
 & Coefficient & Coefficient & APE & Coefficient & APE & Coefficient & APE & Coefficient \\
\hline
kids & -. 0389 (.0092) & -. 199 (.015) & -. 0660 (.0048) & -. 117 (.027) & -. 0389 (.0085) & -. 317 (.062) & -. 0403 (*) & -. 644 (.125) \\
\hline
lhinc & -. 0089 (.0046) & -. 211 (.024) & -. 0701 (.0079) & -. 029 (.014) & -. 0095 (.0048) & -. 078 (.041) & -. 0099 (*) & -. 184 (.083) \\
\hline
$\overline{k i d s}$ & - & - & - & -. 086 & - & -. 210 & - & - \\
\hline
$\overline{\text { lhinc }}$ & - & - & - & -. 250 & - & -. 646 & - & - \\
\hline
$\left(1+\hat{\sigma}_{a}^{2}\right)^{-1 / 2}$ & - & - &  & - &  &  & . 387 & - \\
\hline
Log likelihood & - &  & -16,556.67 &  & -16,516.44 & -8,990.09 &  & -2,003.42 \\
\hline
Number of women & 5,663 & \multicolumn{2}{|c|}{5,663} & \multicolumn{2}{|c|}{5,663} & \multicolumn{2}{|c|}{5,663} & 1,055 \\
\hline
\end{tabular}
\end{center}
\end{table}

Standard errors are in parentheses below all coefficients or APEs. For the linear model, pooled probit, and Chamberlain's RE probit estimated by pooled MLE, the standard errors are robust to arbitrary serial correlation.\\
All models include a full set of period dummy variables (unreported).\\
Columns (2), (3), and (4) also include the time-constant variables educ, black, age, and age ${ }^{2}$ (unreported).\\
Because of prohibitive computation times, bootstrap standard errors were not computed for the RE probit APEs estimated by MLE.

Example 15.5 (Panel Data Models for Women's Labor Force Participation): The data set LFP.RAW contains data on $N=5,663$ married women over $T=5$ periods, where the periods are spaced four months apart. The response variable is $l f p_{i t}$, a labor force participation indicator. The key explanatory variables are $k i d s_{i t}$ (number of children under 18) and lhinc $_{i t}=\log \left(\right.$ hinc $\left._{i t}\right)$ (where husband's income, hinc $_{i t}$, is in dollars per month and is positive for all $i$ and $t$ ). We also include the time-constant variables educ (years of schooling in the first period), black (a binary race indicator), age (age in the first period), and age ${ }^{2}$; these drop out of the linear FE and FE logit estimation.

The findings in Table 15.3 reveal a consistent pattern: allowing unobserved heterogeneity to be correlated with $k i d s_{i t}$ and lhinc ${ }_{i t}$ has important effects on the estimated APEs. Estimating the LPM by FE gives estimated coefficients of roughly -.039 and -. 009 on the $k i d s_{i t}$ and lhinc ${ }_{i t}$ variables, respectively, with the coefficient on lhinc ${ }_{i t}$\\
being marginally statistically significant. Each child is estimated to reduce the labor force participation probability by about .039 , while a $10 \%$ increase in a husband's income lowers the probability by only .0009 . If we use probit and assume $c_{i}$ is independent of $\mathbf{x}_{i}$, the APEs are much higher, especially on the income variable, which is almost 10 times the LPM coefficient. Could this difference be due to the different functional forms? Column (3) makes it clear that the difference in APEs between columns (1) and (2) is due to the restriction in (2) that $c_{i}$ is independent of $\mathbf{x}_{i}$. When we use the Chamberlain-Mundlak device and use pooled probit, we obtain APEs that are very similar to the LPM estimates. In fact, to four decimal places, the APE estimates on $k i d s_{i t}$ are identical. Plus, we can see directly that the coefficients on the time averages are very statistically significant and practically very large; each is much larger than the corresponding coefficient on the time-varying variable.

Using full MLE, rather than PMLE, to estimate Chamberlain's model does not change any conclusions. When we multiply the MLE coefficients by the scale factor .387 , the results are very similar to the pooled MLE coefficients $\left(-.1227\right.$ for kids ${ }_{i t}$ and -.030 for $l h i c_{i t}$ ). In addition, the APE estimates are very similar to the pooled estimation; interestingly, the bootstrap standard errors for the APEs are actually somewhat higher than for the pooled probit estimates.

Finally, column (5) contains the estimates from FE logit. As we discussed, the coefficient magnitudes are difficult to interpret. But the relative size is $.644 / .184=3.50$, which is not terribly different from, say, the ratio of coefficients using the pooled MLE estimates of Chamberlain's model, $.117 / .029=4.03$. But if we do not control for the time averages, the estimated ratios are much different, $.199 / .211=.94$. Because Chamberlain's approach yields APEs, and the estimates in column (3) allow for arbitrary serial dependence, it seems sensible to rely on these estimates (assuming, of course, that we believe $k i d s_{i t}$ and $l h i c_{i t}$ are strictly exogenous conditional on $c_{i}$ ).

Because we have covered several different possibilities for modeling and estimation in the context of unobserved effects with binary responses, it is helpful to have a simple summary of the strengths and weaknesses of each approach. Table 15.4 contains such a summary. Each method is evaluated by answers to five questions: (1) Is the response probability contained in the unit interval? (2) Is the distribution of the heterogeneity, given the covariates, restricted? (3) Is serial dependence allowed after accounting for $c_{i}$ ? (4) Are the partial effects at the mean heterogeneity identified? (5) Are the APEs identified?

Table 15.4 reveals a simple but important point that is easily missed: no procedure strictly dominates the others. When deciding on a method or methods, one needs to determine how important each factor is likely to be. For some features, this is easy. For example, hopefully we know whether having consistent parameter estimates is

\begin{table}[h]
\begin{center}
\captionsetup{labelformat=empty}
\caption{Table 15.4\\
Summary of Features of Models and Estimation Methods for Unobserved Effects Binary Response Models}
\begin{tabular}{|l|l|l|l|l|l|}
\hline
Model, Estimation Method & $\mathrm{P}\left(y_{i t}=\right. \left.1 \mid \mathbf{x}_{i t}, c_{i}\right)$ Bounded in $(0,1)$ ? & Restricts $\mathrm{D}\left(c_{i} \mid \mathbf{x}_{i}\right)$ ? & Idiosyncratic Serial Dependence? & Partial Effects at $\mathrm{E}\left(c_{i}\right)$ ? & APEs? \\
\hline
RE probit, MLE & Yes & Yes (independence, normal) & No & Yes & Yes \\
\hline
RE probit, pooled MLE & Yes & Yes (independence, normal) & Yes & No & Yes \\
\hline
RE probit, GEE & Yes & Yes (independence, normal) & Yes & No & Yes \\
\hline
Chamberlain's RE probit, MLE & Yes & Yes (linear mean, normal) & No & Yes & Yes \\
\hline
Chamberlain's RE probit, pooled MLE & Yes & Yes (linear mean, normal) & Yes & No & Yes \\
\hline
Chamberlain's RE probit, GEE & Yes & Yes (linear mean, normal) & Yes & No & Yes \\
\hline
LPM, within & No & No & Yes & Yes & Yes \\
\hline
FE logit, MLE & Yes & No & No & No & No \\
\hline
\end{tabular}
\end{center}
\end{table}

enough for our purposes or whether we want to estimate the magnitudes of effects. Unfortunately, it is difficult to decide issues that are essentially empirical in nature, such as how important functional form restrictions are on $\mathrm{P}\left(y_{i t}=1 \mid \mathbf{x}_{i t}, c_{i}\right)$ or $\mathrm{D}\left(c_{i} \mid \mathbf{x}_{i}\right)$.

\subsection*{15.8.4 Dynamic Unobserved Effects Models}
Dynamic models that also contain unobserved effects are important in testing theories and evaluating policies. Here we cover one class of models that illustrates the important points for general dynamic models and is of considerable interest in its own right. Our treatment follows Wooldridge (2005b).

Suppose we date our observations starting at $t=0$, so that $y_{i 0}$ is the first observation on $y$. For $t=1, \ldots, T$ we are interested in the dynamic unobserved effects model


\begin{equation*}
\mathrm{P}\left(y_{i t}=1 \mid y_{i, t-1}, \ldots, y_{i 0}, \mathbf{z}_{i}, c_{i}\right)=G\left(\mathbf{z}_{i t} \boldsymbol{\delta}+\rho y_{i, t-1}+c_{i}\right), \tag{15.84}
\end{equation*}


where $\mathbf{z}_{i t}$ is a vector of contemporaneous explanatory variables, $\mathbf{z}_{i}=\left(\mathbf{z}_{i 1}, \ldots, \mathbf{z}_{i T}\right)$, and $G$ can be the probit or logit function. There are several important points about this model. First, the $\mathbf{z}_{i t}$ are assumed to satisfy a strict exogeneity assumption (conditional on $c_{i}$ ), since $\mathbf{z}_{i}$ appears in the conditioning set on the left-hand side of equation (15.84), but only $\mathbf{z}_{i t}$ appears on the right-hand side. Second, the probability of success at time $t$ is allowed to depend on the outcome in $t-1$ as well as unobserved\\
heterogeneity, $c_{i}$. We saw the linear version in Section 11.6.2. Of particular interest is the hypothesis $\mathrm{H}_{0}: \rho=0$. Under this null, the response probability at time $t$ does not depend on past outcomes once $c_{i}$ (and $\mathbf{z}_{i}$ ) have been controlled for. Even if $\rho=0$, $\mathrm{P}\left(y_{i t}=1 \mid y_{i, t-1}, \mathbf{z}_{i}\right) \neq \mathrm{P}\left(y_{i t}=1 \mid \mathbf{z}_{i}\right)$ owing to the presence of $c_{i}$. But economists are interested in whether there is state dependence-that is, $\rho \neq 0$ in equation (15.84)after controlling for the unobserved heterogeneity, $c_{i}$.

We might also be interested in the effects of $\mathbf{z}_{t}$, as it may contain policy variables. Then, equation (15.84) simply captures the fact that, in addition to an unobserved effect, behavior may depend on past observed behavior.

How can we estimate $\boldsymbol{\delta}$ and $\rho$ in equation (15.84), in addition to quantities such as APEs? First, we can always write


\begin{align*}
f\left(y_{1}, y_{2}, \ldots, y_{T} \mid y_{0}, \mathbf{z}, c ; \boldsymbol{\beta}\right) & =\prod_{t=1}^{T} f\left(y_{t} \mid y_{t-1}, \ldots y_{1}, y_{0}, \mathbf{z}_{t}, c ; \boldsymbol{\beta}\right) \\
& =\prod_{t=1}^{T} G\left(\mathbf{z}_{t} \boldsymbol{\delta}+\rho y_{t-1}+c\right)^{y_{t}}\left[1-G\left(\mathbf{z}_{t} \boldsymbol{\delta}+\rho y_{t-1}+c\right)\right]^{1-y_{t}} \tag{15.85}
\end{align*}


With fixed- $T$ asymptotics, this density, because of the unobserved effect $c$, does not allow us to construct a log-likelihood function that can be used to estimate $\boldsymbol{\beta}$ consistently. Just as in the case with strictly exogenous explanatory variables, treating the $c_{i}$ as parameters to be estimated does not result in consistent estimators of $\boldsymbol{\delta}$ and $\rho$ as $N \rightarrow \infty$. In fact, the simulations in Heckman (1981) show that the incidental parameters problem is even more severe in dynamic models. What we should do is integrate out the unobserved effect $c$, as we discussed generally in Section 13.9.2.

Our need to integrate $c$ out of the distribution raises the issue of how we treat the initial observations, $y_{i 0}$; this is usually called the initial conditions problem. One possibility is to treat each $y_{i 0}$ as a nonstochastic starting position for each $i$. Then, if $c_{i}$ is assumed to be independent of $\mathbf{z}_{i}$ (as in a pure RE environment), equation (15.85) can be integrated against the density of $c$ to obtain the density of $\left(y_{1}, y_{2}, \ldots, y_{T}\right)$ given $\mathbf{z}$; this density also depends on $y_{0}$ through $f\left(y_{1} \mid y_{0}, c, \mathbf{z}_{1} ; \boldsymbol{\beta}\right)$. We can then apply CMLE. Although treating the $y_{i 0}$ as nonrandom simplifies estimation, it is undesirable because it effectively means that $c_{i}$ and $y_{i 0}$ are independent, a very strong assumption.

There seems to be confusion in the literature about when it is plausible to treat the $y_{i 0}$ as fixed, and therefore independent of $c_{i}$. Some have justified this assumption when one observes the process generating $y_{i t}$ from its beginning, as would happen if we follow the employment history of a cohort of high school graduates who do not pursue additional education, with $y_{i 0}$ being an employment indicator in the initial\\
postgraduation period. Does it make sense to assume $y_{i 0}$ and $c_{i}$ are independent? Almost certainly not. Because $c_{i}$ contains unobserved attributes that affect $y_{i t}$ in periods $t \geq 1$, it is almost certain that an individual's initial employment status is related to $c_{i}$. In this example and most others, treating the $y_{i 0}$ as independent of $c_{i}$ is risky, regardless of when the process underlying the panel data actually started.

Another possibility is to first specify a density for $y_{i 0}$ given ( $\mathbf{z}_{i}, c_{i}$ ) and to multiply this density by equation (15.85) to obtain $f\left(y_{0}, y_{1}, y_{2}, \ldots, y_{T} \mid \mathbf{z}, c ; \boldsymbol{\beta}, \boldsymbol{\gamma}\right)$. Next, a density for $c_{i}$ given $\mathbf{z}_{i}$ can be specified. Finally, $f\left(y_{0}, y_{1}, y_{2}, \ldots, y_{T} \mid \mathbf{z}, c ; \boldsymbol{\beta}, \boldsymbol{\gamma}\right)$ is integrated against the density $h(c \mid \mathbf{z} ; \boldsymbol{\alpha})$ to obtain the density of $\left(y_{i 0}, y_{i 1}, y_{i 2}, \ldots, y_{i T}\right)$ given $\mathbf{z}_{i}$. This density can then be used in an MLE analysis. The problem with this approach is that finding the density of $y_{i 0}$ given $\left(\mathbf{z}_{i}, c_{i}\right)$ is very difficult, if not impossible, even if the process is assumed to be in equilibrium. For discussion, see Hsiao (2003, Section 7.5).

Heckman (1981) suggests approximating the conditional density of $y_{i 0}$ given $\left(\mathbf{z}_{i}, c_{i}\right)$ and then specifying a density for $c_{i}$ given $\mathbf{z}_{i}$. For example, we might assume that $y_{i 0}$ follows a probit model with success probability $\Phi\left(\eta+\mathbf{z}_{i} \boldsymbol{\pi}+\gamma c_{i}\right)$ and specify the density of $c_{i}$ given $\mathbf{z}_{i}$ as normal. Once these two densities are given, they can be multiplied by equation (15.85), and $c$ can be integrated out to approximate the density of $\left(y_{i 0}, y_{i 1}, y_{i 2}, \ldots, y_{i T}\right)$ given $\mathbf{z}_{i}$; see Hsiao (2003, Section 7.5) or Wooldridge (2005b).

Heckman's (1981) approach attempts to find or approximate the joint distribution of $\left(y_{i 0}, y_{i 1}, y_{i 2}, \ldots, y_{i T}\right)$ given $\mathbf{z}_{i}$. We discussed an alternative approach in Section 13.9.2: obtain the joint distribution of $\left(y_{i 1}, y_{i 2}, \ldots, y_{i T}\right)$ conditional on $\left(y_{i 0}, \mathbf{z}_{i}\right)$. This allows us to remain agnostic about the distribution of $y_{i 0}$ given $\left(\mathbf{z}_{i}, c_{i}\right)$, which is the primary source of difficulty in Heckman's approach. If we can find the density of $\left(y_{i 1}, y_{i 2}, \ldots, y_{i T}\right)$ given $\left(y_{i 0}, \mathbf{z}_{i}\right)$, in terms of $\boldsymbol{\beta}$ and other parameters, then we can use standard CMLE methods: we are simply conditioning on $y_{i 0}$ in addition to $\mathbf{z}_{i}$. It is important to see that using the density of $\left(y_{i 1}, y_{i 2}, \ldots, y_{i T}\right)$ given $\left(y_{i 0}, \mathbf{z}_{i}\right)$ is not the same as treating $y_{i 0}$ as nonrandom. Indeed, the model with $c_{i}$ independent of $y_{i 0}$, given $\mathbf{z}_{i}$, is a special case.

To obtain $f\left(y_{1}, y_{2}, \ldots, y_{T} \mid y_{i 0}, \mathbf{z}_{i}\right)$, we need to propose a density for $c_{i}$ given $\left(y_{i 0}, \mathbf{z}_{i}\right)$. This approach is very much like Chamberlain's (1980) approach to static probit models with unobserved effects, except that we now condition on $y_{i 0}$ as well. (Since the density of $c_{i}$ given $\mathbf{z}_{i}$ is not restricted by the specification (15.85), our choice of the density of $c_{i}$ given ( $y_{i 0}, \mathbf{z}_{i}$ ) is not logically restricted in any way.) Given a density $h\left(c \mid y_{0}, \mathbf{z} ; \gamma\right)$, which depends on a vector of parameters $\gamma$, we have

$$
f\left(y_{1}, y_{2}, \ldots, y_{T} \mid y_{0}, \mathbf{z}, \boldsymbol{\theta}\right)=\int_{-\infty}^{\infty} f\left(y_{1}, y_{2}, \ldots, y_{T} \mid y_{0}, \mathbf{z}, c ; \boldsymbol{\beta}\right) h\left(c \mid y_{0}, \mathbf{z} ; \boldsymbol{\gamma}\right) \mathrm{d} c
$$

See Property CD. 2 in Chapter 13. The integral can be replaced with a weighted average if the distribution of $c$ is discrete. When $G=\Phi$ in the model (15.84)-the leading case-a very convenient choice for $h\left(c \mid y_{0}, \mathbf{z} ; \boldsymbol{\gamma}\right)$ is $\operatorname{Normal}\left(\psi+\xi_{0} y_{i 0}+\mathbf{z}_{i} \boldsymbol{\xi}\right.$, $\sigma_{a}^{2}$ ), which follows by writing $c_{i}=\psi+\xi_{0} y_{i 0}+\mathbf{z}_{i} \boldsymbol{\xi}+a_{i}$, where $a_{i} \sim \operatorname{Normal}\left(0, \sigma_{a}^{2}\right)$ and independent of $\left(y_{i 0}, \mathbf{z}_{i}\right)$. Then we can write\\
$y_{i t}=1\left[\psi+\mathbf{z}_{i t} \boldsymbol{\delta}+\rho y_{i, t-1}+\xi_{0} y_{i 0}+\mathbf{z}_{i} \boldsymbol{\xi}+a_{i}+e_{i t}>0\right]$,\\
so that $y_{i t}$ given $\left(y_{i, t-1}, \ldots, y_{i 0}, \mathbf{z}_{i}, a_{i}\right)$ follows a probit model and $a_{i}$ given ( $y_{i 0}, \mathbf{z}_{i}$ ) is distributed as $\operatorname{Normal}\left(0, \sigma_{a}^{2}\right)$. Therefore, the density of $\left(y_{i 1}, \ldots, y_{i T}\right)$ given $\left(y_{i 0}, \mathbf{z}_{i}\right)$ has exactly the form in equation (15.70), where $\mathbf{x}_{i t}=\left(1, \mathbf{z}_{i t}, y_{i, t-1}, y_{i 0}, \mathbf{z}_{i}\right)$ and with $a$ and $\sigma_{a}$ replacing $c$ and $\sigma_{c}$, respectively. Conveniently, this finding means that we can use standard RE probit software to estimate $\psi, \boldsymbol{\delta}, \rho, \xi_{0}, \boldsymbol{\xi}$, and $\sigma_{a}^{2}$ : we simply expand the list of explanatory variables to include $y_{i 0}$ and $\mathbf{z}_{i}$ in each time period. (The approach that treats $y_{i 0}$ and $\mathbf{z}_{i}$ as fixed omits $y_{i 0}$ and $\mathbf{z}_{i}$ in each time period.) It is simple to test $\mathrm{H}_{0}: \rho=0$, which means there is no state dependence once we control for an unobserved effect.

We can estimate APEs as in Chamberlain's unobserved effects probit model with strictly exogenous explanatory variables, but now we average out the initial condition along with the leads and lags of all strictly exogenous variables. Let $\mathbf{z}_{t}$ and $y_{t-1}$ be given values of the explanatory variables. Then the ASF, $\mathrm{E}\left[\Phi\left(\mathbf{z}_{t} \boldsymbol{\delta}+\rho y_{t-1}+c_{i}\right)\right]= \mathrm{E}\left[\Phi\left(\psi_{a}+\mathbf{z}_{t} \boldsymbol{\delta}_{a}+\rho_{a} y_{t-1}+\xi_{a 0} y_{i 0}+\mathbf{z}_{i} \xi_{a}\right)\right]$, is consistently estimated as $\widehat{\operatorname{ASF}}\left(\mathbf{z}_{t}, y_{t-1}\right)= N^{-1} \sum_{i=1}^{N} \Phi\left(\hat{\psi}_{a}+\mathbf{z}_{t} \hat{\boldsymbol{\delta}}_{a}+\hat{\rho}_{a} y_{t-1}+\hat{\xi}_{a 0} y_{i 0}+\mathbf{z}_{i} \hat{\boldsymbol{\xi}}_{a}\right)$, where the $a$ subscript denotes that the original coefficients have been multiplied by $\left(1+\hat{\sigma}_{a}^{2}\right)^{-1 / 2}$, and $\hat{\psi}, \hat{\boldsymbol{\delta}}, \hat{\rho}, \hat{\xi}, \hat{\boldsymbol{\xi}}$, and $\hat{\sigma}_{a}^{2}$ are the CMLEs that would be reported directly by any econometrics package that estimates RE probit models. We can take derivatives of this expression with respect to continuous elements of $\mathbf{z}_{t}$, or take differences with respect to discrete elements. Of particular interest is to alternatively set $y_{t-1}=1$ and $y_{t-1}=0$ and obtain the change in the probability that $y_{i t}=1$ when $y_{t-1}$ goes from zero to one. Probably we would average across $\mathbf{z}_{i t}$, too. To obtain a single APE, we can also average across all time periods. The key is that we average across $\left(y_{i 0}, \mathbf{z}_{i}\right)$ to estimate the ASF, and then from there we decide which partial effects are of interest. Wooldridge (2005b) contains further discussion, and we compute the APE for the lagged dependent variable in the example below.

In estimating the dynamic model just described, it is important to understand that a simpler, pooled estimation method does not consistently estimate the scaled parameters or the APEs. In other words, we cannot simply use pooled probit of $y_{i t}$ on $1, \mathbf{z}_{i t}, y_{i, t-1}, y_{i 0}, \mathbf{z}_{i}$. The problem is that, while $\mathrm{P}\left(y_{i t}=1 \mid \mathbf{z}_{i}, y_{i, t-1}, \ldots, y_{i 0}, a_{i}\right)= \Phi\left(\psi+\mathbf{z}_{i t} \boldsymbol{\delta}+\rho y_{i, t-1}+\xi_{0} y_{i 0}+\mathbf{z}_{i} \boldsymbol{\xi}+a_{i}\right)$, it is not true that $\mathrm{P}\left(y_{i t}=1 \mid \mathbf{z}_{i}, y_{i, t-1}, \ldots, y_{i 0}\right)$\\
$=\Phi\left(\psi_{a}+\mathbf{z}_{i t} \boldsymbol{\delta}_{a}+\rho_{a} y_{i, t-1}+\xi_{a 0} y_{i 0}+\mathbf{z}_{i} \xi_{a}\right)$ unless $a_{i}$ is identically zero, which means that $c_{i}$ is a deterministic linear function of $\left(y_{i 0}, \mathbf{z}_{i}\right)$. Correlation between $y_{i, t-1}$ and $a_{i}$ means that $\mathrm{P}\left(y_{i t}=1 \mid \mathbf{z}_{i}, y_{i, t-1}, \ldots, y_{i 0}\right)$ does not follow a probit model with index that depends on the scaled coefficients of interest. Therefore, we should use the RE probit approach, and not pooled probit (as in Section 15.8.1). For comparison purposes, we can estimate the dynamic model without the unobserved effect, $c_{i}$, and then pooled probit is the appropriate estimation method.

Example 15.6 (Dynamic Women's LFP Equation): We now use the data in LFP. RAW to estimate a model for $\mathrm{P}\left(\right.$ lfp $_{i t}=1 \mid$ kids $_{i t}$, lhinc $_{i t}$, lfp $\left._{i, t-1}, c_{i}\right)$, where one lag of labor force participation is assumed to suffice for the dynamics and $\left\{\left(\right.\right.$ kids $_{i t}$, lhinc $\left.\left._{i t}\right): t=1, \ldots, T\right\}$ is assumed to be strictly exogenous conditional on $c_{i}$. Also, we include the time-constant variables educ, black, age, and age ${ }^{2}$ and a full set of time-period dummies. (We start with five periods and lose one with the lag. Therefore, we estimate the model using four years of data.) We include among the regressors the initial value, $l f p_{i 0}$, kids $_{i 1}$ through kids $_{i 4}$, and $\operatorname{lhinc}_{i 1}$ through $\operatorname{lhinc}_{i 4}$. (To keep the notation consistent with the previous development, we implicitly relabel the time periods in the data set.) Estimating the model by RE probit gives $\hat{\rho}=1.541$ (se $=.067$ ), and so, even after controlling for unobserved heterogeneity, there is strong evidence of state dependence. But to obtain the size of the effect, we compute the APE for lfp $p_{t-1}$. The calculation involves averaging $\Phi\left(\hat{\psi}_{a}+\mathbf{z}_{i t} \hat{\boldsymbol{\delta}}_{a}+\hat{\rho}_{a}+\hat{\xi}_{a 0} y_{i 0}+\right. \left.\mathbf{z}_{i} \hat{\boldsymbol{\xi}}_{a}\right)-\Phi\left(\hat{\psi}_{a}+\mathbf{z}_{i t} \hat{\boldsymbol{\delta}}_{a}+\hat{\xi}_{a 0} y_{i 0}+\mathbf{z}_{i} \hat{\boldsymbol{\xi}}_{a}\right)$ across all $t$ and $i$; we must be sure to scale the original coefficients by $\left(1+\hat{\sigma}_{a}^{2}\right)^{-1 / 2}$, where, in this application, $\hat{\sigma}_{a}^{2}=1.103$. The APE estimated from this method is about .260 (panel bootstrap standard error is .026 with 500 replications). In other words, averaged across all women and all time periods, the probability of being in the labor force at time $t$ is about .26 higher if the woman was in the labor force at time $t-1$ than if she was not. This estimate controls for unobserved heterogeneity, number of young children, husband's income, and the woman's education, race, and age.

It is instructive to compare the APE with the estimate of a dynamic probit model that ignores $c_{i}$. In this case, we just use pooled probit of lfp $p_{i t}$ on 1, kids ${ }_{i t}$, lhinc ${ }_{i t}$, $l f p_{i, t-1}$, educ $_{i}$, black $_{i}$, age $_{i}$, and age $_{i}^{2}$ and include a full set of period dummies. The coefficient on $l f p_{i, t-1}$ is 2.876 ( $\mathrm{se}=.027$ ), which is much higher than in the dynamic RE probit model. More important, the APE for state dependence is about .837 (panel bootstrap se $=.005$ ), which is much higher than when heterogeneity is controlled for. Therefore, in this example, much of the persistence in labor force participation of married women is accounted for by the unobserved heterogeneity. There is still some state dependence, but its value is much smaller than a simple dynamic probit indicates.

As mentioned earlier, Wooldridge (2000) proposes an extension of the preceding approach to allow other not strictly exogenous explanatory variables to appear in unobserved effects probit models. Generally, if $w_{i t}$ is a variable that is sequentially but not strictly exogenous, we model $\mathrm{P}\left(y_{i t}=1 \mid w_{i t}, \mathbf{z}_{i t}, y_{i, t-1}, c_{i 1}\right)$ and assume this is the same probability when we include all lags of $w_{i t}$, further lags of $y_{i t}$, and the entire history of the strictly exogenous variables, $\mathbf{z}_{i t}$. (Allowing a lagged value of $w_{i t}$ only changes the notation.) We then add a model for the density of $\mathbf{D}\left(w_{i t} \mid \mathbf{z}_{i t}, w_{i, t-1}, y_{i, t-1}\right.$, $c_{i 2}$ ), assuming that one lag each of $w_{i t}$ and $y_{i t}$ is sufficient to capture the dynamics, and that $\mathbf{z}_{i t}$ is also strictly exogenous in this distribution. The joint density of $\left(y_{i t}, w_{i t}\right)$ given $\quad\left(\mathbf{z}_{i}, w_{i, t-1}, \ldots, w_{i 0}, y_{i, t-1}, \ldots, y_{i 0}, \mathbf{c}_{i}\right) \quad$ is $\quad f_{t 1}\left(y_{t} \mid w_{t}, \mathbf{z}_{i t}, y_{i, t-1}, c_{i 1}\right) \cdot f_{t 2}\left(w_{t} \mid \mathbf{z}_{i t}\right.$, $\left.w_{i, t-1}, y_{i, t-1}, c_{i 2}\right) \equiv g_{t}\left(y_{t}, w_{t} \mid \mathbf{z}_{i}, w_{i, t-1}, \ldots, w_{i 0}, y_{i, t-1}, \ldots, y_{i 0}, c_{i 1}, c_{i 2}\right)$. By multiplying these densities from $t=1$ to $T$ we obtain the density of $\left(y_{i 1}, w_{i 1}, \ldots, y_{i T}, w_{i T}\right)$ given $\left(\mathbf{z}_{i}, y_{i 0}, w_{i 0}, \mathbf{c}_{i}\right)$. Now, as before, if we specify a density for $\mathrm{D}\left(\mathbf{c}_{i} \mid \mathbf{z}_{i}, y_{i 0}, w_{i 0}\right)$, we can obtain a density for $\mathrm{D}\left(\mathbf{y}_{i}, \mathbf{w}_{i} \mid \mathbf{z}_{i}, y_{i 0}, w_{i 0}\right)$ by integrating out $\mathbf{c}_{i}$. We can construct a log-likelihood function for estimating the parameters in both conditional densities, as well as the parameters in the heterogeneity distribution. The problem is significantly harder if we truly allow two sources of heterogeneity in $\mathbf{c}_{i}=\left(c_{i 1}, c_{i 2}\right)$, especially if we allow them to be correlated, but conditioning on the initial conditions $\left(y_{i 0}, w_{i 0}\right)$ helps simplify the estimation. Wooldridge (2000) provides more details when $w_{i t}$ is a binary response that may react to changes in past values of $y_{i t}$, thereby causing it to violate the strict exogeneity assumption.

\subsection*{15.8.5 Probit Models with Heterogeneity and Endogenous Explanatory Variables}
The previous methods assumed that the explanatory variables satisfy a strict exogeneity assumption or a sequential exogeneity assumption. We can also use probit models to account for unobserved heterogeneity and contemporaneous endogeneity. We focus on two cases: (1) the endogenous explanatory variable has a conditional normal distribution and (2) the endogenous explanatory variable follows a reduced form probit. The reasons for these restrictions will become clear, as we must adapt the methods from Section 15.7.

We can write the model as


\begin{equation*}
y_{i t 1}=1\left[\mathbf{z}_{i t 1} \boldsymbol{\delta}_{1}+\alpha_{1} y_{i t 2}+c_{i 1}+u_{i t 1} \geq 0\right], \quad u_{i t 1} \mid \mathbf{z}_{i}, c_{i 1} \sim \operatorname{Normal}(0,1), \tag{15.86}
\end{equation*}


where $y_{i t 1}$ is the binary response, $y_{i t 2}$ is the endogenous explanatory variable, $c_{i 1}$ is the unobserved heterogeneity, $\left\{u_{i t 1}: t=1, \ldots, T\right\}$ is the sequence of idiosyncratic errors, and $\mathbf{z}_{i}=\left(\mathbf{z}_{i 1}, \ldots, \mathbf{z}_{i T}\right)$ is the sequence of strictly exogenous variables (conditional on $c_{i 1}$ ). Notice that $u_{i t 1}$ is assumed to be independent of $\left(\mathbf{z}_{i}, c_{i 1}\right)$. We would want to include a full set of period dummies in $\mathbf{z}_{i t}$.

Our focus will be on estimating APEs. Regardless of the nature of $y_{t 2}$, these are easily obtained from the ASF. At time $t$,\\
$\operatorname{ASF}\left(\mathbf{z}_{t 1}, y_{t 2}\right)=\mathrm{E}_{c_{i 1}}\left[\Phi\left(\mathbf{z}_{t 1} \boldsymbol{\delta}_{1}+\alpha_{1} y_{t 2}+c_{i 1}\right)\right]$,\\
where $\mathrm{E}_{c_{i 1}}(\cdot)$ denotes the expected value with respect to the distribution of $c_{i 1}$. As we did previously, we specify a model for the conditional distribution $\mathrm{D}\left(c_{i 1} \mid \mathbf{z}_{i}\right)$, the distribution conditional only on the strictly exogenous explanatory variables. Not suprisingly, a normal distribution is convenient:\\
$c_{i 1}=\psi_{1}+\overline{\mathbf{z}}_{i} \xi_{1}+a_{i 1}, \quad a_{i 1} \mid \mathbf{z}_{i} \sim \operatorname{Normal}\left(0, \sigma_{a_{1}}^{2}\right)$,\\
where $\overline{\mathbf{z}}_{i}$ contains the time averages of all strictly exogenous variables (except any aggregate time effects, such as period dummies). In particular, if at time $t$ we have $\mathbf{z}_{i t}=\left(\mathbf{z}_{i t 1}, \mathbf{z}_{i t 2}\right)$, then the time averages of $\mathbf{z}_{i t 2}$ are also in equation (15.88). Plugging into (15.86) and doing simple algebra gives


\begin{align*}
y_{i t 1} & =1\left[\mathbf{z}_{i t 1} \boldsymbol{\delta}_{1}+\alpha_{1} y_{i t 2}+\psi_{1}+\overline{\mathbf{z}}_{i} \xi_{1}+a_{i 1}+u_{i t 1} \geq 0\right] \\
& \equiv 1\left[\mathbf{z}_{i t 1} \boldsymbol{\delta}_{a 1}+\alpha_{a 1} y_{i t 2}+\psi_{a 1}+\overline{\mathbf{z}}_{i} \xi_{a 1}+e_{i t 1} \geq 0\right] \tag{15.89}
\end{align*}


where $e_{i t 1} \equiv\left(a_{i 1}+u_{i t 1}\right) /\left(1+\sigma_{a_{1}}^{2}\right)^{1 / 2}$ has a standard normal distribution conditional on $\mathbf{z}_{i}$. On the parameters, the " $a$ " subscript denotes division by $\left(1+\sigma_{a_{1}}^{2}\right)^{1 / 2}$, for example, $\boldsymbol{\delta}_{a 1} \equiv \boldsymbol{\delta}_{1} /\left(1+\sigma_{a_{1}}^{2}\right)^{1 / 2}$. Now the average structural function can be obtained as

$$
\begin{gathered}
\mathrm{E}_{\left(\overline{\mathbf{z}}_{i}, e_{i t 1}\right)}\left\{1\left[\mathbf{z}_{t 1} \boldsymbol{\delta}_{a 1}+\alpha_{a 1} y_{t 2}+\psi_{a 1}+\overline{\mathbf{z}}_{i} \boldsymbol{\xi}_{a 1}+e_{i t 1} \geq 0\right]\right\} \\
=\mathrm{E}_{\overline{\mathbf{z}}_{i}}\left[\Phi\left(\mathbf{z}_{t 1} \boldsymbol{\delta}_{a 1}+\alpha_{a 1} y_{t 2}+\psi_{a 1}+\overline{\mathbf{z}}_{i} \boldsymbol{\xi}_{a 1}\right)\right]
\end{gathered}
$$

where the equality follows by iterated expectations. It follows that we can consistently estimate the APEs by averaging out $\overline{\mathbf{z}}_{i}$ if we can estimate the scaled coefficients. But these scaled coefficients are precisely what we estimate if we apply one of the MLE methods from Section 15.7.2 or 15.7.3, depending on whether $y_{i t 2}$ has a conditional normal distribution or follows a probit. In the former case, we can write a reduced form as\\
$y_{i t 2}=\mathbf{z}_{i t} \boldsymbol{\delta}_{2}+\psi_{2}+\overline{\mathbf{z}}_{i} \xi_{2}+v_{i t 2}, \quad v_{i t 2} \mid \mathbf{z}_{i} \sim \operatorname{Normal}\left(0, \tau_{2}^{2}\right), t=1, \ldots, T$.

Notice that the instruments for $y_{i t 2}$ omitted from the estimating equation (15.89) are $\mathbf{z}_{i t 2}$; all elements of $\overline{\mathbf{z}}_{i}$ are in equations (15.89) and the reduced form (15.90). Then, we use pooled MLE directly on equations (15.89) and (15.90), and then compute the APEs exactly as we did before, possibly computing them for each time period. (We might call this approach pooled IV probit.) Bootstrapping is a simple (though\\
computationally expensive) way to obtain proper standard errors. If $y_{i t 2}$ is binary, we replace the right hand side of (15.90) with $1\left[\mathbf{z}_{i t} \boldsymbol{\delta}_{2}+\psi_{2}+\overline{\mathbf{z}}_{i} \boldsymbol{\xi}_{2}+v_{i t 2} \geq 0\right]$, set $\tau_{2}^{2}=1$, and apply pooled bivariate probit.

Pooled estimation is very convenient because any routine that allows estimation of the models for cross section data can be used for panel data, provided robust standard errors and test statistics are computed to account for the neglected time dependence. A full MLE method, which would account for the presence of $a_{i 1}$ and assume, in addition, that the $u_{i t 1}$ are serially independent, would be much more computationally intensive and would not be robust to serial correlation in the idiosyncratic errors. More importantly, MLE jointly across time would require specification of the joint distribution of $\left\{\left(e_{i t 1}, v_{i t 2}\right): t=1, \ldots, T\right\}$, and not just the bivariate distribution of ( $e_{i t 1}, v_{i t 2}$ ) for each $t$. Assuming independence across $t$ would be unrealistic, and allowing for realistic correlation over time would be computationally expensive.

When $y_{i t 2}$ is continuous, a control function approach is also available. Using manipulations similar to those above and in Section 15.7.2, we can write\\
$y_{i t 1}=1\left[\mathbf{z}_{i t 1} \boldsymbol{\delta}_{g 1}+\alpha_{g 1} y_{i t 2}+\theta_{g 1} v_{i t 2}+\psi_{g 1}+\overline{\mathbf{z}}_{i} \boldsymbol{\xi}_{g 1}+r_{i t 1} \geq 0\right]$,\\
where the coefficients have been scaled by a different variance and so we index the parameters by " $g$ " just to distinguish them from the previous scaled parameters. Now, $r_{i t 1}$ is independent of $\left(\mathbf{z}_{i}, y_{i t 2}, v_{i t 2}\right)$ with a standard normal distribution. After obtaining the $\hat{v}_{i t 2}$ as the residuals from pooled OLS estimation of the reduced form, we use pooled probit of $y_{i t 1}$ on $\mathbf{z}_{i t 1}, y_{i t 2}, \hat{v}_{i t 2}, 1$, and $\overline{\mathbf{z}}_{i}$. A simple test of the null hypothesis that $y_{i t 2}$ is contemporaneously endogenous is obtained by a $t$ test of $\mathrm{H}_{0}: \theta_{g 1}=0$, which is directly available from the pooled probit provided we make the statistic robust to arbitrary serial dependence.

\subsection*{15.8.6 Semiparametric Approaches}
Under strict exogeneity of the explanatory variables, it is possible to consistently estimate $\boldsymbol{\beta}$ up to scale under very weak assumptions. Manski (1987) derives an objective function that identifies $\boldsymbol{\beta}$ up to scale in the $T=2$ case when $e_{i 1}$ and $e_{i 2}$ in the model (15.72) are identically distributed conditional on ( $\mathbf{x}_{i 1}, \mathbf{x}_{i 2}, c_{i}$ ) and $\mathbf{x}_{i t}$ is strictly exogenous. The estimator is the maximum score estimator applied to the differences $\Delta y_{i}$ and $\Delta \mathbf{x}_{i}$. As in the cross-sectional case, it is not known how to estimate the average response probabilities.

Honoré and Kyriazidou (2000a) show how to estimate the parameters in the unobserved effects logit model with a lagged dependent variable and strictly exogenous explanatory variables without making distributional assumptions about the\\
unobserved effect. Unfortunately, the estimators, which are consistent and asymptotically normal, do not generally converge at the usual $\sqrt{N}$ rate. In addition, as with many semiparametric approaches, discrete explanatory variables such as time dummies are ruled out, and it is not possible to estimate the APEs. See also Arellano and Honoré (2001).

A middle ground between parametrically specifying $\mathrm{D}\left(c_{i} \mid \mathbf{x}_{i}\right)$ and allowing it to be completely unrestricted is to impose substantive assumptions on $\mathrm{D}\left(c_{i} \mid \mathbf{x}_{i}\right)$ but without making parametric assumptions. As a special case of Altonji and Matzkin's (2005) "exchangeability" assumption, we might impose the restriction


\begin{equation*}
\mathrm{D}\left(c_{i} \mid \mathbf{x}_{i}\right)=\mathrm{D}\left(c_{i} \mid \overline{\mathbf{x}}_{i}\right) \tag{15.91}
\end{equation*}


without specifying $\mathrm{D}\left(c_{i} \mid \overline{\mathbf{x}}_{i}\right)$. Why is this restriction useful? Consider a general specification where the only restriction is strict exogeneity conditional on the heterogeneity:\\
$\mathrm{P}\left(y_{i t}=1 \mid \mathbf{x}_{i}, \mathbf{c}_{i}\right)=\mathrm{P}\left(y_{i t}=1 \mid \mathbf{x}_{i t}, \mathbf{c}_{i}\right)=G_{t}\left(\mathbf{x}_{i t}, \mathbf{c}_{i}\right)$,\\
where $G_{t}$ is an unknown function taking values in $(0,1)$ and we let $\mathbf{c}_{i}$ be a vector of unobserved heterogeneity to emphasize the generality of the setup. We allow for a $t$ subscript on $G$ as a general way of allowing aggregate time effects (analogous to including different period dummies in a linear model, probit, or logit). The ASF at time $t$ can be written as\\
$\operatorname{ASF}_{t}\left(\mathbf{x}_{t}\right)=\mathrm{E}_{\mathbf{c}_{i}}\left[G_{t}\left(\mathbf{x}_{t}, \mathbf{c}_{i}\right)\right]=\mathrm{E}_{\overline{\mathbf{x}}_{i}}\left\{\mathrm{E}\left[G_{t}\left(\mathbf{x}_{t}, \mathbf{c}_{i}\right) \mid \overline{\mathbf{x}}_{i}\right]\right\} \equiv \mathrm{E}_{\overline{\mathbf{x}}_{i}}\left[R_{t}\left(\mathbf{x}_{t}, \overline{\mathbf{x}}_{i}\right)\right]$,\\
where $R_{t}\left(\mathbf{x}_{t}, \overline{\mathbf{x}}_{i}\right) \equiv \mathrm{E}\left[G_{t}\left(\mathbf{x}_{t}, \mathbf{c}_{i}\right) \mid \overline{\mathbf{x}}_{i}\right]$. It follows that, given an estimator $\hat{R}_{t}(\cdot, \cdot)$ of the function $R_{t}(\cdot, \cdot)$, the ASF can be estimated as $N^{-1} \sum_{i=1}^{N} \hat{R}_{t}\left(\mathbf{x}_{t}, \overline{\mathbf{x}}_{i}\right)$, and then we can take derivatives or changes with respect to the entries in $\mathbf{x}_{t}$.

How can we estimate $R_{t}(\cdot, \cdot)$ ? This is where assumption (15.91) comes into play. If we combine (15.91) and (15.92) we have

$$
\begin{aligned}
\mathrm{E}\left(y_{i t} \mid \mathbf{x}_{i}\right) & =\mathrm{E}\left[\mathrm{E}\left(y_{i t} \mid \mathbf{x}_{i}, \mathbf{c}_{i}\right) \mid \mathbf{x}_{i}\right]=\mathrm{E}\left[G_{t}\left(\mathbf{x}_{i t}, \mathbf{c}_{i}\right) \mid \mathbf{x}_{i}\right]=\int G_{t}\left(\mathbf{x}_{i t}, \mathbf{c}\right) \mathrm{d} F\left(\mathbf{c} \mid \mathbf{x}_{i}\right) \\
& =\int G_{t}\left(\mathbf{x}_{i t}, \mathbf{c}\right) \mathrm{d} F\left(\mathbf{c} \mid \overline{\mathbf{x}}_{i}\right)=R_{t}\left(\mathbf{x}_{i t}, \overline{\mathbf{x}}_{i}\right)
\end{aligned}
$$

where $F\left(\mathbf{c} \mid \mathbf{x}_{i}\right)$ denotes the cdf of $\mathrm{D}\left(\mathbf{c}_{i} \mid \mathbf{x}_{i}\right)$ (which can be a discrete, continuous, or mixed distribution), the second equality follows from (15.92), the fourth equality follows from assumption (15.91), and the last equality follows from the definition of $R_{t}(\cdot, \cdot)$. Of course, because $\mathrm{E}\left(y_{i t} \mid \mathbf{x}_{i}\right)$ depends only on $\left(\mathbf{x}_{i t}, \overline{\mathbf{x}}_{i}\right)$, we must have\\
$\mathrm{E}\left(y_{i t} \mid \mathbf{x}_{i t}, \overline{\mathbf{x}}_{i}\right)=R_{t}\left(\mathbf{x}_{i t}, \overline{\mathbf{x}}_{i}\right)$. Further, $\left\{\mathbf{x}_{i t}: t=1, \ldots, T\right\}$ is assumed to have time variation, and so $\mathbf{x}_{i t}$ and $\overline{\mathbf{x}}_{i}$ can be used as separate regressors even in a fully nonparametric setting.

We do not generally treat nonparametric methods in this text, but the preceding discussion suggests some simple yet flexible parametric approaches. The key is that, under specification (15.92) with assumption (15.91), we can specify flexible binary models for $\mathrm{P}\left(\mathrm{y}_{i t}=1 \mid \mathbf{x}_{i t}, \overline{\mathbf{x}}_{i}\right)$, estimate these models using MLE methods, and then average out $\overline{\mathbf{x}}_{i}$ to obtain estimated APEs. Without a very large $N$ or with many elements of $\mathbf{x}_{i t}$, we probably would economize by assuming that at least some parameters are constant across $t$. For example, a flexible probit model is\\
$\mathrm{P}\left(y_{i t}=1 \mid \mathbf{x}_{i t}, \overline{\mathbf{x}}_{i}\right)=\Phi\left[\theta_{t}+\mathbf{x}_{i t} \boldsymbol{\beta}+\overline{\mathbf{x}}_{i} \gamma+\left(\overline{\mathbf{x}}_{i} \otimes \overline{\mathbf{x}}_{i}\right) \boldsymbol{\delta}+\left(\mathbf{x}_{i t} \otimes \overline{\mathbf{x}}_{i}\right) \boldsymbol{\eta}\right], \quad t=1, \ldots, T$,\\
where the Kronecker products simply mean we include all squares and nonredundant interactions among $\overline{\mathbf{x}}_{i}$ and interactions among $\mathbf{x}_{i t}$ and $\overline{\mathbf{x}}_{i}$. Aggregate time effects are allowed through the $\theta_{t}$, and we can estimate this model by pooled probit, GEE, or minimum distance methods. Using a logit instead of a probit would likely change very little in terms of estimated partial effects. With large $N$, one might use a cdf that depends on extra parameters. The point here is that, under (15.91) and (15.92), the focus on APEs liberates one from having to specify specific functional forms in equation (15.92). Because we can only identify APEs anyway, rather than subjectspecific effects, we might as well start from models for $\mathrm{P}\left(y_{i t}=1 \mid \mathbf{x}_{i t}, \overline{\mathbf{x}}_{i}\right)$.

We can go further. For example, suppose that we think the heterogeneity $\mathbf{c}_{i}$ is correlated with features of the covariates other than just the time average. Altonji and Matzkin (2005) allow for $\overline{\mathbf{x}}_{i}$ in equation (15.91) to be replaced by other functions $\mathbf{w}_{i}$ of $\left\{\mathbf{x}_{i t}: t=1, \ldots, T\right\}$, such as sample variances and covariance. These are examples of "exchangeable" functions of $\left\{\mathbf{x}_{i t}: t=1, \ldots, T\right\}$-that is, statistics whose value is the same regardless of the ordering of the $\mathbf{x}_{i t}$. Nonexchangeable functions can be used, too. For example, we might think that $\mathbf{c}_{i}$ is correlated with individualspecific trends, and so we obtain $\mathbf{w}_{i}$ as the intercept and slope from the unit-specific regressions $\mathbf{x}_{i t}$ on $1, t, t=1, \ldots, T$ (for $T \geq 3$ ); we can also add the error variance from this individual specific regression if we have sufficient time periods.

Altonji and Matzkin (2005) focus on what they call the local average response (LAR) as opposed to the APE. In specification (15.92), the LAR at $\mathbf{x}_{t}$ for a continuous variable $x_{t j}$ is


\begin{equation*}
\int \frac{\partial G_{t}\left(\mathbf{x}_{t}, \mathbf{c}\right)}{\partial x_{t j}} \mathrm{~d} H_{t}\left(\mathbf{c} \mid \mathbf{x}_{t}\right) \tag{15.95}
\end{equation*}


where $H_{t}\left(\mathbf{c} \mid \mathbf{x}_{t}\right)$ denotes the cdf of $\mathrm{D}\left(\mathbf{c}_{i} \mid \mathbf{x}_{i t}=\mathbf{x}_{t}\right)$. This is a "local" partial effect because it averages out the heterogeneity for the slice of the population given by the vector $\mathbf{x}_{t}$. The APE, which by comparison could be called a "global average response," averages out over the entire distribution of $\mathbf{c}_{i}$.

When $\mathrm{D}\left(\mathbf{c}_{i} \mid \mathbf{x}_{i}\right)=\mathrm{D}\left(\mathbf{c}_{i} \mid \mathbf{w}_{i}\right)$ for $\mathbf{w}_{i}$ a function of $\mathbf{x}_{i}$, Altonji and Matzkin show that the LAR can be obtained as


\begin{equation*}
\int \frac{\partial R_{t}\left(\mathbf{x}_{t}, \mathbf{w}\right)}{\partial x_{t j}} \mathrm{~d} K_{t}\left(\mathbf{w} \mid \mathbf{x}_{t}\right), \tag{15.96}
\end{equation*}


where $R\left(\mathbf{x}_{t}, \mathbf{w}\right)=\mathrm{E}\left(y_{i t} \mid \mathbf{x}_{i t}=\mathbf{x}_{t}, \mathbf{w}_{i}=\mathbf{w}\right)$ and $K_{t}\left(\mathbf{w} \mid \mathbf{x}_{t}\right)$ is the cdf of $\mathrm{D}\left(\mathbf{w}_{i} \mid \mathbf{x}_{i t}=\mathbf{x}_{t}\right)$. Altonji and Matzkin demonstrate how to estimate the LAR based on nonparametric estimation of $\mathrm{E}\left(y_{i t} \mid \mathbf{x}_{i t}, \mathbf{w}_{i}\right)$ followed by "local" averaging, that is, averaging $\partial \hat{\mathrm{E}}\left(y_{i t} \mid \mathbf{x}_{i t}=\mathbf{x}_{t}, \mathbf{w}_{i}\right) / \partial x_{t j}$ over observations $i$ with $\mathbf{x}_{i t}$ "close" to $\mathbf{x}_{t}$. In the binary response context, the expected value is simply $\mathrm{P}\left(y_{i t}=1 \mid \mathbf{x}_{i t}, \mathbf{w}_{i}\right)$. Of course, the LAR can even be estimated using parametric models such as that in equation (15.94). For a continuous $x_{t j}$, we would simply average the derivative of $\Phi\left[\theta_{t}+\mathbf{x}_{t} \boldsymbol{\beta}+\overline{\mathbf{x}}_{i} \gamma+\right. \left.\left(\overline{\mathbf{x}}_{i} \otimes \overline{\mathbf{x}}_{i}\right) \boldsymbol{\delta}+\left(\mathbf{x}_{t} \otimes \overline{\mathbf{x}}_{i}\right) \boldsymbol{\eta}\right]$ with respect to $x_{t j}$ over $i$ with $\mathbf{x}_{i t}$ "close" to $\mathbf{x}_{t}$. Because defining the appropriate notion of "closeness" requires some care, the LAR is more difficult to estimate than the APE, but LAR is perhaps more relevant because it is the average response to an exogenous change in $x_{t j}$ for units already starting from $\mathbf{x}_{t}$. See Altonji and Matzkin (2005) for further discussion concerning identification and estimation of LARs.

An interesting possibility is to combine the Blundell and Powell (2004) approach to endogeneity with the Altonji and Matzkin (2005) approach to unobserved heterogeneity, resulting in semiparametric versions of the methods in Section 15.8.5.

\section*{Problems}
15.1. Suppose that $y$ is a binary outcome and $d 1, d 2, \ldots, d M$ are dummy variables for exhaustive and mutually exclusive categories; that is, each person in the population falls into one and only one category.\\
a. Show that the fitted values from the regression (without an intercept)\\
$y_{i}$ on $d 1_{i}, d 2_{i}, \ldots, d M_{i}, \quad i=1,2, \ldots, N$\\
are always in the unit interval. In particular, carefully describe the coefficient on each dummy variable and the fitted value for each $i$.\\
b. What happens if $y_{i}$ is regressed on $M$ linearly independent, linear combinations of $d 1_{i}, \ldots, d M_{i}$, for example, $1, d 2_{i}, d 3_{i}, \ldots, d M_{i}$ ?\\
15.2. Suppose that family $i$ chooses annual consumption $c_{i}$ (in dollars) and charitable contributions $q_{i}$ (in dollars) to solve the problem\\
$\max _{c, q} c+a_{i} \log (1+q)$\\
subject to $c+p_{i} q \leq m_{i}, \quad c, q \geq 0$\\
where $m_{i}$ is income of family $i, p_{i}$ is the price of one dollar of charitable contributions -where $p_{i}<1$ because of the tax deductability of charitable contributions, and this price differs across families because of different marginal tax rates and different state tax codes-and $a_{i} \geq 0$ determines the marginal utility of charitable contributions. Take $m_{i}$ and $p_{i}$ as exogenous to the family in this problem.\\
a. Show that the optimal solution is $q_{i}=0$ if $a_{i} \leq p_{i}$, and $q_{i}=a_{i} / p_{i}-1$ if $a_{i}>p_{i}$.\\
b. Define $y_{i}=1$ if $q_{i}>0$ and $y_{i}=0$ if $q_{i}=0$, and suppose that $a_{i}=\exp \left(\mathbf{z}_{i} \gamma+v_{i}\right)$, where $\mathbf{z}_{i}$ is a $J$-vector of observable family traits and $v_{i}$ is unobservable. Assume that $v_{i}$ is independent of ( $\mathbf{z}_{i}, m_{i}, p_{i}$ ) and $v_{i} / \sigma$ has symmetric distribution function $G(\cdot)$, where $\sigma^{2}=\operatorname{Var}\left(v_{i}\right)$. Show that\\
$\mathrm{P}\left(y_{i}=1 \mid \mathbf{z}_{i}, m_{i}, p_{i}\right)=G\left[\left(\mathbf{z}_{i} \gamma-\log p_{i}\right) / \sigma\right]$\\
so that $y_{i}$ follows an index model.\\
15.3. Let $\mathbf{z}_{1}$ be a vector of variables, let $z_{2}$ be a continuous variable, and let $d_{1}$ be a dummy variable.\\
a. In the model\\
$\mathbf{P}\left(y=1 \mid \mathbf{z}_{1}, z_{2}\right)=\Phi\left(\mathbf{z}_{1} \boldsymbol{\delta}_{1}+\gamma_{1} z_{2}+\gamma_{2} z_{2}^{2}\right)$,\\
find the partial effect of $z_{2}$ on the response probability. How would you estimate this partial effect?\\
b. In the model\\
$\mathbf{P}\left(y=1 \mid \mathbf{z}_{1}, z_{2}, d_{1}\right)=\Phi\left(\mathbf{z}_{1} \boldsymbol{\delta}_{1}+\gamma_{1} z_{2}+\gamma_{2} d_{1}+\gamma_{3} z_{2} d_{1}\right)$,\\
find the partial effect of $z_{2}$. How would you measure the effect of $d_{1}$ on the response probability? How would you estimate these effects?\\
c. Describe how you would obtain the standard errors of the estimated partial effects from parts a and b .\\
15.4. Evaluate the following statement: "Estimation of a linear probability model is more robust than probit or logit because the LPM does not assume homoskedasticity or a distributional assumption."\\
15.5. Consider the probit model\\
$\mathrm{P}(y=1 \mid \mathbf{z}, q)=\Phi\left(\mathbf{z}_{1} \boldsymbol{\delta}_{1}+\gamma_{1} z_{2} q\right)$,\\
where $q$ is independent of $\mathbf{z}$ and distributed as $\operatorname{Normal}(0,1)$; the vector $\mathbf{z}$ is observed but the scalar $q$ is not.\\
a. Find the partial effect of $z_{2}$ on the response probability, namely,\\
$\frac{\partial \mathbf{P}(y=1 \mid \mathbf{z}, q)}{\partial z_{2}}$.\\
b. Show that $\mathbf{P}(y=1 \mid \mathbf{z})=\Phi\left[\mathbf{z}_{1} \boldsymbol{\delta}_{1} /\left(1+\gamma_{1}^{2} z_{2}^{2}\right)^{1 / 2}\right]$.\\
c. Define $\rho_{1} \equiv \gamma_{1}^{2}$. How would you test $\mathrm{H}_{0}: \rho_{1}=0$ ?\\
d. If you have reason to believe $\rho_{1}>0$, how would you estimate $\boldsymbol{\delta}_{1}$ along with $\rho_{1}$ ?\\
15.6. Consider taking a large random sample of workers at a given point in time. Let $\operatorname{sick}_{i}=1$ if person $i$ called in sick during the last 90 days, and zero otherwise. Let $\mathbf{Z}_{i}$ be a vector of individual and employer characteristics. Let $c i g s_{i}$ be the number of cigarettes individual $i$ smokes per day (on average).\\
a. Explain the underlying experiment of interest when we want to examine the effects of cigarette smoking on workdays lost.\\
b. Why might $c i g s_{i}$ be correlated with unobservables affecting $\operatorname{sick}_{i}$ ?\\
c. One way to write the model of interest is\\
$\mathbf{P}\left(\right.$ sick $=1 \mid \mathbf{z}$, cigs,$\left.q_{1}\right)=\Phi\left(\mathbf{z}_{1} \boldsymbol{\delta}_{1}+\gamma_{1}\right.$ cigs $\left.+q_{1}\right)$,\\
where $\mathbf{z}_{1}$ is a subset of $\mathbf{z}$ and $q_{1}$ is an unobservable variable that is possibly correlated with cigs. What happens if $q_{1}$ is ignored and you estimate the probit of sick on $\mathbf{z}_{1}$, cigs?\\
d. Can cigs have a conditional normal distribution in the population? Explain.\\
e. Explain how to test whether cigs is exogenous. Does this test rely on cigs having a conditional normal distribution?\\
f. Suppose that some of the workers live in states that recently implemented nosmoking laws in the workplace. Does the presence of the new laws suggest a good IV candidate for cigs?\\
15.7. Use the data in GROGGER.RAW for this question.\\
a. Define a binary variable, say arr86, equal to unity if a man was arrested at least once during 1986, and zero otherwise. Estimate an LPM relating arr86 to pcnv, avgsen, tottime, ptime86, inc86, black, hispan, and born60. Report the usual and hetero-skedasticity-robust standard errors. What is the estimated effect on the probability of arrest if pcnv goes from .25 to .75 ?\\
b. Test the joint significance of avgsen and tottime, using a nonrobust and robust test.\\
c. Now estimate the model by probit. At the average values of avgsen, tottime, inc86, and ptime86 in the sample, and with black $=1$, hispan $=0$, and born60 $=1$, what is the estimated effect on the probability of arrest if pcnv goes from .25 to .75 ? Compare this result with the answer from part a.\\
d. For the probit model estimated in part c, obtain the percent correctly predicted. What is the percent correctly predicted when narr $86=0$ ? When narr $86=1$ ? What do you make of these findings?\\
e. In the probit model, add the terms $p c n v^{2}$, $p t i m e 86^{2}$, and $i n c 86^{2}$ to the model. Are these individually or jointly significant? Describe the estimated relationship between the probability of arrest and pcnv. In particular, at what point does the probability of conviction have a negative effect on probability of arrest?\\
15.8. Use the data set BWGHT.RAW for this problem.\\
a. Define a binary variable, smokes, if the woman smokes during pregnancy. Estimate a probit model relating smokes to motheduc, white, and $\log ($ faminc $)$. At white $=0$ and faminc evaluated at the average in the sample, what is the estimated difference in the probability of smoking for a woman with 16 years of education and one with 12 years of education?\\
b. Do you think faminc is exogenous in the smoking equation? What about motheduc?\\
c. Assume that motheduc and white are exogenous in the probit from part a. Also assume that fatheduc is exogenous to this equation. Estimate the reduced form for $\log ($ faminc $)$ to see if fatheduc is partially correlated with $\log ($ faminc $)$.\\
d. Test the null hypothesis that $\log ($ faminc $)$ is exogenous in the probit from part a.\\
15.9. Assume that the binary variable $y$ follows an LPM.\\
a. Write down the log-likelihood function for observation $i$.\\
b. Why might MLE of the LPM be difficult?\\
c. Assuming that you can estimate the LPM by MLE, explain why it is valid, as a model selection device, to compare the log likelihood from the LPM with that from logit or probit.\\
15.10. Suppose you wish to use goodness-of-fit measures to compare the LPM with a model such as logit or probit, after estimating the LPM by ordinary least squares. The usual $R$-squared from OLS estimation measures the proportion of the variance in $y$ that is explained by $\hat{\mathrm{P}}(y=1 \mid \mathbf{x})=\mathbf{x} \hat{\boldsymbol{\beta}}$.\\
a. Explain how to obtain a comparable $R$-squared measured for the general index $\operatorname{model} \mathrm{P}(y=1 \mid \mathbf{x})=G(\mathbf{x} \boldsymbol{\beta})$.\\
b. Compute the $R$-squared measures using the data in GROGGER.RAW, where the dependent variable is arr86 and the explanatory variables are $p c n v, p c n v^{2}, a v g s e n$, tottime, ptime86, ptime86 ${ }^{2}$, inc86, inc86 ${ }^{2}$, black, hispan, and born60. Are the $R$ squareds substantially different?\\
15.11. List assumptions under which the pooled probit estimator is a conditional MLE based on the distribution of $\mathbf{y}_{i}$ given $\mathbf{x}_{i}$, where $\mathbf{y}_{i}$ is the $T \times 1$ vector of binary outcomes and $\mathbf{x}_{i}$ is the vector of all explanatory variables across all $T$ time periods.\\
15.12. Find $\mathrm{P}\left(y_{i 1}=1, y_{i 2}=0, y_{i 3}=0 \mid \mathbf{x}_{i}, c_{i}, n_{i}=1\right)$ in the fixed effects logit model with $T=3$.\\
15.13. Suppose that you have a control group, $A$, and a treatment group, $B$, and two periods of data. Between the two years, a new policy is implemented that affects group B; see Section 6.5.\\
a. If your outcome variable is binary (for example, an employment indicator), and you have no covariates, how would you estimate the effect of the policy?\\
b. If you have covariates, write down a probit model that allows you to estimate the effect of the policy change. Explain in detail how you would estimate this effect.\\
c. How would you get an asymptotic 95 percent confidence interval for the estimate in part $b$ ?\\
15.14. Consider the binary response model with interactions between the (continuous) endogenous explanatory variable and the exogenous variables:\\
$y_{1}=1\left[\mathbf{z}_{1} \boldsymbol{\delta}_{1}+y_{2} \mathbf{z}_{1} \boldsymbol{\alpha}_{1}+u_{1}>0\right]$,\\
and assume that assumption (15.40) holds with $\left(u_{1}, v_{2}\right)$ jointly normal and independent of $\mathbf{z}$, with mean zero and $\operatorname{Var}\left(u_{1}\right)=1$.\\
a. Find $\mathrm{P}\left(y_{1}=1 \mid \mathbf{z}\right)$ and conclude that it has the form of a particular heteroskedastic probit model.\\
b. Consider the following two-step procedure. In the first step, regress $y_{i 2}$ on $\mathbf{z}_{i}$ and obtain the fitted values, $\hat{y}_{i 2}$. In the second step, estimate probit of $y_{i 1}$ on $\mathbf{z}_{i 1}, \hat{y}_{i 2} \mathbf{z}_{i 1}$. How come this does not produce consistent estimators of $\boldsymbol{\delta}_{1}$ and $\boldsymbol{\alpha}_{1}$ ?\\
c. How would you consistently estimate $\boldsymbol{\delta}_{1}$ and $\boldsymbol{\alpha}_{1}$ ?\\
15.15. Consider the problem of obtaining standard errors for the coefficient and APE estimators in two-step control function estimation of binary response models. For simplicity, let $y_{2}$ be a continuous, univariate endogenous variable following the reduced form $y_{2}=\mathbf{z} \boldsymbol{\delta}_{2}+v_{2}, y_{1}=1\left[\mathbf{x}_{1} \boldsymbol{\beta}_{1}+u_{1}>0\right], \mathbf{x}_{1}=\mathbf{g}_{1}\left(\mathbf{z}_{1}, y_{2}\right)$ is any function of the exogenous and endogenous variables. As usual, we assume that $\mathbf{z}=\left(\mathbf{z}_{1}, \mathbf{z}_{2}\right)$, where one element of $\mathbf{z}_{2}$ has nonzero element of $\boldsymbol{\delta}_{2}$. Assume that the standard RiversVuong assumptions hold, so that\\
$\mathrm{P}\left(y_{1}=1 \mid \mathbf{z}, y_{2}\right)=\Phi\left(\mathbf{x}_{1} \boldsymbol{\beta}_{\rho 1}+\theta_{\rho 1} v_{2}\right)$,\\
where $\boldsymbol{\beta}_{\rho 1}$ and $\theta_{\rho 1}$ are the scaled coefficients that appear in the APEs. For simplicity, let $\gamma_{1}=\left(\boldsymbol{\beta}_{\rho 1}^{\prime}, \theta_{\rho 1}\right)^{\prime}$, let $\hat{\gamma}_{1}$ be the second-step Rivers-Vuong estimator, and let $\hat{\boldsymbol{\delta}}_{2}$ be the first-stage OLS estimator.\\
a. Find $\mathrm{Avar} \sqrt{N}\left(\hat{\gamma}_{1}-\gamma_{1}\right)$ using the material in Section 12.4.2.\\
b. Show that the asymptotic variance that properly accounts for the first-step estimation of $\boldsymbol{\delta}_{2}$ is no smaller (in the matrix sense) than the asymptotic variance that ignores the estimation error in $\hat{\boldsymbol{\delta}}_{2}$.\\
c. Define the vector of estimated average partial effects (based on partial derivatives) as\\
$\hat{\boldsymbol{\eta}}_{1}=\left[N^{-1} \sum_{i=1}^{N} \Phi\left(\hat{\mathbf{w}}_{i 1} \hat{\gamma}_{1}\right)\right] \hat{\gamma}_{1}$,\\
where $\hat{\mathbf{w}}_{i 1}=\left(\mathbf{x}_{i 1}, \hat{v}_{i 2}\right)$. Find Avar $\sqrt{N}\left(\hat{\boldsymbol{\eta}}_{1}-\boldsymbol{\eta}_{1}\right)$. It may help to use Problem 12.17.\\
d. Show how to consistently estimate the asymptotic variance in part c.\\
15.16. Suppose that the binary variable $y$ follows the model\\
$\mathrm{P}(y=1 \mid \mathbf{x})=1-[1+\exp (\mathbf{x} \boldsymbol{\beta})]^{-\alpha}$,\\
where $\alpha>0$ is a parameter that can be estimated along with the vector $\boldsymbol{\beta}$. The $1 \times K$ vector $\mathbf{x}$ contains unity as its first element. This model is sometimes called the skewed logit model. Note that $\alpha=1$ corresponds to the usual logit model.\\
a. For a continuous variable $x_{K}$, find the partial effect of $x_{K}$ on $\mathrm{P}(y=1 \mid \mathbf{x})$.\\
b. Write down the log-likelihood function for a random draw $i$.\\
c. Use the data from MROZ.RAW to estimate the model, using the same explanatory variables in Table 15.1. Using a standard $t$ test, can you reject $\mathrm{H}_{0}: \log (\alpha)=0$ ?\\
d. Compute the likelihood ratio test of $\mathrm{H}_{0}: \alpha=1$. How does this compare with the test in part c?\\
e. Overall, would you say that the more complicated model is justified? What other statistics might you compute to support your answer?\\
15.17. In Example 15.4, add samesex as an explanatory variable to the worked equation in bivariate probit, and include samesex in the morekids probit, too.\\
a. What is the estimated APE with respect to morekids, and how does it compare with those in Table 15.2?\\
b. Is samesex significant in the worked equation?\\
15.18. Consider Chamberlain's random effects probit model under assumptions (15.65), (15.66), and (15.73), but replace assumption (15.67) with\\
$c_{i} \mid \mathbf{x}_{i} \sim \operatorname{Normal}\left[\psi+\overline{\mathbf{x}}_{i} \boldsymbol{\xi}, \sigma_{a}^{2} \exp \left(\overline{\mathbf{x}}_{i} \boldsymbol{\lambda}\right)\right]$,\\
so that $c_{i}$ given $\overline{\mathbf{x}}_{i}$ has exponential heteroskedasticity.\\
a. Find $\mathrm{P}\left(y_{i t}=1 \mid \mathbf{x}_{i}, a_{i}\right)$, where $a_{i}=c_{i}-\mathrm{E}\left(c_{i} \mid \mathbf{x}_{i}\right)$. Does this probability differ from the probability under assumption (15.73)? Explain.\\
b. Derive the log-likelihood function by first finding the density of $\left(y_{i 1}, \ldots, y_{i T}\right)$ given $\mathbf{x}_{i}$. Does it have similarities with the log-likelihood function under assumption (15.73)?\\
c. Assuming you have estimated $\boldsymbol{\beta}, \psi, \boldsymbol{\xi}, \sigma_{a}^{2}$, and $\lambda$ by CMLE, how would you estimate the average partial effects? \{Hint: First show that $\mathrm{E}\left[\Phi\left(\mathbf{x}^{\mathrm{o}} \boldsymbol{\beta}+\psi+\overline{\mathbf{x}}_{i} \boldsymbol{\xi}+a_{i}\right) \mid \mathbf{x}_{i}\right] =\Phi\left(\left\{\mathbf{x}^{\mathrm{o}} \boldsymbol{\beta}+\psi+\overline{\mathbf{x}}_{i} \xi\right\} /\left\{1+\sigma_{a}^{2} \exp \left(\overline{\mathbf{x}}_{i} \boldsymbol{\lambda}\right)\right\}^{1 / 2}\right)$, and then use the appropriate average across $i$.\\
15.19. Use the data in KEANE.RAW for this question, and restrict your attention to black men who are in the sample all 11 years.\\
a. Use pooled probit to estimate the model $\mathrm{P}\left(\right.$ employ $_{i t}=1 \mid$ employ $\left._{i, t-1}\right)= \Phi\left(\delta_{0}+\right.$ pemploy $\left._{i, t-1}\right)$. What assumption is needed to ensure that the usual standard errors and test statistics from pooled probit are asymptotically valid?\\
b. Estimate $\mathrm{P}\left(\right.$ employ $_{t}=1 \mid$ employ $\left._{t-1}=1\right)$ and $\mathrm{P}\left(\right.$ employ $_{t}=1 \mid$ employ $\left._{t-1}=0\right)$. Explain how you would obtain standard errors of these estimates.\\
c. Add a full set of year dummies to the analysis in part a, and estimate the probabilities in part $b$ for 1987. Are there important differences with the estimates in part $b$ ?\\
d. Now estimate a dynamic unobserved effects model using the method described in Section 15.8.4. In particular, add employ ${ }_{i, 81}$ as an additional explanatory variable, and use random effects probit software. Use a full set of year dummies.\\
e. Is there evidence of state dependence, conditional on $c_{i}$ ? Explain.\\
f. Average the estimated probabilities across employ ${ }_{i, 81}$ to get the average partial effect for 1987. Compare the estimates with the effects estimated in part c.

\section*{\begin{center}

\end{document}
